\documentclass[a4paper,11pt]{article}
        \usepackage[top=15mm,bottom=18mm,left=16mm,right=16mm]{geometry}
        \usepackage{fontspec}
        \usepackage{xeCJK}
        \setmainfont{Times New Roman}
        \setCJKmainfont{Yu Gothic}
        \setmonofont{Consolas}
        \setCJKmonofont{Yu Gothic}
        \usepackage{booktabs}
        \usepackage{longtable}
        \usepackage{tabularx}
        \usepackage{array}
        \usepackage{enumitem}
        \usepackage{hyperref}
        \usepackage{listings}
        \usepackage{xcolor}
        \usepackage{amsmath}
        \usepackage{amssymb}
        \hypersetup{
          colorlinks=true,
          linkcolor=blue,
          urlcolor=blue,
          pdftitle={フル MLE 同定本体},
          pdfauthor={Codex}
        }
        \lstset{
          basicstyle=\ttfamily\small,
          breaklines=true,
          columns=fullflexible,
          keepspaces=true,
          frame=single,
          framerule=0.2pt,
          rulecolor=\color{black},
          showstringspaces=false
        }
        \setlength{\parskip}{0.42em}
        \setlength{\parindent}{1em}
        \setlength{\emergencystretch}{3em}
        \renewcommand{\arraystretch}{1.15}
        \title{35. フル MLE 同定本体}
        \author{solar\_ws0129-main}
        \date{2026-07-29}
        \begin{document}
        \maketitle
        \tableofcontents
        \clearpage

        \section{対象}
        \begin{tabularx}{\linewidth}{>{\raggedright\arraybackslash}p{0.18\linewidth} X}
        \toprule
        項目 & 内容 \\
        \midrule
        ファイル & \path|scripts/run_vehicle_identification.py| \\
        ソースSHA-256 & \path|beea09d61241299016d808d93b3ce656b2bccebb2d44dade39cebff6710f289a| \\
        種別 & Python \\
        区分 & identification \\
        役割 & 実車ログ、weather、grounded maps、battery/PV/vehicle モデルを用いて MLE 同定を実行する大型スクリプト。 \\
        起動文脈 & fit action の本丸。 \\
        \bottomrule
        \end{tabularx}

        \section{このファイルを読む理由}
        実車ログ、weather、grounded maps、battery/PV/vehicle モデルを用いて MLE 同定を実行する大型スクリプト。

        \begin{itemize}[leftmargin=1.6em]
        \item 出力 tag を持つ immutable run を作る。
\item adopt までは canonical profile を上書きしない流れに対応する。
        \end{itemize}

        \section{呼び出し関係}
        \subsection{どこから呼ばれるか}
        \begin{itemize}[leftmargin=1.6em]
        \item \path|scripts/solar_control.sh|
        \end{itemize}

        \subsection{次に読むと理解が進むファイル}
        \begin{itemize}[leftmargin=1.6em]
        \item \path|scripts/generate_fit_fullsim_report.py|
\item \path|scripts/tune_upper_planner_weights.py|
        \end{itemize}

        \subsection{同一リポジトリ内の主要依存}
        \begin{itemize}[leftmargin=1.6em]
        \item 同一リポジトリ内の明示 import は少ない。
        \end{itemize}

        \section{ソース冒頭の把握}
        モジュール docstring または先頭意図の要約:

        モジュール docstring は明示されていない。

        \subsection{代表コード断片}
        \begin{lstlisting}[language=Python]
}


def neutralize_identification_scalars(model):
    """Fit each calibration factor once on top of the declared physical maps."""
    model.p.panel_gain = 1.0
    model.p.drive_eff_scale = 1.0
    model.p.regen_eff_scale = 1.0
    model.p.regen_utilization = 1.0
    model.p.rint_scale = 1.0
    model.p.r_polarization_ohm = 0.0
    model.aux_power_override_w = None
    return model


def resolve_relative(base_dir: Path, raw: str) -> Path:
    path = Path(str(raw or "").strip())
    if not path:
        return base_dir
    if path.is_absolute():
        return path
    return (base_dir / path).resolve()
\end{lstlisting}


        \section{主要構造の読み取り}
        主要関数は neutralize\_identification\_scalars, resolve\_relative, stage, load\_manifest, relpath\_from, tex\_path\_fragment, tex\_text\_fragment, load\_yaml\_if\_exists。 CLI 引数宣言は 9 件。

        \subsection{主要クラス・関数}
            \begin{longtable}{p{0.07\linewidth} p{0.16\linewidth} p{0.25\linewidth} p{0.40\linewidth}}
            \toprule
            行 & 種別 & 名前 & 補足 \\
            \midrule
            123 & function & \path|neutralize_identification_scalars| & args: model \\
135 & function & \path|resolve_relative| & args: base\_dir, raw \\
144 & function & \path|stage| & args: message \\
148 & function & \path|load_manifest| & args: package\_dir, manifest\_arg \\
168 & function & \path|relpath_from| & args: base\_dir, target \\
177 & function & \path|tex_path_fragment| & args: value \\
199 & function & \path|tex_text_fragment| & args: value \\
216 & function & \path|load_yaml_if_exists| & args: path \\
223 & function & \path|declared_control_stop_km| & args: profile\_cfg, profile\_path \\
245 & function & \path|_terminal_anchor_from_payload| & args: payload \\
292 & function & \path|_append_reason_column| & args: frame, mask, reason \\
300 & function & \path|resolve_manifest_context| & args: package\_dir, manifest \\
306 & function & \path|opt_path| & args: raw\_value \\
413 & function & \path|hampel_mask| & args: series \\
425 & function & \path|apply_sensor_quality_annotations| & args: work \\
520 & function & \path|polarization_current_trace| & args: replay \\
544 & function & \path|fit_battery_polarization| & args: replay \\
682 & function & \path|apply_battery_polarization| & args: replay, dynamic\_fit \\
706 & function & \path|resolve_fit_plan| & args: options \\
750 & function & \path|resolve_identification_output_layout| & args: package\_dir, profile\_cfg \\
777 & function & \path|identification_profile_output_path| & args: canonical\_profile, run\_output\_dir \\
789 & function & \path|load_ocv_df| & args: profile\_cfg, profile\_yaml \\
797 & function & \path|build_source_map_assets| & args: profile\_cfg, profile\_yaml \\
801 & function & \path|rel| & args: key, fallback \\
821 & function & \path|apply_actual_event_annotations| & args: work, payload \\
863 & function & \path|truncate_at_retire_event| & args: work, payload \\
892 & function & \path|normalize_generic_log| & args: log\_csv \\
994 & function & \path|build_terminal_anchor| & args: log\_df, ocv\_df, base\_model, anchor\_km \\
1029 & function & \path|terminal_metrics| & args: replay\_df, ocv\_df, base\_model, batt, anchor\_km \\
1078 & function & \path|replay_day_metrics| & args: replay\_df \\
1099 & function & \path|identification_selection_score| & args: metrics \\
1135 & function & \path|condition_vehicle_fit_on_measured_pv| & args: frame \\
1155 & function & \path|calibrate_solar_measurement_to_pack| & args: frame \\
1271 & function & \path|_shift_acceleration_within_segments| & args: frame, sample\_shift, acceleration \\
1283 & function & \path|_acceleration_trace_from_speed| & args: frame \\
1321 & function & \path|_bilinear_interp_array| & args: x\_grid, y\_grid, values, x, y \\
1341 & function & \path|_map_efficiency_array| & args: base\_model, maps, speed\_ms, torque\_nm \\
1364 & function & \path|_motion_predictions_for_acceleration| & args: frame, acceleration, base\_model, mot \\
1431 & function & \path|fit_acceleration_timestamp_alignment| & args: frame, base\_model, mot, options \\
1541 & function & \path|rmse| & args: mask \\
1631 & function & \path|_grade_from_smoothed_elevation| & args: distance\_km, elevation\_m, smoothing\_window\_km \\
1653 & function & \path|fit_grade_observation_alignment| & args: frame, route\_profile\_csv, output\_csv, base\_model, mot, options \\
1750 & function & \path|rmse| & args: predicted, mask \\
1886 & function & \path|pv_leave_one_day_out_validation| & args: frame, model \\
1950 & function & \path|rmse| & args: values \\
1964 & function & \path|add_end_to_end_metrics| & args: primary, end\_to\_end \\
1973 & function & \path|add_battery_conditional_metrics| & args: primary, conditional \\
1981 & function & \path|write_replay_csv| & args: frame, output\_path \\
1998 & function & \path|apply_fit_to_cfg| & args: cfg \\
2113 & function & \path|update_profile| & args: profile\_yaml, cfg, map\_assets, pv, batt, mot, observed\_log\_csv, battery\_dynamic\_fit, solar\_measurement\_calibration, route\_profile\_asset, package\_dir \\
2157 & function & \path|update_profile_artifact_references| & args: profile\_yaml, package\_dir \\
2180 & function & \path|write_terminal_consistency_from_anchor| & args: output\_path, terminal\_anchor \\
2318 & function & \path|sync_canonical_fullsim_profile| & args: package\_dir \\
2353 & function & \path|write_generic_summary| & args: package\_dir, manifest\_path, profile\_yaml, map\_assets, pv, batt, mot, metrics, terminal\_anchor, stage\_anchors, map\_shape\_fit, post\_refine, day\_metrics, battery\_dynamic\_fit, fit\_plan, manifest\_context, output\_dir, replay\_csv, battery\_conditioned\_replay\_csv, end\_to\_end\_replay\_csv \\
2422 & function & \path|write_generic_report| & args: package\_dir, profile\_yaml, manifest\_path, summary\_yaml, observed\_log\_csv, pv, batt, mot, metrics, post\_refine, map\_assets \\
3075 & function & \path|main| &  \\
            \bottomrule
            \end{longtable}

        \subsection{ファイルを上から読んだときの定義順}
            \begin{longtable}{p{0.07\linewidth} p{0.84\linewidth}}
            \toprule
            行 & 内容 \\
            \midrule
            20 & SCRIPT\_DIR に Path(\_\_file\_\_).resolve().parent の結果を代入する。 \\
21 & ROOT に SCRIPT\_DIR.parents[0] の結果を代入する。 \\
22 & (SCRIPT\_DIR, ROOT) を順に走査し、各要素を candidate に入れて処理する。 \\
67 & FIT\_QUALITY\_PRESETS に \{'quick': \{'battery\_restart\_count': 2, 'battery\_maxiter': 50, 'motion\_restart\_count': 3, 'motion\_maxiter': 70, 'joint\_restart\_count': 2, 'joint\_random\_start\_count': 8, 'joint\_local\_topk': 2, 'joint\_maxiter': 24, 'fit\_stride': 3, 'allow\_map\_shape\_fit': False, 'post\_refine\_enabled': False\}, 'standard': \{'battery\_restart\_count': 4, 'battery\_maxiter': 80, 'motion\_restart\_count': 5, 'motion\_maxiter': 100, 'joint\_restart\_count': 4, 'joint\_random\_start\_count': 6, 'joint\_local\_topk': 3, 'joint\_maxiter': 40, 'fit\_stride': 2, 'allow\_map\_shape\_fit': True, 'post\_refine\_enabled': False\}, 'full': \{'battery\_restart\_count': 6, 'battery\_maxiter': 140, 'motion\_restart\_count': 8, 'motion\_maxiter': 180, 'joint\_restart\_count': 6, 'joint\_random\_start\_count': 18, 'joint\_local\_topk': 6, 'joint\_maxiter': 80, 'fit\_stride': 2, 'allow\_map\_shape\_fit': True, 'post\_refine\_enabled': False\}, 'ultra': \{'battery\_restart\_count': 8, 'battery\_maxiter': 220, 'motion\_restart\_count': 10, 'motion\_maxiter': 260, 'joint\_restart\_count': 8, 'joint\_random\_start\_count': 28, 'joint\_local\_topk': 8, 'joint\_maxiter': 120, 'fit\_stride': 1, 'allow\_map\_shape\_fit': True, 'post\_refine\_enabled': True\}\} を代入する。 \\
123 & 関数 neutralize\_identification\_scalars を定義する。 \\
135 & 関数 resolve\_relative を定義する。 \\
144 & 関数 stage を定義する。 \\
148 & 関数 load\_manifest を定義する。 \\
168 & 関数 relpath\_from を定義する。 \\
177 & 関数 tex\_path\_fragment を定義する。 \\
199 & 関数 tex\_text\_fragment を定義する。 \\
216 & 関数 load\_yaml\_if\_exists を定義する。 \\
223 & 関数 declared\_control\_stop\_km を定義する。 \\
245 & 関数 \_terminal\_anchor\_from\_payload を定義する。 \\
292 & 関数 \_append\_reason\_column を定義する。 \\
300 & 関数 resolve\_manifest\_context を定義する。 \\
413 & 関数 hampel\_mask を定義する。 \\
425 & 関数 apply\_sensor\_quality\_annotations を定義する。 \\
520 & 関数 polarization\_current\_trace を定義する。 \\
544 & 関数 fit\_battery\_polarization を定義する。 \\
682 & 関数 apply\_battery\_polarization を定義する。 \\
706 & 関数 resolve\_fit\_plan を定義する。 \\
750 & 関数 resolve\_identification\_output\_layout を定義する。 \\
777 & 関数 identification\_profile\_output\_path を定義する。 \\
789 & 関数 load\_ocv\_df を定義する。 \\
797 & 関数 build\_source\_map\_assets を定義する。 \\
821 & 関数 apply\_actual\_event\_annotations を定義する。 \\
863 & 関数 truncate\_at\_retire\_event を定義する。 \\
892 & 関数 normalize\_generic\_log を定義する。 \\
994 & 関数 build\_terminal\_anchor を定義する。 \\
1029 & 関数 terminal\_metrics を定義する。 \\
1078 & 関数 replay\_day\_metrics を定義する。 \\
1099 & 関数 identification\_selection\_score を定義する。 \\
1135 & 関数 condition\_vehicle\_fit\_on\_measured\_pv を定義する。 \\
1155 & 関数 calibrate\_solar\_measurement\_to\_pack を定義する。 \\
1271 & 関数 \_shift\_acceleration\_within\_segments を定義する。 \\
1283 & 関数 \_acceleration\_trace\_from\_speed を定義する。 \\
1321 & 関数 \_bilinear\_interp\_array を定義する。 \\
1341 & 関数 \_map\_efficiency\_array を定義する。 \\
1364 & 関数 \_motion\_predictions\_for\_acceleration を定義する。 \\
            \bottomrule
            \end{longtable}

        \subsection{import 群の詳細}
            \begin{longtable}{p{0.07\linewidth} p{0.33\linewidth} p{0.50\linewidth}}
            \toprule
            行 & import 文 & このファイルで読む理由 \\
            \midrule
            2 & \path|from __future__ import annotations| & 型ヒントの遅延評価を有効にし、自己参照型や forward reference を安全に扱うため。 このファイル内での使用位置は少ないか、間接利用である。 \\
4 & \path|import argparse| & CLI 引数を宣言し、実行時パラメータを外から受け取るため。 このファイル内での主な使用位置は L3076。 \\
5 & \path|import json| & manifest、checkpoint、UDP payload をやり取りするため。 このファイル内での主な使用位置は L3181, L3841。 \\
6 & \path|import math| & Python標準のスカラー数値関数と定数を使うため。実際に参照する名前と行は使用位置欄で確認する。 このファイル内での主な使用位置は L539, L1595, L1596, L1834, L3619。 \\
7 & \path|import os| & パス、環境変数、プロセス外部状態を扱うため。 このファイル内での主な使用位置は L23, L172, L174, L2013, L2014, L2380, L2381, L2382, ...。 \\
8 & \path|import re| & re モジュールを利用するため。 このファイル内での主な使用位置は L764。 \\
9 & \path|import sys| & sys モジュールを利用するため。 このファイル内での主な使用位置は L24, L25。 \\
10 & \path|import textwrap| & textwrap モジュールを利用するため。 このファイル内での主な使用位置は L3070。 \\
11 & \path|from pathlib import Path| & ファイルやディレクトリを安全に扱うため。 このファイル内での主な使用位置は L20, L135, L136, L148, L151, L158, L168, L216, ...。 \\
12 & \path|from typing import Any, Dict, Tuple| & typing から Any, Dict, Tuple を読み込み、このファイルの処理を組み立てるため。 このファイル内での主な使用位置は L67, L148, L245, L251, L331, L366, L544, L600, ...。 \\
14 & \path|import numpy as np| & 数値配列、clip、補間、統計計算を行うため。 このファイル内での主な使用位置は L296, L422, L458, L466, L492, L502, L512, L525, ...。 \\
15 & \path|import pandas as pd| & CSV や時刻列の読込・整形・再サンプリングに使うため。 このファイル内での主な使用位置は L292, L413, L414, L426, L430, L464, L465, L466, ...。 \\
16 & \path|import yaml| & profile、stop、schedule、summary YAML を読み書きするため。 このファイル内での主な使用位置は L164, L220, L2154, L2164, L2174, L2311, L2333, L2349, ...。 \\
17 & \path|from scipy.optimize import least_squares| & scipy.optimize から least\_squares を読み込み、このファイルの処理を組み立てるため。 このファイル内での主な使用位置は L1206, L1213, L1232。 \\
18 & \path|from scipy.signal import savgol_filter| & scipy.signal から savgol\_filter を読み込み、このファイルの処理を組み立てるため。 このファイル内での主な使用位置は L1304, L1643。 \\
27 & \path|from build_bwsc2025_fitted_package import BATTERY_PACK_MAX_CHARGE_V, BATTERY_NOMINAL_VOLTAGE_V, INVALID_PACK_VOLTAGE_MIN_V, ROOT, TIMEZONE_LOCAL, BatteryFitResult, MotionFitResult, PostRefineResult, PvFitResult, attach_archive_pv_model, build_grounded_map_assets, build_stage_anchors, build_model_from_map_assets, build_model_from_profile_cfg, compile_tex, dcir_observations, ensure_dir, fit_battery_parameters, fit_map_shapes, fit_motion_parameters, fit_pv_parameters, fit_regen_utilization, fit_stop_tilt_fraction, infer_soc_from_loaded_state, joint_refine_parameters, joint_replay, load_profile_yaml, metrics_from_replay, motion_power_prediction, post_refine_replay_scalars, replay_segment_start_mask, resample_for_fit, soc_fit_upper_bound, write_current_maps_and_coefficients, write_scaled_maps| & build\_bwsc2025\_fitted\_package から BATTERY\_PACK\_MAX\_CHARGE\_V, BATTERY\_NOMINAL\_VOLTAGE\_V, INVALID\_PACK\_VOLTAGE\_MIN\_V, ROOT, TIMEZONE\_LOCAL, BatteryFitResult, MotionFitResult, PostRefineResult, PvFitResult, attach\_archive\_pv\_model, build\_grounded\_map\_assets, build\_stage\_anchors, build\_model\_from\_map\_assets, build\_model\_from\_profile\_cfg, compile\_tex, dcir\_observations, ensure\_dir, fit\_battery\_parameters, fit\_map\_shapes, fit\_motion\_parameters, fit\_pv\_parameters, fit\_regen\_utilization, fit\_stop\_tilt\_fraction, infer\_soc\_from\_loaded\_state, joint\_refine\_parameters, joint\_replay, load\_profile\_yaml, metrics\_from\_replay, motion\_power\_prediction, post\_refine\_replay\_scalars, replay\_segment\_start\_mask, resample\_for\_fit, soc\_fit\_upper\_bound, write\_current\_maps\_and\_coefficients, write\_scaled\_maps を読み込み、このファイルの処理を組み立てるため。 このファイル内での主な使用位置は L21, L22, L157, L381, L444, L531, L567, L842, ...。 \\
64 & \path|from audit_identification_residuals import run_audit as run_residual_audit| & audit\_identification\_residuals から run\_audit as run\_residual\_audit を読み込み、このファイルの処理を組み立てるため。 このファイル内での主な使用位置は L3671。 \\
            \bottomrule
            \end{longtable}

        \section{このファイルを読む前に必要な基礎知識}
        次の章は、構文やROS用語を既知と仮定しないための説明である。

        \subsection{プログラム、プロセス、メモリ、オブジェクトを区別する}
ソースファイルはディスク上の文字列であり、それ自体は走っていない。Python実行可能プログラムがソースを読み、OSがその実行を一つのプロセスとして管理する。プロセスには仮想メモリ、開いているファイル、スレッド、終了コードなどが対応する。


\begin{lstlisting}
ソースファイル -> Pythonインタプリタが読み込む -> OS上のプロセス -> Pythonオブジェクトをメモリに生成 -> 関数やcallbackを実行
\end{lstlisting}


「メモリ上に生成する」とは、実行中プロセスが使う記憶領域に、その値の型、属性、参照関係を表す実体を用意することである。変数は実体そのものというより、そのオブジェクトを指す名前として理解するとPythonの代入が読みやすい。


\begin{lstlisting}[language=python]
node = MPCNode()
alias = node
# nodeとaliasは同じオブジェクトを参照する。
\end{lstlisting}


一つのプロセスには複数スレッドを持たせられる。スレッドは同じプロセスのメモリを共有するため、受信callbackと最適化callbackが同じself.zを書き換える場合は実行順と排他を検討する必要がある。

\paragraph{根拠資料}
\begin{itemize}[leftmargin=1.6em]
\item \href{https://docs.python.org/3/tutorial/classes.html}{Python公式チュートリアル: Classes}
\end{itemize}

\subsection{名前、代入、型、参照}
Pythonの代入文は、右辺を先に評価し、その結果のオブジェクトへ左辺の名前を結び付ける。`x = f()`では、まずfを呼び、その戻り値が得られてからxが更新される。


`float(...)`、`int(...)`、`bool(...)`は型名を呼び出して値を変換する。外部CSV、YAML、ROS parameterから来た値は型が期待どおりとは限らないため、このリポジトリでは明示変換が多い。


`None`は値が無いことを示す単一のオブジェクト、`math.nan`は浮動小数点値ではあるが有効な数値ではないことを示す。両者は用途が異なるため、`is None`と`math.isfinite`を使い分ける。

\paragraph{根拠資料}
\begin{itemize}[leftmargin=1.6em]
\item \href{https://docs.python.org/3/tutorial/classes.html}{Python公式チュートリアル: Classes}
\end{itemize}

\subsection{関数、引数、戻り値、スコープ}
`def`は関数本体をその場で実行する文ではない。関数オブジェクトを作り、その名前を現在の名前空間へ登録する。`f`は関数そのもの、`f()`はその関数を呼んだ結果である。


仮引数は関数定義側の受け取り名、実引数は呼出側から渡す値である。位置引数、キーワード引数、既定値付き引数、`*args`、`**kwargs`は渡し方が異なる。


\begin{lstlisting}[language=python]
def clip_speed(value, minimum=0.0, maximum=110.0):
    return min(maximum, max(minimum, value))

v = clip_speed(120.0, maximum=90.0)
\end{lstlisting}


関数内で作った通常の名前はローカルスコープに属する。メソッドが`self.z`を更新すればオブジェクトの状態が残るが、単に`z`を更新しただけならその呼出しのローカル値である。

\paragraph{根拠資料}
\begin{itemize}[leftmargin=1.6em]
\item \href{https://docs.python.org/3/tutorial/classes.html}{Python公式チュートリアル: Classes}
\end{itemize}

\subsection{module、package、import、console\_scripts}
Pythonファイルはmoduleとして読み込める。複数moduleをディレクトリにまとめたものがpackageである。`from .model import SolarCarModel`の先頭の点は、現在と同じpackage内のmodel moduleを指す相対importである。


import時にはファイルのトップレベル文が上から一度実行される。`def`や`class`は関数・クラスを登録するが、その本体の通常処理は呼び出すまで走らない。


setuptoolsの`console\_scripts`は、端末で使う実行可能名と`package.module:function`を対応付けるインストール時メタデータである。生成される実行用ラッパーはmoduleをimportし、指定関数を呼び、戻り値を終了コードとして扱う小さな入口である。


\begin{lstlisting}
ROS launchのexecutable名 -> install済みconsole script -> mpc_solarcar.mpc_nodeをimport -> main()を呼ぶ
\end{lstlisting}

\paragraph{根拠資料}
\begin{itemize}[leftmargin=1.6em]
\item \href{https://setuptools.pypa.io/en/latest/userguide/entry_point.html}{setuptools公式: Entry Points / Console Scripts}
\item \href{https://docs.python.org/3/tutorial/classes.html}{Python公式チュートリアル: Classes}
\end{itemize}

\subsection{例外、try/finally、with、資源解放}
例外は通常の戻り値とは別経路で異常を呼出元へ伝える。`try/except`は想定した異常を処理し、`finally`は成功・失敗にかかわらず後始末を行う。


`with open(...) as f:`はcontext managerを使い、ブロックを出るとファイルを閉じる。CSVやログの破損を避けるため、開いた資源の所有者と閉じる場所を明確にする。


`except Exception: pass`はノードを止めない利点がある一方、入力異常を隠して原因追跡を難しくする。安全に関係する値では、少なくとも頻度制限付き警告、異常カウンタ、fallback状態のpublishを検討する。



\subsection{list、dict、deque、NumPy配列、pandas DataFrame}
listは順序付き可変列、tupleは順序付きで通常変更しない列、dictはキーから値を引く対応表、dequeは両端追加・削除が効率的なキューである。どの構造を選ぶかはアクセス方法と更新方法で決まる。


NumPy配列は同種数値を連続的に扱い、要素ごとの演算、clip、補間、線形代数を簡潔に書く。shapeは各次元の要素数であり、速度系列なら通常`(N,)`、候補集団なら`(population, N)`となる。


pandas DataFrameは列名を持つ表である。CSV読込後は列型、欠損、単位、timezone、並び順、重複を明示的に処理しなければ、数値計算が動いても意味が誤る。



\subsection{CLI、PowerShell、Bash、環境変数、終了コード}
CLIは端末からプログラム名と引数を渡す操作界面である。`argparse`は文字列として届く引数を名前、型、既定値、必須性に従って解析する。


PowerShellとBashは別のshellであり、変数記法、改行継続、引用、パス表記が異なる。このプロジェクトではWindows側のSolarSim.ps1がWSL側のsolar\_control.shへ処理を渡す。


環境変数は親プロセスから子プロセスへ受け渡される名前付き文字列である。ROS\_DOMAIN\_ID、RMW\_IMPLEMENTATION、Pythonの数値スレッド数などはコード外から動作を変えるため、実行記録へ残す必要がある。


終了コード0は一般に成功、0以外は失敗を示す。shellルータは子プロセスの終了コードを握り潰さず上位へ返すことで、自動運用が失敗を検知できる。

\paragraph{根拠資料}
\begin{itemize}[leftmargin=1.6em]
\item \href{https://setuptools.pypa.io/en/latest/userguide/entry_point.html}{setuptools公式: Entry Points / Console Scripts}
\end{itemize}

\subsection{目的関数、制約、L-BFGS-B、SHGO、有限grid証明}
数値最適化器は、利用者が与えた目的関数を複数の候補点で評価し、より小さい値を持つ候補を探す。solverが物理を理解するのではなく、物理と運用価値はcost関数へ書かれる。


L-BFGS-Bは変数ごとの上下限を扱える局所最適化法である。初期値の近くの谷へ収束し得るため、非凸問題では複数seedや大域探索と組み合わせる。successがFalseでも有限な候補が返る場合があるため、採用条件をコード側で決める。


SHGOは定めたsamplingと局所最適化を組み合わせる大域最適化法である。有限Cartesian gridの全列挙は、そのgrid上の最良を証明できるが、連続領域全体の最良を自動的に証明しない。資料ではこの証明範囲を区別する。

\paragraph{根拠資料}
\begin{itemize}[leftmargin=1.6em]
\item \href{https://docs.scipy.org/doc/scipy/reference/generated/scipy.optimize.minimize.html}{SciPy公式: scipy.optimize.minimize}
\item \href{https://docs.scipy.org/doc/scipy/reference/generated/scipy.optimize.shgo.html}{SciPy公式: scipy.optimize.shgo}
\end{itemize}

\subsection{天候、route、補間、時刻、単位}
予報は時刻だけの系列か、時刻とroute距離の2次元gridになり得る。現在時刻tと距離sがgrid点の間なら、周囲の値から補間して日射、温度、風を求める。


UTC、現地timezone、naive datetimeを混ぜると数時間のずれが発生する。入力でtimezoneを確定し、内部比較はUTC、表示は現地時刻という役割分担が安全である。


route距離のkm、速度のkm/hとm/s、時間の秒、energyのWhとJを跨ぐため、変換係数3.6、3600、1000の意味を各式で確認する。



\subsection{freshness、filter、guard、fallback、fail-safe}
分散システムでは最後に受け取った値が現在も有効とは限らない。受信時刻とtimeoutからfreshnessを判定し、stale値を計画状態へ無条件に同期しない。


filterはnoiseと一時的な飛び値を抑えるが、遅れを生む。slew limitは指令変化率を制限する。安全guardはsolverのcost罰則とは別に、現在出力へ強制制約を適用する最後の防波堤である。


fallbackは失敗時の代替動作を事前に決める設計である。前回計画保持、物理に基づく決定論的入力、停止、低速制限などから、故障modeごとに選ぶ。fallback発生はstatusとlogへ残し、正常解と区別する。



\subsection{CSV/YAMLのdata contractと検証}
data contractは列名、型、単位、timezone、欠損可否、並び順、重複、許容範囲、先頭行、encodingを事前に決めた仕様である。単にCSVとして読めることは、モデル入力として正しいことを意味しない。


同定用実測、map、route、forecast、stop、scheduleはそれぞれgrainが異なる。生成時にschema validation、物理範囲、時間単調性、route範囲、coverageを検査し、検査結果をartifactとして残す。


学習用データと独立検証データを分離し、RMSEだけでなくbias、時系列残差、energy積算誤差、終端SoC、温度・電圧制約、外挿領域を評価する。



\subsection{単体試験、SILS、replay、観測可能性}
単体試験は関数やモデル項の局所契約、SILSは複数Nodeと時間進行を含む閉ループ、historical replayは実測入力への再現性、本番preflightは実機通信と停止動作を確認する。目的が異なるため一つで代用しない。


再現可能な診断には、入力、出力、内部状態、solver成否、候補数、計算時間、fallback、source revision、profile hashを同じrun IDで記録する。


非同期不具合は最終値だけでは追えない。topic周期、message age、callback所要時間、queue遅延、Node生存、publisher数を同時に観測する。

\paragraph{根拠資料}
\begin{itemize}[leftmargin=1.6em]
\item \href{https://docs.ros.org/en/ros2_packages/humble/api/rqt_graph/index.html}{ROS 2 Humble公式: rqt\_graph}
\item \href{https://docs.ros.org/en/humble/Tutorials/Beginner-CLI-Tools/Recording-And-Playing-Back-Data/Recording-And-Playing-Back-Data.html}{ROS 2 Humble公式: Recording and playing back data}
\item \href{https://docs.ros.org/en/rolling/Concepts/Intermediate/About-Executors.html}{ROS 2公式: Executors}
\end{itemize}



        \section{関数・クラスを上から順に解説}
        \subsection{L123 関数 neutralize\_identification\_scalars}
            \begin{itemize}[leftmargin=1.6em]
              \item 行範囲: L123--L132
              \item 所属: module直下
              \item 定義: \path|neutralize_identification_scalars(model)|
              \item docstring: Fit each calibration factor once on top of the declared physical maps.
              \item このブロックが直接呼ぶ主な関数・メソッド: 特に明示されていない。
              \item 戻り値の要点: model
              \item 主なローカル代入名: 特になし。
              \item 読み取る主なインスタンス属性: 特になし。
              \item 更新する主なインスタンス属性: 特になし。
              \item 明示的に送出する例外: 明示的raiseはない。
              \item 制御構造の規模: 条件分岐 0、ループ 0、try 0
            \end{itemize}
            \paragraph{この定義を読むためのPython構文}
            \begin{itemize}[leftmargin=1.6em]
            \item 特別な構文上の補足は少ない。
            \end{itemize}
            \paragraph{上から順の処理}
            \begin{enumerate}[leftmargin=1.8em]
            \item model.p.panel\_gain に 1.0 の結果を代入する。
\item model.p.drive\_eff\_scale に 1.0 の結果を代入する。
\item model.p.regen\_eff\_scale に 1.0 の結果を代入する。
\item model.p.regen\_utilization に 1.0 の結果を代入する。
\item model.p.rint\_scale に 1.0 の結果を代入する。
\item model.p.r\_polarization\_ohm に 0.0 の結果を代入する。
\item model.aux\_power\_override\_w に None の結果を代入する。
\item model を返す。
            \end{enumerate}
            \paragraph{代表コード断片}
            \begin{lstlisting}[language=Python]
def neutralize_identification_scalars(model):
    """Fit each calibration factor once on top of the declared physical maps."""
    model.p.panel_gain = 1.0
    model.p.drive_eff_scale = 1.0
    model.p.regen_eff_scale = 1.0
    model.p.regen_utilization = 1.0
    model.p.rint_scale = 1.0
    model.p.r_polarization_ohm = 0.0
    model.aux_power_override_w = None
    return model
\end{lstlisting}

\subsection{L135 関数 resolve\_relative}
            \begin{itemize}[leftmargin=1.6em]
              \item 行範囲: L135--L141
              \item 所属: module直下
              \item 定義: \path|resolve_relative(base_dir: Path, raw: str) -> Path|
              \item docstring: 明示されていない。
              \item このブロックが直接呼ぶ主な関数・メソッド: Path, is\_absolute, resolve, str, strip
              \item 戻り値の要点: (base\_dir / path).resolve() / base\_dir / path
              \item 主なローカル代入名: path
              \item 読み取る主なインスタンス属性: 特になし。
              \item 更新する主なインスタンス属性: 特になし。
              \item 明示的に送出する例外: 明示的raiseはない。
              \item 制御構造の規模: 条件分岐 2、ループ 0、try 0
            \end{itemize}
            \paragraph{この定義を読むためのPython構文}
            \begin{itemize}[leftmargin=1.6em]
            \item `->`以後は戻り値の型注釈であり、実行時の値を自動変換する命令ではない。
            \end{itemize}
            \paragraph{上から順の処理}
            \begin{enumerate}[leftmargin=1.8em]
            \item path に Path(str(raw or '').strip()) の結果を代入する。
\item 条件 not path を判定し、真なら内部処理を行う。
\item   base\_dir を返す。
\item 条件 path.is\_absolute() を判定し、真なら内部処理を行う。
\item   path を返す。
\item (base\_dir / path).resolve() を返す。
            \end{enumerate}
            \paragraph{代表コード断片}
            \begin{lstlisting}[language=Python]
def resolve_relative(base_dir: Path, raw: str) -> Path:
    path = Path(str(raw or "").strip())
    if not path:
        return base_dir
    if path.is_absolute():
        return path
    return (base_dir / path).resolve()
\end{lstlisting}

\subsection{L144 関数 stage}
            \begin{itemize}[leftmargin=1.6em]
              \item 行範囲: L144--L145
              \item 所属: module直下
              \item 定義: \path|stage(message: str) -> None|
              \item docstring: 明示されていない。
              \item このブロックが直接呼ぶ主な関数・メソッド: print
              \item 戻り値の要点: 戻り値が無いか、return 文が明示されていない。
              \item 主なローカル代入名: 特になし。
              \item 読み取る主なインスタンス属性: 特になし。
              \item 更新する主なインスタンス属性: 特になし。
              \item 明示的に送出する例外: 明示的raiseはない。
              \item 制御構造の規模: 条件分岐 0、ループ 0、try 0
            \end{itemize}
            \paragraph{この定義を読むためのPython構文}
            \begin{itemize}[leftmargin=1.6em]
            \item `->`以後は戻り値の型注釈であり、実行時の値を自動変換する命令ではない。
            \end{itemize}
            \paragraph{上から順の処理}
            \begin{enumerate}[leftmargin=1.8em]
            \item print(...) を実行する。
            \end{enumerate}
            \paragraph{代表コード断片}
            \begin{lstlisting}[language=Python]
def stage(message: str) -> None:
    print(f"[generic-id] {message}", flush=True)
\end{lstlisting}

\subsection{L148 関数 load\_manifest}
            \begin{itemize}[leftmargin=1.6em]
              \item 行範囲: L148--L165
              \item 所属: module直下
              \item 定義: \path!load_manifest(package_dir: Path, manifest_arg: str | None) -> Tuple[Path, dict]!
              \item docstring: 明示されていない。
              \item このブロックが直接呼ぶ主な関数・メソッド: Path, cwd, is\_absolute, is\_file, next, open, resolve, safe\_load, str, strip
              \item 戻り値の要点: (manifest\_path, payload)
              \item 主なローカル代入名: candidates, default\_path, f, manifest\_path, path, payload, raw\_path
              \item 読み取る主なインスタンス属性: 特になし。
              \item 更新する主なインスタンス属性: 特になし。
              \item 明示的に送出する例外: 明示的raiseはない。
              \item 制御構造の規模: 条件分岐 2、ループ 0、try 0
            \end{itemize}
            \paragraph{この定義を読むためのPython構文}
            \begin{itemize}[leftmargin=1.6em]
            \item `->`以後は戻り値の型注釈であり、実行時の値を自動変換する命令ではない。
\item 内包表記は、反復・条件・値生成を一つの式で記述してcollectionを作る。
\item with文はcontext managerに開始・終了処理を任せ、ファイルなどの資源を確実に解放する。
            \end{itemize}
            \paragraph{上から順の処理}
            \begin{enumerate}[leftmargin=1.8em]
            \item default\_path に package\_dir / 'data' / 'identification' / 'identification\_manifest.yaml' の結果を代入する。
\item 条件 manifest\_arg を判定し、真なら内部処理を行う。
\item   raw\_path に Path(str(manifest\_arg).strip()) の結果を代入する。
\item   条件 raw\_path.is\_absolute() を判定し、真なら内部処理を行う。
\item     manifest\_path に raw\_path の結果を代入する。
\item     上の条件が偽の場合:
\item     candidates に [(package\_dir / raw\_path).resolve(), (ROOT / raw\_path).resolve(), (Path.cwd() / raw\_path).resolve()] の結果を代入する。
\item     manifest\_path に next((path for path in candidates if path.is\_file()), candidates[0]) の結果を代入する。
\item   上の条件が偽の場合:
\item   manifest\_path に default\_path の結果を代入する。
\item with 文で manifest\_path.open('r', encoding='utf-8') を管理しながら処理する。
\item   payload に yaml.safe\_load(f) or \{\} の結果を代入する。
\item (manifest\_path, payload) を返す。
            \end{enumerate}
            \paragraph{代表コード断片}
            \begin{lstlisting}[language=Python]
def load_manifest(package_dir: Path, manifest_arg: str | None) -> Tuple[Path, dict]:
    default_path = package_dir / "data" / "identification" / "identification_manifest.yaml"
    if manifest_arg:
        raw_path = Path(str(manifest_arg).strip())
        if raw_path.is_absolute():
            manifest_path = raw_path
        else:
            candidates = [
                (package_dir / raw_path).resolve(),
                (ROOT / raw_path).resolve(),
                (Path.cwd() / raw_path).resolve(),
            ]
            manifest_path = next((path for path in candidates if path.is_file()), candidates[0])
    else:
        manifest_path = default_path
    with manifest_path.open("r", encoding="utf-8") as f:
        payload = yaml.safe_load(f) or {}
    return manifest_path, payload
\end{lstlisting}

\subsection{L168 関数 relpath\_from}
            \begin{itemize}[leftmargin=1.6em]
              \item 行範囲: L168--L174
              \item 所属: module直下
              \item 定義: \path!relpath_from(base_dir: Path, target: Path | None) -> str!
              \item docstring: 明示されていない。
              \item このブロックが直接呼ぶ主な関数・メソッド: fspath, relpath, replace
              \item 戻り値の要点: '' / os.path.relpath(target, base\_dir).replace('\textbackslash{}\textbackslash{}', '/') / os.fspath(target)
              \item 主なローカル代入名: 特になし。
              \item 読み取る主なインスタンス属性: 特になし。
              \item 更新する主なインスタンス属性: 特になし。
              \item 明示的に送出する例外: 明示的raiseはない。
              \item 制御構造の規模: 条件分岐 1、ループ 0、try 1
            \end{itemize}
            \paragraph{この定義を読むためのPython構文}
            \begin{itemize}[leftmargin=1.6em]
            \item `->`以後は戻り値の型注釈であり、実行時の値を自動変換する命令ではない。
\item try文は例外が起き得る処理と、異常時または終了時の経路を分ける。
            \end{itemize}
            \paragraph{上から順の処理}
            \begin{enumerate}[leftmargin=1.8em]
            \item 条件 target is None を判定し、真なら内部処理を行う。
\item   '' を返す。
\item 例外処理を伴う try ブロックを実行する。
\item   os.path.relpath(target, base\_dir).replace('\textbackslash{}\textbackslash{}', '/') を返す。
\item   Exceptionを捕捉した場合:
\item   os.fspath(target) を返す。
            \end{enumerate}
            \paragraph{代表コード断片}
            \begin{lstlisting}[language=Python]
def relpath_from(base_dir: Path, target: Path | None) -> str:
    if target is None:
        return ""
    try:
        return os.path.relpath(target, base_dir).replace("\\", "/")
    except Exception:
        return os.fspath(target)
\end{lstlisting}

\subsection{L177 関数 tex\_path\_fragment}
            \begin{itemize}[leftmargin=1.6em]
              \item 行範囲: L177--L196
              \item 所属: module直下
              \item 定義: \path|tex_path_fragment(value: object) -> str|
              \item docstring: Render ASCII and Unicode paths without losing CJK glyphs or TeX syntax.
              \item このブロックが直接呼ぶ主な関数・メソッド: get, isascii, join, replace, str
              \item 戻り値の要点: f'\textbackslash{}\textbackslash{}texttt\{\{\{escaped\}\}\}' / f'\textbackslash{}\textbackslash{}path\{\{\{text\}\}\}'
              \item 主なローカル代入名: char, escaped, replacements, text
              \item 読み取る主なインスタンス属性: 特になし。
              \item 更新する主なインスタンス属性: 特になし。
              \item 明示的に送出する例外: 明示的raiseはない。
              \item 制御構造の規模: 条件分岐 1、ループ 0、try 0
            \end{itemize}
            \paragraph{この定義を読むためのPython構文}
            \begin{itemize}[leftmargin=1.6em]
            \item `->`以後は戻り値の型注釈であり、実行時の値を自動変換する命令ではない。
\item 内包表記は、反復・条件・値生成を一つの式で記述してcollectionを作る。
            \end{itemize}
            \paragraph{上から順の処理}
            \begin{enumerate}[leftmargin=1.8em]
            \item text に str(value) の結果を代入する。
\item 条件 text.isascii() を判定し、真なら内部処理を行う。
\item   f'\textbackslash{}\textbackslash{}path\{\{\{text\}\}\}' を返す。
\item replacements に \{'\textbackslash{}\textbackslash{}': '\textbackslash{}\textbackslash{}textbackslash\{\}', '\{': '\textbackslash{}\textbackslash{}\{', '\}': '\textbackslash{}\textbackslash{}\}', '\$': '\textbackslash{}\textbackslash{}\$', '\&': '\textbackslash{}\textbackslash{}\&', '\#': '\textbackslash{}\textbackslash{}\#', '\%': '\textbackslash{}\textbackslash{}\%', '\_': '\textbackslash{}\textbackslash{}\_', '\textasciicircum{}': '\textbackslash{}\textbackslash{}textasciicircum\{\}', '\textasciitilde{}': '\textbackslash{}\textbackslash{}textasciitilde\{\}'\} の結果を代入する。
\item escaped に ''.join((replacements.get(char, char) for char in text)) の結果を代入する。
\item escaped に escaped.replace('/', '/\textbackslash{}\textbackslash{}allowbreak\{\}') の結果を代入する。
\item f'\textbackslash{}\textbackslash{}texttt\{\{\{escaped\}\}\}' を返す。
            \end{enumerate}
            \paragraph{代表コード断片}
            \begin{lstlisting}[language=Python]
def tex_path_fragment(value: object) -> str:
    """Render ASCII and Unicode paths without losing CJK glyphs or TeX syntax."""
    text = str(value)
    if text.isascii():
        return rf"\path{{{text}}}"
    replacements = {
        "\\": r"\textbackslash{}",
        "{": r"\{",
        "}": r"\}",
        "$": r"\$",
        "&": r"\&",
        "#": r"\#",
        "%": r"\%",
        "_": r"\_",
        "^": r"\textasciicircum{}",
        "~": r"\textasciitilde{}",
    }
    escaped = "".join(replacements.get(char, char) for char in text)
    escaped = escaped.replace("/", r"/\allowbreak{}")
    return rf"\texttt{{{escaped}}}"
\end{lstlisting}

\subsection{L199 関数 tex\_text\_fragment}
            \begin{itemize}[leftmargin=1.6em]
              \item 行範囲: L199--L213
              \item 所属: module直下
              \item 定義: \path|tex_text_fragment(value: object) -> str|
              \item docstring: Escape arbitrary report prose without path-style whitespace handling.
              \item このブロックが直接呼ぶ主な関数・メソッド: get, join, str
              \item 戻り値の要点: ''.join((replacements.get(char, char) for char in str(value)))
              \item 主なローカル代入名: char, replacements
              \item 読み取る主なインスタンス属性: 特になし。
              \item 更新する主なインスタンス属性: 特になし。
              \item 明示的に送出する例外: 明示的raiseはない。
              \item 制御構造の規模: 条件分岐 0、ループ 0、try 0
            \end{itemize}
            \paragraph{この定義を読むためのPython構文}
            \begin{itemize}[leftmargin=1.6em]
            \item `->`以後は戻り値の型注釈であり、実行時の値を自動変換する命令ではない。
\item 内包表記は、反復・条件・値生成を一つの式で記述してcollectionを作る。
            \end{itemize}
            \paragraph{上から順の処理}
            \begin{enumerate}[leftmargin=1.8em]
            \item replacements に \{'\textbackslash{}\textbackslash{}': '\textbackslash{}\textbackslash{}textbackslash\{\}', '\{': '\textbackslash{}\textbackslash{}\{', '\}': '\textbackslash{}\textbackslash{}\}', '\$': '\textbackslash{}\textbackslash{}\$', '\&': '\textbackslash{}\textbackslash{}\&', '\#': '\textbackslash{}\textbackslash{}\#', '\%': '\textbackslash{}\textbackslash{}\%', '\_': '\textbackslash{}\textbackslash{}\_', '\textasciicircum{}': '\textbackslash{}\textbackslash{}textasciicircum\{\}', '\textasciitilde{}': '\textbackslash{}\textbackslash{}textasciitilde\{\}'\} の結果を代入する。
\item ''.join((replacements.get(char, char) for char in str(value))) を返す。
            \end{enumerate}
            \paragraph{代表コード断片}
            \begin{lstlisting}[language=Python]
def tex_text_fragment(value: object) -> str:
    """Escape arbitrary report prose without path-style whitespace handling."""
    replacements = {
        "\\": r"\textbackslash{}",
        "{": r"\{",
        "}": r"\}",
        "$": r"\$",
        "&": r"\&",
        "#": r"\#",
        "%": r"\%",
        "_": r"\_",
        "^": r"\textasciicircum{}",
        "~": r"\textasciitilde{}",
    }
    return "".join(replacements.get(char, char) for char in str(value))
\end{lstlisting}

\subsection{L216 関数 load\_yaml\_if\_exists}
            \begin{itemize}[leftmargin=1.6em]
              \item 行範囲: L216--L220
              \item 所属: module直下
              \item 定義: \path!load_yaml_if_exists(path: Path | None) -> dict!
              \item docstring: 明示されていない。
              \item このブロックが直接呼ぶ主な関数・メソッド: exists, open, safe\_load
              \item 戻り値の要点: \{\} / yaml.safe\_load(f) or \{\}
              \item 主なローカル代入名: f
              \item 読み取る主なインスタンス属性: 特になし。
              \item 更新する主なインスタンス属性: 特になし。
              \item 明示的に送出する例外: 明示的raiseはない。
              \item 制御構造の規模: 条件分岐 1、ループ 0、try 0
            \end{itemize}
            \paragraph{この定義を読むためのPython構文}
            \begin{itemize}[leftmargin=1.6em]
            \item `->`以後は戻り値の型注釈であり、実行時の値を自動変換する命令ではない。
\item with文はcontext managerに開始・終了処理を任せ、ファイルなどの資源を確実に解放する。
            \end{itemize}
            \paragraph{上から順の処理}
            \begin{enumerate}[leftmargin=1.8em]
            \item 条件 path is None or not path.exists() を判定し、真なら内部処理を行う。
\item   \{\} を返す。
\item with 文で path.open('r', encoding='utf-8') を管理しながら処理する。
\item   yaml.safe\_load(f) or \{\} を返す。
            \end{enumerate}
            \paragraph{代表コード断片}
            \begin{lstlisting}[language=Python]
def load_yaml_if_exists(path: Path | None) -> dict:
    if path is None or not path.exists():
        return {}
    with path.open("r", encoding="utf-8") as f:
        return yaml.safe_load(f) or {}
\end{lstlisting}

\subsection{L223 関数 declared\_control\_stop\_km}
            \begin{itemize}[leftmargin=1.6em]
              \item 行範囲: L223--L242
              \item 所属: module直下
              \item 定義: \path|declared_control_stop_km(profile_cfg: dict, profile_path: Path) -> list[float]|
              \item docstring: Load control-stop distances declared by the vehicle package.
              \item このブロックが直接呼ぶ主な関数・メソッド: append, bool, float, get, isinstance, load\_yaml\_if\_exists, resolve\_relative, set, sorted, str, strip
              \item 戻り値の要点: [] / sorted(set(distances))
              \item 主なローカル代入名: distances, item, key, paths, payload, raw, stops
              \item 読み取る主なインスタンス属性: 特になし。
              \item 更新する主なインスタンス属性: 特になし。
              \item 明示的に送出する例外: 明示的raiseはない。
              \item 制御構造の規模: 条件分岐 3、ループ 2、try 1
            \end{itemize}
            \paragraph{この定義を読むためのPython構文}
            \begin{itemize}[leftmargin=1.6em]
            \item `->`以後は戻り値の型注釈であり、実行時の値を自動変換する命令ではない。
\item try文は例外が起き得る処理と、異常時または終了時の経路を分ける。
            \end{itemize}
            \paragraph{上から順の処理}
            \begin{enumerate}[leftmargin=1.8em]
            \item paths に profile\_cfg.get('paths', \{\}) if isinstance(profile\_cfg, dict) else \{\} の結果を代入する。
\item ('actual\_stop\_yaml', 'stop\_yaml') を順に走査し、各要素を key に入れて処理する。
\item   raw に str(paths.get(key, '') or '').strip() の結果を代入する。
\item   条件 not raw を判定し、真なら内部処理を行う。
\item     Continue 文を実行する。
\item   payload に load\_yaml\_if\_exists(resolve\_relative(profile\_path.parent, raw)) の結果を代入する。
\item   stops に payload.get('stops', []) if isinstance(payload, dict) else [] の結果を代入する。
\item   distances に [] の結果を代入する。
\item   stops を順に走査し、各要素を item に入れて処理する。
\item     条件 not isinstance(item, dict) or not bool(item.get('is\_control\_stop', False)) を判定し、真なら内部処理を行う。
\item       Continue 文を実行する。
\item     例外処理を伴う try ブロックを実行する。
\item       distances.append(...) を実行する。
\item       (KeyError, TypeError, ValueError)を捕捉した場合:
\item       Continue 文を実行する。
\item   条件 distances を判定し、真なら内部処理を行う。
\item     sorted(set(distances)) を返す。
\item [] を返す。
            \end{enumerate}
            \paragraph{代表コード断片}
            \begin{lstlisting}[language=Python]
def declared_control_stop_km(profile_cfg: dict, profile_path: Path) -> list[float]:
    """Load control-stop distances declared by the vehicle package."""
    paths = profile_cfg.get("paths", {}) if isinstance(profile_cfg, dict) else {}
    for key in ("actual_stop_yaml", "stop_yaml"):
        raw = str(paths.get(key, "") or "").strip()
        if not raw:
            continue
        payload = load_yaml_if_exists(resolve_relative(profile_path.parent, raw))
        stops = payload.get("stops", []) if isinstance(payload, dict) else []
        distances = []
        for item in stops:
            if not isinstance(item, dict) or not bool(item.get("is_control_stop", False)):
                continue
            try:
                distances.append(float(item["s_km"]))
            except (KeyError, TypeError, ValueError):
                continue
        if distances:
            return sorted(set(distances))
    return []
\end{lstlisting}

\subsection{L245 関数 \_terminal\_anchor\_from\_payload}
            \begin{itemize}[leftmargin=1.6em]
              \item 行範囲: L245--L289
              \item 所属: module直下
              \item 定義: \path|_terminal_anchor_from_payload(payload: dict) -> Dict[str, Any]|
              \item docstring: 明示されていない。
              \item このブロックが直接呼ぶ主な関数・メソッド: bool, float, get, isinstance, str, strip
              \item 戻り値の要点: out / \{\} / \{\}
              \item 主なローカル代入名: anchor, key, out, raw, raw\_time
              \item 読み取る主なインスタンス属性: 特になし。
              \item 更新する主なインスタンス属性: 特になし。
              \item 明示的に送出する例外: 明示的raiseはない。
              \item 制御構造の規模: 条件分岐 6、ループ 3、try 1
            \end{itemize}
            \paragraph{この定義を読むためのPython構文}
            \begin{itemize}[leftmargin=1.6em]
            \item `->`以後は戻り値の型注釈であり、実行時の値を自動変換する命令ではない。
\item try文は例外が起き得る処理と、異常時または終了時の経路を分ける。
            \end{itemize}
            \paragraph{上から順の処理}
            \begin{enumerate}[leftmargin=1.8em]
            \item 条件 not isinstance(payload, dict) を判定し、真なら内部処理を行う。
\item   \{\} を返す。
\item anchor に payload.get('terminal\_anchor', payload) の結果を代入する。
\item 条件 not isinstance(anchor, dict) を判定し、真なら内部処理を行う。
\item   \{\} を返す。
\item out に \{\} を代入する。
\item ('count', 's\_km', 'voltage\_v', 'current\_a', 'temp\_c', 'soc\_target', 'soc\_sigma', 'soc\_evidence\_min', 'soc\_evidence\_max', 'voltage\_sigma\_v', 'ocv\_terminal\_v', 'series\_resistance\_ohm', 'ocv\_statistical\_sigma\_v', 'ocv\_systematic\_sigma\_v', 'ocv\_total\_sigma\_v') を順に走査し、各要素を key に入れて処理する。
\item   raw に anchor.get(key, None) の結果を代入する。
\item   条件 raw in (None, '') を判定し、真なら内部処理を行う。
\item     Continue 文を実行する。
\item   例外処理を伴う try ブロックを実行する。
\item     out[key] に float(raw) の結果を代入する。
\item     Exceptionを捕捉した場合:
\item     Pass 文を実行する。
\item raw\_time に str(anchor.get('time\_utc', '') or '').strip() の結果を代入する。
\item 条件 raw\_time を判定し、真なら内部処理を行う。
\item   out['time\_utc'] に raw\_time の結果を代入する。
\item ('notes', 'source\_documents', 'method', 'soc\_target\_basis') を順に走査し、各要素を key に入れて処理する。
\item   条件 key in anchor を判定し、真なら内部処理を行う。
\item     out[key] に anchor[key] の結果を代入する。
\item ('quality\_gate\_pass', 'conditional\_on\_grounded\_ocv\_map', 'weak\_channel\_cross\_consistency\_gate\_pass') を順に走査し、各要素を key に入れて処理する。
\item   条件 key in anchor を判定し、真なら内部処理を行う。
\item     out[key] に bool(anchor[key]) の結果を代入する。
\item out を返す。
            \end{enumerate}
            \paragraph{代表コード断片}
            \begin{lstlisting}[language=Python]
def _terminal_anchor_from_payload(payload: dict) -> Dict[str, Any]:
    if not isinstance(payload, dict):
        return {}
    anchor = payload.get("terminal_anchor", payload)
    if not isinstance(anchor, dict):
        return {}
    out: Dict[str, Any] = {}
    for key in (
        "count",
        "s_km",
        "voltage_v",
        "current_a",
        "temp_c",
        "soc_target",
        "soc_sigma",
        "soc_evidence_min",
        "soc_evidence_max",
        "voltage_sigma_v",
        "ocv_terminal_v",
        "series_resistance_ohm",
        "ocv_statistical_sigma_v",
        "ocv_systematic_sigma_v",
        "ocv_total_sigma_v",
    ):
        raw = anchor.get(key, None)
        if raw in (None, ""):
            continue
        try:
            out[key] = float(raw)
        except Exception:
            pass
    raw_time = str(anchor.get("time_utc", "") or "").strip()
    if raw_time:
        out["time_utc"] = raw_time
    for key in ("notes", "source_documents", "method", "soc_target_basis"):
...
\end{lstlisting}

\subsection{L292 関数 \_append\_reason\_column}
            \begin{itemize}[leftmargin=1.6em]
              \item 行範囲: L292--L297
              \item 所属: module直下
              \item 定義: \path|_append_reason_column(frame: pd.DataFrame, mask: pd.Series, reason: str) -> None|
              \item docstring: 明示されていない。
              \item このブロックが直接呼ぶ主な関数・メソッド: any, astype, fillna, len, where
              \item 戻り値の要点: 戻り値が無いか、return 文が明示されていない。
              \item 主なローカル代入名: current, merged
              \item 読み取る主なインスタンス属性: 特になし。
              \item 更新する主なインスタンス属性: 特になし。
              \item 明示的に送出する例外: 明示的raiseはない。
              \item 制御構造の規模: 条件分岐 1、ループ 0、try 0
            \end{itemize}
            \paragraph{この定義を読むためのPython構文}
            \begin{itemize}[leftmargin=1.6em]
            \item `->`以後は戻り値の型注釈であり、実行時の値を自動変換する命令ではない。
            \end{itemize}
            \paragraph{上から順の処理}
            \begin{enumerate}[leftmargin=1.8em]
            \item 条件 not mask.any() を判定し、真なら内部処理を行う。
\item    を返す。
\item current に frame.loc[mask, 'exclude\_reason'].fillna('').astype(str) の結果を代入する。
\item merged に np.where(current.str.len() > 0, current + ';' + reason, reason) の結果を代入する。
\item frame.loc[mask, 'exclude\_reason'] に merged の結果を代入する。
            \end{enumerate}
            \paragraph{代表コード断片}
            \begin{lstlisting}[language=Python]
def _append_reason_column(frame: pd.DataFrame, mask: pd.Series, reason: str) -> None:
    if not mask.any():
        return
    current = frame.loc[mask, "exclude_reason"].fillna("").astype(str)
    merged = np.where(current.str.len() > 0, current + ";" + reason, reason)
    frame.loc[mask, "exclude_reason"] = merged
\end{lstlisting}

\subsection{L300 関数 resolve\_manifest\_context}
            \begin{itemize}[leftmargin=1.6em]
              \item 行範囲: L300--L410
              \item 所属: module直下
              \item 定義: \path|resolve_manifest_context(package_dir: Path, manifest: dict) -> dict|
              \item docstring: 明示されていない。
              \item このブロックが直接呼ぶ主な関数・メソッド: FileNotFoundError, Path, ValueError, \_terminal\_anchor\_from\_payload, append, difference, exists, get, is\_absolute, isinstance, items, join
              \item 戻り値の要点: \{'inputs': inputs, 'options': options, 'grounded\_sources': grounded, 'evidence': evidence, 'actual\_event\_path': actual\_event\_path, 'counterfactual\_event\_path': counterfactual\_event\_path, 'terminal\_anchor\_path': terminal\_anchor\_path, 'grounded\_summary\_path': grounded\_summary\_path, 'source\_inventory\_path': source\_inventory\_path, 'notes\_markdown\_path': notes\_markdown\_path, 'explicit\_grounded\_assets': explicit\_grounded\_assets, 'grounded\_summary\_payload': grounded\_summary\_payload, 'actual\_event\_payload': actual\_event\_payload, 'terminal\_anchor\_override': terminal\_anchor\_override, 'external\_documents': external\_documents, 'declared\_evidence': declared\_evidence\} / resolve\_relative(package\_dir, raw) / None
              \item 主なローカル代入名: actual\_event\_path, actual\_event\_payload, candidates, counterfactual\_event\_path, declared\_evidence, evidence, explicit\_grounded\_assets, external\_documents, grounded, grounded\_summary\_path, grounded\_summary\_payload, group\_key, inputs, manifest\_key, missing, missing\_evidence, model\_key, notes\_markdown\_path, options, path
              \item 読み取る主なインスタンス属性: 特になし。
              \item 更新する主なインスタンス属性: 特になし。
              \item 明示的に送出する例外: FileNotFoundError('identification evidence manifest contains missing artifacts:\textbackslash{}n- ' + '\textbackslash{}n- '.join(missing\_evidence)), ValueError('grounded\_sources declares explicit maps but is missing required entries: ' + ', '.join(missing))
              \item 制御構造の規模: 条件分岐 15、ループ 3、try 0
            \end{itemize}
            \paragraph{この定義を読むためのPython構文}
            \begin{itemize}[leftmargin=1.6em]
            \item `->`以後は戻り値の型注釈であり、実行時の値を自動変換する命令ではない。
\item 内包表記は、反復・条件・値生成を一つの式で記述してcollectionを作る。
            \end{itemize}
            \paragraph{上から順の処理}
            \begin{enumerate}[leftmargin=1.8em]
            \item inputs に manifest.get('inputs', \{\}) if isinstance(manifest, dict) else \{\} の結果を代入する。
\item options に manifest.get('options', \{\}) if isinstance(manifest, dict) else \{\} の結果を代入する。
\item grounded に manifest.get('grounded\_sources', \{\}) if isinstance(manifest, dict) else \{\} の結果を代入する。
\item evidence に manifest.get('evidence', \{\}) if isinstance(manifest, dict) else \{\} の結果を代入する。
\item 関数 opt\_path を定義する。
\item actual\_event\_path に opt\_path(inputs.get('actual\_event\_yaml')) の結果を代入する。
\item counterfactual\_event\_path に opt\_path(inputs.get('counterfactual\_event\_yaml')) の結果を代入する。
\item terminal\_anchor\_path に opt\_path(inputs.get('terminal\_anchor\_yaml')) の結果を代入する。
\item grounded\_summary\_path に opt\_path(grounded.get('grounded\_map\_summary\_yaml', '') or options.get('grounded\_map\_summary\_yaml', '')) の結果を代入する。
\item source\_inventory\_path に opt\_path(evidence.get('source\_inventory\_json')) の結果を代入する。
\item notes\_markdown\_path に opt\_path(evidence.get('notes\_markdown')) の結果を代入する。
\item path\_key\_map に \{'drive\_eff\_map': 'drive\_eff\_map\_csv', 'drive\_map\_eco': 'drive\_map\_eco\_csv', 'drive\_map\_power': 'drive\_map\_power\_csv', 'regen\_eff\_map': 'regen\_eff\_map\_csv', 'regen\_map\_eco': 'regen\_map\_eco\_csv', 'regen\_map\_power': 'regen\_map\_power\_csv', 'rint\_map': 'rint\_map\_csv', 'panel\_eff\_map': 'panel\_eff\_map\_csv', 'mppt\_eff\_map': 'mppt\_eff\_map\_csv', 'ocv\_soc\_map': 'ocv\_soc\_map\_csv'\} の結果を代入する。
\item explicit\_grounded\_assets に \{\} を代入する。
\item path\_key\_map.items() を順に走査し、各要素を (model\_key, manifest\_key) に入れて処理する。
\item   raw に str(grounded.get(manifest\_key, '') or '').strip() の結果を代入する。
\item   条件 raw を判定し、真なら内部処理を行う。
\item     explicit\_grounded\_assets[model\_key] に resolve\_relative(package\_dir, raw) の結果を代入する。
\item 条件 explicit\_grounded\_assets を判定し、真なら内部処理を行う。
\item   条件 'drive\_map\_eco' not in explicit\_grounded\_assets and 'drive\_eff\_map' in explicit\_grounded\_assets を判定し、真なら内部処理を行う。
\item     explicit\_grounded\_assets['drive\_map\_eco'] に explicit\_grounded\_assets['drive\_eff\_map'] の結果を代入する。
\item   条件 'drive\_map\_power' not in explicit\_grounded\_assets and 'drive\_eff\_map' in explicit\_grounded\_assets を判定し、真なら内部処理を行う。
\item     explicit\_grounded\_assets['drive\_map\_power'] に explicit\_grounded\_assets['drive\_eff\_map'] の結果を代入する。
\item   条件 'regen\_map\_eco' not in explicit\_grounded\_assets and 'regen\_eff\_map' in explicit\_grounded\_assets を判定し、真なら内部処理を行う。
\item     explicit\_grounded\_assets['regen\_map\_eco'] に explicit\_grounded\_assets['regen\_eff\_map'] の結果を代入する。
\item   条件 'regen\_map\_power' not in explicit\_grounded\_assets and 'regen\_eff\_map' in explicit\_grounded\_assets を判定し、真なら内部処理を行う。
\item     explicit\_grounded\_assets['regen\_map\_power'] に explicit\_grounded\_assets['regen\_eff\_map'] の結果を代入する。
\item   required に \{'drive\_eff\_map', 'regen\_eff\_map', 'rint\_map', 'panel\_eff\_map', 'mppt\_eff\_map', 'ocv\_soc\_map'\} の結果を代入する。
\item   missing に sorted(required.difference(explicit\_grounded\_assets)) の結果を代入する。
\item   条件 missing を判定し、真なら内部処理を行う。
\item     ValueError('grounded\_sources declares explicit maps but is missing required entries: ' + ', '.join(missing)) を送出する。
\item grounded\_summary\_payload に load\_yaml\_if\_exists(grounded\_summary\_path) の結果を代入する。
\item 条件 grounded\_summary\_path is not None を判定し、真なら内部処理を行う。
\item   grounded\_summary\_payload['summary\_yaml'] に relpath\_from(package\_dir, grounded\_summary\_path) の結果を代入する。
\item actual\_event\_payload に load\_yaml\_if\_exists(actual\_event\_path) の結果を代入する。
\item terminal\_anchor\_payload に load\_yaml\_if\_exists(terminal\_anchor\_path) の結果を代入する。
\item terminal\_anchor\_override に \_terminal\_anchor\_from\_payload(terminal\_anchor\_payload) の結果を代入する。
\item external\_documents に evidence.get('external\_documents', []) の結果を代入する。
\item 条件 not isinstance(external\_documents, list) を判定し、真なら内部処理を行う。
\item   external\_documents に [external\_documents] の結果を代入する。
\item declared\_evidence に \{\} を代入する。
\item missing\_evidence に [] を代入する。
\item ('field\_tests', 'normalized\_field\_outputs', 'external\_documents') を順に走査し、各要素を group\_key に入れて処理する。
\item   raw\_values に evidence.get(group\_key, []) の結果を代入する。
\item   条件 not isinstance(raw\_values, list) を判定し、真なら内部処理を行う。
\item     raw\_values に [raw\_values] の結果を代入する。
\item   resolved\_group に [] を代入する。
\item   raw\_values を順に走査し、各要素を raw\_value に入れて処理する。
\item     raw に str(raw\_value or '').strip() の結果を代入する。
\item     条件 not raw を判定し、真なら内部処理を行う。
\item       Continue 文を実行する。
\item     raw\_path に Path(raw) の結果を代入する。
\item     条件 raw\_path.is\_absolute() を判定し、真なら内部処理を行う。
\item       resolved に raw\_path の結果を代入する。
\item       上の条件が偽の場合:
\item       candidates に [(package\_dir / raw\_path).resolve(), (ROOT / raw\_path).resolve()] の結果を代入する。
\item       resolved に next((path for path in candidates if path.exists()), candidates[0]) の結果を代入する。
\item     resolved\_group.append(...) を実行する。
\item     条件 not resolved.exists() を判定し、真なら内部処理を行う。
\item       missing\_evidence.append(...) を実行する。
\item   declared\_evidence[group\_key] に resolved\_group の結果を代入する。
\item 条件 missing\_evidence を判定し、真なら内部処理を行う。
\item   FileNotFoundError('identification evidence manifest contains missing artifacts:\textbackslash{}n- ' + '\textbackslash{}n- '.join(missing\_evidence)) を送出する。
\item \{'inputs': inputs, 'options': options, 'grounded\_sources': grounded, 'evidence': evidence, 'actual\_event\_path': actual\_event\_path, 'counterfactual\_event\_path': counterfactual\_event\_path, 'terminal\_anchor\_path': terminal\_anchor\_path, 'grounded\_summary\_path': grounded\_summary\_path, 'source\_inventory\_path': source\_inventory\_path, 'notes\_markdown\_path': notes\_markdown\_path, 'explicit\_grounded\_assets': explicit\_grounded\_assets, 'grounded\_summary\_payload': grounded\_summary\_payload, 'actual\_event\_payload': actual\_event\_payload, 'terminal\_anchor\_override': terminal\_anchor\_override, 'external\_documents': external\_documents, 'declared\_evidence': declared\_evidence\} を返す。
            \end{enumerate}
            \paragraph{代表コード断片}
            \begin{lstlisting}[language=Python]
def resolve_manifest_context(package_dir: Path, manifest: dict) -> dict:
    inputs = manifest.get("inputs", {}) if isinstance(manifest, dict) else {}
    options = manifest.get("options", {}) if isinstance(manifest, dict) else {}
    grounded = manifest.get("grounded_sources", {}) if isinstance(manifest, dict) else {}
    evidence = manifest.get("evidence", {}) if isinstance(manifest, dict) else {}

    def opt_path(raw_value) -> Path | None:
        raw = str(raw_value or "").strip()
        if not raw:
            return None
        return resolve_relative(package_dir, raw)

    actual_event_path = opt_path(inputs.get("actual_event_yaml"))
    counterfactual_event_path = opt_path(inputs.get("counterfactual_event_yaml"))
    terminal_anchor_path = opt_path(inputs.get("terminal_anchor_yaml"))
    grounded_summary_path = opt_path(grounded.get("grounded_map_summary_yaml", "") or options.get("grounded_map_summary_yaml", ""))
    source_inventory_path = opt_path(evidence.get("source_inventory_json"))
    notes_markdown_path = opt_path(evidence.get("notes_markdown"))

    path_key_map = {
        "drive_eff_map": "drive_eff_map_csv",
        "drive_map_eco": "drive_map_eco_csv",
        "drive_map_power": "drive_map_power_csv",
        "regen_eff_map": "regen_eff_map_csv",
        "regen_map_eco": "regen_map_eco_csv",
        "regen_map_power": "regen_map_power_csv",
        "rint_map": "rint_map_csv",
        "panel_eff_map": "panel_eff_map_csv",
        "mppt_eff_map": "mppt_eff_map_csv",
        "ocv_soc_map": "ocv_soc_map_csv",
    }
    explicit_grounded_assets: Dict[str, Path] = {}
    for model_key, manifest_key in path_key_map.items():
        raw = str(grounded.get(manifest_key, "") or "").strip()
        if raw:
...
\end{lstlisting}

\subsection{L306 関数 resolve\_manifest\_context.opt\_path}
            \begin{itemize}[leftmargin=1.6em]
              \item 行範囲: L306--L310
              \item 所属: \path|resolve_manifest_context|
              \item 定義: \path!opt_path(raw_value) -> Path | None!
              \item docstring: 明示されていない。
              \item このブロックが直接呼ぶ主な関数・メソッド: resolve\_relative, str, strip
              \item 戻り値の要点: resolve\_relative(package\_dir, raw) / None
              \item 主なローカル代入名: raw
              \item 読み取る主なインスタンス属性: 特になし。
              \item 更新する主なインスタンス属性: 特になし。
              \item 明示的に送出する例外: 明示的raiseはない。
              \item 制御構造の規模: 条件分岐 1、ループ 0、try 0
            \end{itemize}
            \paragraph{この定義を読むためのPython構文}
            \begin{itemize}[leftmargin=1.6em]
            \item `->`以後は戻り値の型注釈であり、実行時の値を自動変換する命令ではない。
            \end{itemize}
            \paragraph{上から順の処理}
            \begin{enumerate}[leftmargin=1.8em]
            \item raw に str(raw\_value or '').strip() の結果を代入する。
\item 条件 not raw を判定し、真なら内部処理を行う。
\item   None を返す。
\item resolve\_relative(package\_dir, raw) を返す。
            \end{enumerate}
            \paragraph{代表コード断片}
            \begin{lstlisting}[language=Python]
    def opt_path(raw_value) -> Path | None:
        raw = str(raw_value or "").strip()
        if not raw:
            return None
        return resolve_relative(package_dir, raw)
\end{lstlisting}

\subsection{L413 関数 hampel\_mask}
            \begin{itemize}[leftmargin=1.6em]
              \item 行範囲: L413--L422
              \item 所属: module直下
              \item 定義: \path|hampel_mask(series: pd.Series, *, window: int, n_sigma: float, min_abs: float) -> pd.Series|
              \item docstring: 明示されていない。
              \item このブロックが直接呼ぶ主な関数・メソッド: abs, astype, clip, float, int, max, maximum, median, rolling, to\_numeric
              \item 戻り値の要点: abs\_dev > np.maximum(float(min\_abs), float(n\_sigma) * scale)
              \item 主なローカル代入名: abs\_dev, mad, med, scale, window, x
              \item 読み取る主なインスタンス属性: 特になし。
              \item 更新する主なインスタンス属性: 特になし。
              \item 明示的に送出する例外: 明示的raiseはない。
              \item 制御構造の規模: 条件分岐 1、ループ 0、try 0
            \end{itemize}
            \paragraph{この定義を読むためのPython構文}
            \begin{itemize}[leftmargin=1.6em]
            \item `->`以後は戻り値の型注釈であり、実行時の値を自動変換する命令ではない。
            \end{itemize}
            \paragraph{上から順の処理}
            \begin{enumerate}[leftmargin=1.8em]
            \item x に pd.to\_numeric(series, errors='coerce').astype(float) の結果を代入する。
\item window に max(3, int(window)) の結果を代入する。
\item 条件 window \% 2 == 0 を判定し、真なら内部処理を行う。
\item   window を Add で更新する。
\item med に x.rolling(window, center=True, min\_periods=1).median() の結果を代入する。
\item abs\_dev に (x - med).abs() の結果を代入する。
\item mad に abs\_dev.rolling(window, center=True, min\_periods=1).median() の結果を代入する。
\item scale に (1.4826 * mad).clip(lower=max(1e-06, float(min\_abs) * 0.25)) の結果を代入する。
\item abs\_dev > np.maximum(float(min\_abs), float(n\_sigma) * scale) を返す。
            \end{enumerate}
            \paragraph{代表コード断片}
            \begin{lstlisting}[language=Python]
def hampel_mask(series: pd.Series, *, window: int, n_sigma: float, min_abs: float) -> pd.Series:
    x = pd.to_numeric(series, errors="coerce").astype(float)
    window = max(3, int(window))
    if window % 2 == 0:
        window += 1
    med = x.rolling(window, center=True, min_periods=1).median()
    abs_dev = (x - med).abs()
    mad = abs_dev.rolling(window, center=True, min_periods=1).median()
    scale = (1.4826 * mad).clip(lower=max(1.0e-6, float(min_abs) * 0.25))
    return abs_dev > np.maximum(float(min_abs), float(n_sigma) * scale)
\end{lstlisting}

\subsection{L425 関数 apply\_sensor\_quality\_annotations}
            \begin{itemize}[leftmargin=1.6em]
              \item 行範囲: L425--L517
              \item 所属: module直下
              \item 定義: \path!apply_sensor_quality_annotations(work: pd.DataFrame, *, base_model, options: dict | None = None) -> pd.DataFrame!
              \item docstring: 明示されていない。
              \item このブロックが直接呼ぶ主な関数・メソッド: \_append\_reason\_column, abs, any, astype, copy, diff, fillna, float, get, getattr, hampel\_mask, int
              \item 戻り値の要点: out / work
              \item 主なローカル代入名: charge\_floor\_a, current, current\_limit\_margin\_a, current\_spike, current\_spike\_min\_abs\_a, current\_spike\_sigma, current\_spike\_window, dt\_sec, impossible\_charge, impossible\_voltage\_slew, invalid\_v, invalid\_v\_threshold, key, low\_current\_pair, options, out, power\_spike, power\_spike\_min\_abs\_w, power\_spike\_sigma, power\_spike\_window
              \item 読み取る主なインスタンス属性: 特になし。
              \item 更新する主なインスタンス属性: 特になし。
              \item 明示的に送出する例外: 明示的raiseはない。
              \item 制御構造の規模: 条件分岐 9、ループ 1、try 0
            \end{itemize}
            \paragraph{この定義を読むためのPython構文}
            \begin{itemize}[leftmargin=1.6em]
            \item `->`以後は戻り値の型注釈であり、実行時の値を自動変換する命令ではない。
            \end{itemize}
            \paragraph{上から順の処理}
            \begin{enumerate}[leftmargin=1.8em]
            \item options に options if isinstance(options, dict) else \{\} の結果を代入する。
\item sensor\_cfg に options.get('sensor\_filter', \{\}) if isinstance(options.get('sensor\_filter', \{\}), dict) else \{\} の結果を代入する。
\item 条件 sensor\_cfg.get('enabled', True) is False を判定し、真なら内部処理を行う。
\item   work を返す。
\item out に work.copy() の結果を代入する。
\item ('exclude\_power\_fit', 'exclude\_voltage\_fit', 'exclude\_weather\_fit') を順に走査し、各要素を key に入れて処理する。
\item   条件 key not in out.columns を判定し、真なら内部処理を行う。
\item     out[key] に False の結果を代入する。
\item   out[key] に out[key].fillna(False).astype(bool) の結果を代入する。
\item 条件 'exclude\_reason' not in out.columns を判定し、真なら内部処理を行う。
\item   out['exclude\_reason'] に '' の結果を代入する。
\item invalid\_v\_threshold に float(sensor\_cfg.get('invalid\_voltage\_threshold\_v', INVALID\_PACK\_VOLTAGE\_MIN\_V)) の結果を代入する。
\item current\_limit\_margin\_a に float(sensor\_cfg.get('charge\_current\_limit\_margin\_a', 2.0)) の結果を代入する。
\item current\_spike\_window に int(sensor\_cfg.get('current\_spike\_window', 9)) の結果を代入する。
\item current\_spike\_sigma に float(sensor\_cfg.get('current\_spike\_sigma', 4.0)) の結果を代入する。
\item current\_spike\_min\_abs\_a に float(sensor\_cfg.get('current\_spike\_min\_abs\_a', 4.0)) の結果を代入する。
\item power\_spike\_window に int(sensor\_cfg.get('power\_spike\_window', 9)) の結果を代入する。
\item power\_spike\_sigma に float(sensor\_cfg.get('power\_spike\_sigma', 4.0)) の結果を代入する。
\item power\_spike\_min\_abs\_w に float(sensor\_cfg.get('power\_spike\_min\_abs\_w', 400.0)) の結果を代入する。
\item voltage\_slew\_max\_vps に float(sensor\_cfg.get('voltage\_slew\_max\_vps', 0.2)) の結果を代入する。
\item voltage\_slew\_low\_current\_a に float(sensor\_cfg.get('voltage\_slew\_low\_current\_a', 2.5)) の結果を代入する。
\item voltage\_spike\_window に int(sensor\_cfg.get('voltage\_spike\_window', 9)) の結果を代入する。
\item voltage\_spike\_sigma に float(sensor\_cfg.get('voltage\_spike\_sigma', 4.0)) の結果を代入する。
\item voltage\_spike\_min\_abs\_v に float(sensor\_cfg.get('voltage\_spike\_min\_abs\_v', 2.0)) の結果を代入する。
\item invalid\_v に np.isfinite(out['battery\_voltage\_v']) \& (out['battery\_voltage\_v'] < invalid\_v\_threshold) の結果を代入する。
\item 条件 invalid\_v.any() を判定し、真なら内部処理を行う。
\item   out.loc[invalid\_v, 'exclude\_voltage\_fit'] に True の結果を代入する。
\item   out.loc[invalid\_v, 'exclude\_power\_fit'] に True の結果を代入する。
\item   \_append\_reason\_column(...) を実行する。
\item voltage に pd.to\_numeric(out['battery\_voltage\_v'], errors='coerce') の結果を代入する。
\item current に pd.to\_numeric(out['battery\_current\_a'], errors='coerce') の結果を代入する。
\item dt\_sec に pd.to\_numeric(out['dt\_sec'], errors='coerce').replace(0.0, np.nan) の結果を代入する。
\item voltage\_slew\_vps に voltage.diff().abs() / dt\_sec の結果を代入する。
\item low\_current\_pair に (current.abs() <= voltage\_slew\_low\_current\_a) \& (current.shift(1).abs() <= voltage\_slew\_low\_current\_a) の結果を代入する。
\item impossible\_voltage\_slew に (voltage\_slew\_vps > voltage\_slew\_max\_vps) \& low\_current\_pair \& \textasciitilde{}invalid\_v \& \textasciitilde{}invalid\_v.shift(1, fill\_value=True) の結果を代入する。
\item 条件 impossible\_voltage\_slew.any() を判定し、真なら内部処理を行う。
\item   out.loc[impossible\_voltage\_slew, 'exclude\_voltage\_fit'] に True の結果を代入する。
\item   out.loc[impossible\_voltage\_slew, 'exclude\_power\_fit'] に True の結果を代入する。
\item   \_append\_reason\_column(...) を実行する。
\item voltage\_spike に hampel\_mask(voltage, window=voltage\_spike\_window, n\_sigma=voltage\_spike\_sigma, min\_abs=voltage\_spike\_min\_abs\_v) \& (current.abs() <= voltage\_slew\_low\_current\_a) \& \textasciitilde{}invalid\_v の結果を代入する。
\item 条件 voltage\_spike.any() を判定し、真なら内部処理を行う。
\item   out.loc[voltage\_spike, 'exclude\_voltage\_fit'] に True の結果を代入する。
\item   out.loc[voltage\_spike, 'exclude\_power\_fit'] に True の結果を代入する。
\item   \_append\_reason\_column(...) を実行する。
\item charge\_floor\_a に float(getattr(base\_model.p, 'I\_chg\_min', -16.5)) - current\_limit\_margin\_a の結果を代入する。
\item impossible\_charge に np.isfinite(out['battery\_current\_a']) \& (out['battery\_current\_a'] < charge\_floor\_a) の結果を代入する。
\item 条件 impossible\_charge.any() を判定し、真なら内部処理を行う。
\item   out.loc[impossible\_charge, 'exclude\_power\_fit'] に True の結果を代入する。
\item   \_append\_reason\_column(...) を実行する。
\item current\_spike に hampel\_mask(out['battery\_current\_a'], window=current\_spike\_window, n\_sigma=current\_spike\_sigma, min\_abs=current\_spike\_min\_abs\_a) \& np.isfinite(out['battery\_current\_a']) の結果を代入する。
\item 条件 current\_spike.any() を判定し、真なら内部処理を行う。
\item   out.loc[current\_spike, 'exclude\_power\_fit'] に True の結果を代入する。
\item   \_append\_reason\_column(...) を実行する。
\item power\_spike に hampel\_mask(out['battery\_power\_w\_obs'], window=power\_spike\_window, n\_sigma=power\_spike\_sigma, min\_abs=power\_spike\_min\_abs\_w) \& np.isfinite(out['battery\_power\_w\_obs']) の結果を代入する。
\item 条件 power\_spike.any() を判定し、真なら内部処理を行う。
\item   out.loc[power\_spike, 'exclude\_power\_fit'] に True の結果を代入する。
\item   \_append\_reason\_column(...) を実行する。
\item out を返す。
            \end{enumerate}
            \paragraph{代表コード断片}
            \begin{lstlisting}[language=Python]
def apply_sensor_quality_annotations(
    work: pd.DataFrame,
    *,
    base_model,
    options: dict | None = None,
) -> pd.DataFrame:
    options = options if isinstance(options, dict) else {}
    sensor_cfg = options.get("sensor_filter", {}) if isinstance(options.get("sensor_filter", {}), dict) else {}
    if sensor_cfg.get("enabled", True) is False:
        return work

    out = work.copy()
    for key in ("exclude_power_fit", "exclude_voltage_fit", "exclude_weather_fit"):
        if key not in out.columns:
            out[key] = False
        out[key] = out[key].fillna(False).astype(bool)
    if "exclude_reason" not in out.columns:
        out["exclude_reason"] = ""

    invalid_v_threshold = float(sensor_cfg.get("invalid_voltage_threshold_v", INVALID_PACK_VOLTAGE_MIN_V))
    current_limit_margin_a = float(sensor_cfg.get("charge_current_limit_margin_a", 2.0))
    current_spike_window = int(sensor_cfg.get("current_spike_window", 9))
    current_spike_sigma = float(sensor_cfg.get("current_spike_sigma", 4.0))
    current_spike_min_abs_a = float(sensor_cfg.get("current_spike_min_abs_a", 4.0))
    power_spike_window = int(sensor_cfg.get("power_spike_window", 9))
    power_spike_sigma = float(sensor_cfg.get("power_spike_sigma", 4.0))
    power_spike_min_abs_w = float(sensor_cfg.get("power_spike_min_abs_w", 400.0))
    voltage_slew_max_vps = float(sensor_cfg.get("voltage_slew_max_vps", 0.20))
    voltage_slew_low_current_a = float(sensor_cfg.get("voltage_slew_low_current_a", 2.5))
    voltage_spike_window = int(sensor_cfg.get("voltage_spike_window", 9))
    voltage_spike_sigma = float(sensor_cfg.get("voltage_spike_sigma", 4.0))
    voltage_spike_min_abs_v = float(sensor_cfg.get("voltage_spike_min_abs_v", 2.0))

    invalid_v = np.isfinite(out["battery_voltage_v"]) & (out["battery_voltage_v"] < invalid_v_threshold)
    if invalid_v.any():
...
\end{lstlisting}

\subsection{L520 関数 polarization\_current\_trace}
            \begin{itemize}[leftmargin=1.6em]
              \item 行範囲: L520--L541
              \item 所属: module直下
              \item 定義: \path|polarization_current_trace(replay: pd.DataFrame, *, current_column: str, tau_sec: float) -> np.ndarray|
              \item docstring: Return the 1-RC branch current state immediately before each sample.
              \item このブロックが直接呼ぶ主な関数・メソッド: diff, exp, float, isfinite, len, max, min, range, replay\_segment\_start\_mask, to\_datetime, to\_numeric, to\_numpy
              \item 戻り値の要点: state
              \item 主なローカル代入名: alpha, current, dt, dt\_sec, idx, segment\_starts, state, tau, time\_utc
              \item 読み取る主なインスタンス属性: 特になし。
              \item 更新する主なインスタンス属性: 特になし。
              \item 明示的に送出する例外: 明示的raiseはない。
              \item 制御構造の規模: 条件分岐 1、ループ 1、try 0
            \end{itemize}
            \paragraph{この定義を読むためのPython構文}
            \begin{itemize}[leftmargin=1.6em]
            \item `->`以後は戻り値の型注釈であり、実行時の値を自動変換する命令ではない。
            \end{itemize}
            \paragraph{上から順の処理}
            \begin{enumerate}[leftmargin=1.8em]
            \item current に pd.to\_numeric(replay[current\_column], errors='coerce').to\_numpy(dtype=float) の結果を代入する。
\item current に np.where(np.isfinite(current), current, 0.0) の結果を代入する。
\item time\_utc に pd.to\_datetime(replay['time\_utc'], utc=True, errors='coerce') の結果を代入する。
\item dt\_sec に time\_utc.diff().dt.total\_seconds().to\_numpy(dtype=float) の結果を代入する。
\item segment\_starts に replay\_segment\_start\_mask(replay) の結果を代入する。
\item state に np.zeros(len(replay), dtype=float) の結果を代入する。
\item tau に max(float(tau\_sec), 1e-06) の結果を代入する。
\item range(1, len(replay)) を順に走査し、各要素を idx に入れて処理する。
\item   条件 segment\_starts[idx] or not np.isfinite(dt\_sec[idx]) or dt\_sec[idx] <= 0.0 を判定し、真なら内部処理を行う。
\item     state[idx] に 0.0 の結果を代入する。
\item     Continue 文を実行する。
\item   dt に min(float(dt\_sec[idx]), 60.0) の結果を代入する。
\item   alpha に math.exp(-dt / tau) の結果を代入する。
\item   state[idx] に alpha * state[idx - 1] + (1.0 - alpha) * current[idx - 1] の結果を代入する。
\item state を返す。
            \end{enumerate}
            \paragraph{代表コード断片}
            \begin{lstlisting}[language=Python]
def polarization_current_trace(
    replay: pd.DataFrame,
    *,
    current_column: str,
    tau_sec: float,
) -> np.ndarray:
    """Return the 1-RC branch current state immediately before each sample."""
    current = pd.to_numeric(replay[current_column], errors="coerce").to_numpy(dtype=float)
    current = np.where(np.isfinite(current), current, 0.0)
    time_utc = pd.to_datetime(replay["time_utc"], utc=True, errors="coerce")
    dt_sec = time_utc.diff().dt.total_seconds().to_numpy(dtype=float)
    segment_starts = replay_segment_start_mask(replay)
    state = np.zeros(len(replay), dtype=float)
    tau = max(float(tau_sec), 1.0e-6)
    for idx in range(1, len(replay)):
        if segment_starts[idx] or not np.isfinite(dt_sec[idx]) or dt_sec[idx] <= 0.0:
            state[idx] = 0.0
            continue
        dt = min(float(dt_sec[idx]), 60.0)
        alpha = math.exp(-dt / tau)
        state[idx] = alpha * state[idx - 1] + (1.0 - alpha) * current[idx - 1]
    return state
\end{lstlisting}

\subsection{L544 関数 fit\_battery\_polarization}
            \begin{itemize}[leftmargin=1.6em]
              \item 行範囲: L544--L679
              \item 所属: module直下
              \item 定義: \path|fit_battery_polarization(replay: pd.DataFrame) -> Dict[str, Any]|
              \item docstring: Fit a bounded one-RC branch and gate it on the last independent day.
              \item このブロックが直接呼ぶ主な関数・メソッド: Series, astype, bool, clip, dot, fillna, float, geomspace, get, int, isfinite, len
              \item 戻り値の要点: \{'adopted': adopted, 'reason': 'bounded\_1rc\_improves\_training\_and\_last\_day\_holdout\_rmse' if adopted else 'training\_or\_last\_day\_holdout\_gate\_failed', 'sample\_count': int(base\_valid.sum()), 'training\_sample\_count': int(training\_valid.sum()), 'validation\_sample\_count': int(validation\_valid.sum()), 'holdout\_group': holdout\_group, 'r\_polarization\_ohm': float(best['r\_polarization\_ohm'] if adopted else 0.0), 'tau\_sec': float(best['tau\_sec']), 'rmse\_before\_v': rmse\_before, 'rmse\_after\_v': float(best['rmse\_after\_v'] if adopted else rmse\_before), 'rmse\_improvement\_v': float(improvement if adopted else 0.0), 'validation\_rmse\_before\_v': validation\_before, 'validation\_rmse\_after\_v': validation\_after, 'validation\_rmse\_ratio': validation\_ratio, 'validation\_rmse\_ratio\_max': 1.0, 'method': 'bounded deterministic tau grid with closed-form least-squares Rp on earlier race days; last race day is an untouched adoption holdout'\} / \{'adopted': False, 'reason': 'insufficient\_valid\_voltage\_samples', 'sample\_count': int(base\_valid.sum()), 'r\_polarization\_ohm': 0.0, 'tau\_sec': 60.0, 'rmse\_before\_v': float('nan'), 'rmse\_after\_v': float('nan')\} / \{'adopted': False, 'reason': 'no\_independent\_day\_holdout', 'sample\_count': int(base\_valid.sum()), 'training\_sample\_count': 0, 'validation\_sample\_count': 0, 'r\_polarization\_ohm': 0.0, 'tau\_sec': 60.0, 'rmse\_before\_v': float(np.sqrt(np.mean(residual[base\_valid] ** 2))), 'rmse\_after\_v': float(np.sqrt(np.mean(residual[base\_valid] ** 2)))\} / \{'adopted': False, 'reason': 'insufficient\_independent\_day\_holdout\_samples', 'sample\_count': int(base\_valid.sum()), 'training\_sample\_count': int(training\_valid.sum()), 'validation\_sample\_count': int(validation\_valid.sum()), 'holdout\_group': holdout\_group, 'r\_polarization\_ohm': 0.0, 'tau\_sec': 60.0, 'rmse\_before\_v': float(np.sqrt(np.mean(residual[base\_valid] ** 2))), 'rmse\_after\_v': float(np.sqrt(np.mean(residual[base\_valid] ** 2)))\}
              \item 主なローカル代入名: adopted, base\_valid, best, best\_state, corrected, denom, excluded, groups, holdout\_group, improvement, residual, rmse, rmse\_before, rp\_ohm, state, tau\_sec, time\_local, training\_valid, unique\_groups, valid
              \item 読み取る主なインスタンス属性: 特になし。
              \item 更新する主なインスタンス属性: 特になし。
              \item 明示的に送出する例外: 明示的raiseはない。
              \item 制御構造の規模: 条件分岐 7、ループ 1、try 0
            \end{itemize}
            \paragraph{この定義を読むためのPython構文}
            \begin{itemize}[leftmargin=1.6em]
            \item `->`以後は戻り値の型注釈であり、実行時の値を自動変換する命令ではない。
\item 内包表記は、反復・条件・値生成を一つの式で記述してcollectionを作る。
            \end{itemize}
            \paragraph{上から順の処理}
            \begin{enumerate}[leftmargin=1.8em]
            \item residual に (pd.to\_numeric(replay['battery\_voltage\_v\_obs'], errors='coerce') - pd.to\_numeric(replay['battery\_voltage\_v\_pred'], errors='coerce')).to\_numpy(dtype=float) の結果を代入する。
\item excluded に replay.get('exclude\_voltage\_fit', pd.Series(False, index=replay.index)).fillna(True).astype(bool).to\_numpy() の結果を代入する。
\item base\_valid に np.isfinite(residual) \& \textasciitilde{}excluded の結果を代入する。
\item 条件 int(base\_valid.sum()) < 500 を判定し、真なら内部処理を行う。
\item   \{'adopted': False, 'reason': 'insufficient\_valid\_voltage\_samples', 'sample\_count': int(base\_valid.sum()), 'r\_polarization\_ohm': 0.0, 'tau\_sec': 60.0, 'rmse\_before\_v': float('nan'), 'rmse\_after\_v': float('nan')\} を返す。
\item 条件 'day' in replay.columns を判定し、真なら内部処理を行う。
\item   groups に pd.to\_numeric(replay['day'], errors='coerce').to\_numpy(dtype=float) の結果を代入する。
\item   上の条件が偽の場合:
\item   time\_local に pd.to\_datetime(replay['time\_utc'], utc=True, errors='coerce').dt.tz\_convert(TIMEZONE\_LOCAL) の結果を代入する。
\item   groups に time\_local.dt.strftime('\%Y\%m\%d').astype(float).to\_numpy() の結果を代入する。
\item unique\_groups に sorted((float(value) for value in np.unique(groups[base\_valid \& np.isfinite(groups)]))) の結果を代入する。
\item 条件 len(unique\_groups) < 2 を判定し、真なら内部処理を行う。
\item   \{'adopted': False, 'reason': 'no\_independent\_day\_holdout', 'sample\_count': int(base\_valid.sum()), 'training\_sample\_count': 0, 'validation\_sample\_count': 0, 'r\_polarization\_ohm': 0.0, 'tau\_sec': 60.0, 'rmse\_before\_v': float(np.sqrt(np.mean(residual[base\_valid] ** 2))), 'rmse\_after\_v': float(np.sqrt(np.mean(residual[base\_valid] ** 2)))\} を返す。
\item holdout\_group に unique\_groups[-1] の結果を代入する。
\item training\_valid に base\_valid \& np.isfinite(groups) \& (groups != holdout\_group) の結果を代入する。
\item validation\_valid に base\_valid \& np.isfinite(groups) \& (groups == holdout\_group) の結果を代入する。
\item 条件 int(training\_valid.sum()) < 500 or int(validation\_valid.sum()) < 100 を判定し、真なら内部処理を行う。
\item   \{'adopted': False, 'reason': 'insufficient\_independent\_day\_holdout\_samples', 'sample\_count': int(base\_valid.sum()), 'training\_sample\_count': int(training\_valid.sum()), 'validation\_sample\_count': int(validation\_valid.sum()), 'holdout\_group': holdout\_group, 'r\_polarization\_ohm': 0.0, 'tau\_sec': 60.0, 'rmse\_before\_v': float(np.sqrt(np.mean(residual[base\_valid] ** 2))), 'rmse\_after\_v': float(np.sqrt(np.mean(residual[base\_valid] ** 2)))\} を返す。
\item best に None を代入する。
\item np.geomspace(10.0, 600.0, 81) を順に走査し、各要素を tau\_sec に入れて処理する。
\item   state に polarization\_current\_trace(replay, current\_column='battery\_current\_a\_obs', tau\_sec=float(tau\_sec)) の結果を代入する。
\item   valid に training\_valid \& np.isfinite(state) の結果を代入する。
\item   denom に float(np.dot(state[valid], state[valid])) の結果を代入する。
\item   条件 denom <= 1e-12 を判定し、真なら内部処理を行う。
\item     Continue 文を実行する。
\item   rp\_ohm に float(np.clip(-np.dot(residual[valid], state[valid]) / denom, 0.0, 0.12)) の結果を代入する。
\item   corrected に residual[valid] + rp\_ohm * state[valid] の結果を代入する。
\item   rmse に float(np.sqrt(np.mean(corrected ** 2))) の結果を代入する。
\item   条件 best is None or rmse < best['rmse\_after\_v'] を判定し、真なら内部処理を行う。
\item     best に \{'r\_polarization\_ohm': rp\_ohm, 'tau\_sec': float(tau\_sec), 'rmse\_after\_v': rmse\} の結果を代入する。
\item rmse\_before に float(np.sqrt(np.mean(residual[training\_valid] ** 2))) の結果を代入する。
\item 条件 best is None を判定し、真なら内部処理を行う。
\item   \{'adopted': False, 'reason': 'no\_finite\_candidate', 'sample\_count': int(base\_valid.sum()), 'training\_sample\_count': int(training\_valid.sum()), 'validation\_sample\_count': int(validation\_valid.sum()), 'holdout\_group': holdout\_group, 'r\_polarization\_ohm': 0.0, 'tau\_sec': 60.0, 'rmse\_before\_v': rmse\_before, 'rmse\_after\_v': rmse\_before\} を返す。
\item validation\_before に float(np.sqrt(np.mean(residual[validation\_valid] ** 2))) の結果を代入する。
\item best\_state に polarization\_current\_trace(replay, current\_column='battery\_current\_a\_obs', tau\_sec=float(best['tau\_sec'])) の結果を代入する。
\item validation\_corrected に residual[validation\_valid] + float(best['r\_polarization\_ohm']) * best\_state[validation\_valid] の結果を代入する。
\item validation\_after に float(np.sqrt(np.mean(validation\_corrected ** 2))) の結果を代入する。
\item validation\_ratio に validation\_after / max(validation\_before, 1e-12) の結果を代入する。
\item improvement に rmse\_before - float(best['rmse\_after\_v']) の結果を代入する。
\item adopted に bool(best['r\_polarization\_ohm'] >= 0.001 and improvement >= 0.005 and np.isfinite(validation\_ratio) and (validation\_ratio <= 1.0)) の結果を代入する。
\item \{'adopted': adopted, 'reason': 'bounded\_1rc\_improves\_training\_and\_last\_day\_holdout\_rmse' if adopted else 'training\_or\_last\_day\_holdout\_gate\_failed', 'sample\_count': int(base\_valid.sum()), 'training\_sample\_count': int(training\_valid.sum()), 'validation\_sample\_count': int(validation\_valid.sum()), 'holdout\_group': holdout\_group, 'r\_polarization\_ohm': float(best['r\_polarization\_ohm'] if adopted else 0.0), 'tau\_sec': float(best['tau\_sec']), 'rmse\_before\_v': rmse\_before, 'rmse\_after\_v': float(best['rmse\_after\_v'] if adopted else rmse\_before), 'rmse\_improvement\_v': float(improvement if adopted else 0.0), 'validation\_rmse\_before\_v': validation\_before, 'validation\_rmse\_after\_v': validation\_after, 'validation\_rmse\_ratio': validation\_ratio, 'validation\_rmse\_ratio\_max': 1.0, 'method': 'bounded deterministic tau grid with closed-form least-squares Rp on earlier race days; last race day is an untouched adoption holdout'\} を返す。
            \end{enumerate}
            \paragraph{代表コード断片}
            \begin{lstlisting}[language=Python]
def fit_battery_polarization(replay: pd.DataFrame) -> Dict[str, Any]:
    """Fit a bounded one-RC branch and gate it on the last independent day."""
    residual = (
        pd.to_numeric(replay["battery_voltage_v_obs"], errors="coerce")
        - pd.to_numeric(replay["battery_voltage_v_pred"], errors="coerce")
    ).to_numpy(dtype=float)
    excluded = replay.get("exclude_voltage_fit", pd.Series(False, index=replay.index)).fillna(True).astype(bool).to_numpy()
    base_valid = np.isfinite(residual) & ~excluded
    if int(base_valid.sum()) < 500:
        return {
            "adopted": False,
            "reason": "insufficient_valid_voltage_samples",
            "sample_count": int(base_valid.sum()),
            "r_polarization_ohm": 0.0,
            "tau_sec": 60.0,
            "rmse_before_v": float("nan"),
            "rmse_after_v": float("nan"),
        }

    if "day" in replay.columns:
        groups = pd.to_numeric(replay["day"], errors="coerce").to_numpy(dtype=float)
    else:
        time_local = pd.to_datetime(replay["time_utc"], utc=True, errors="coerce").dt.tz_convert(
            TIMEZONE_LOCAL
        )
        groups = time_local.dt.strftime("%Y%m%d").astype(float).to_numpy()
    unique_groups = sorted(float(value) for value in np.unique(groups[base_valid & np.isfinite(groups)]))
    if len(unique_groups) < 2:
        return {
            "adopted": False,
            "reason": "no_independent_day_holdout",
            "sample_count": int(base_valid.sum()),
            "training_sample_count": 0,
            "validation_sample_count": 0,
            "r_polarization_ohm": 0.0,
...
\end{lstlisting}

\subsection{L682 関数 apply\_battery\_polarization}
            \begin{itemize}[leftmargin=1.6em]
              \item 行範囲: L682--L703
              \item 所属: module直下
              \item 定義: \path|apply_battery_polarization(replay: pd.DataFrame, dynamic_fit: Dict[str, Any], *, current_column: str) -> pd.DataFrame|
              \item docstring: 明示されていない。
              \item このブロックが直接呼ぶ主な関数・メソッド: bool, copy, float, get, polarization\_current\_trace, to\_numeric
              \item 戻り値の要点: out / out
              \item 主なローカル代入名: out, polarization\_v, state
              \item 読み取る主なインスタンス属性: 特になし。
              \item 更新する主なインスタンス属性: 特になし。
              \item 明示的に送出する例外: 明示的raiseはない。
              \item 制御構造の規模: 条件分岐 1、ループ 0、try 0
            \end{itemize}
            \paragraph{この定義を読むためのPython構文}
            \begin{itemize}[leftmargin=1.6em]
            \item `->`以後は戻り値の型注釈であり、実行時の値を自動変換する命令ではない。
            \end{itemize}
            \paragraph{上から順の処理}
            \begin{enumerate}[leftmargin=1.8em]
            \item out に replay.copy() の結果を代入する。
\item out['battery\_voltage\_v\_pred\_static'] に pd.to\_numeric(out['battery\_voltage\_v\_pred'], errors='coerce') の結果を代入する。
\item 条件 not bool(dynamic\_fit.get('adopted', False)) を判定し、真なら内部処理を行う。
\item   out['battery\_polarization\_v'] に 0.0 の結果を代入する。
\item   out を返す。
\item state に polarization\_current\_trace(out, current\_column=current\_column, tau\_sec=float(dynamic\_fit['tau\_sec'])) の結果を代入する。
\item polarization\_v に float(dynamic\_fit['r\_polarization\_ohm']) * state の結果を代入する。
\item out['battery\_polarization\_v'] に polarization\_v の結果を代入する。
\item out['battery\_voltage\_v\_pred'] に out['battery\_voltage\_v\_pred\_static'] - polarization\_v の結果を代入する。
\item out を返す。
            \end{enumerate}
            \paragraph{代表コード断片}
            \begin{lstlisting}[language=Python]
def apply_battery_polarization(
    replay: pd.DataFrame,
    dynamic_fit: Dict[str, Any],
    *,
    current_column: str,
) -> pd.DataFrame:
    out = replay.copy()
    out["battery_voltage_v_pred_static"] = pd.to_numeric(
        out["battery_voltage_v_pred"], errors="coerce"
    )
    if not bool(dynamic_fit.get("adopted", False)):
        out["battery_polarization_v"] = 0.0
        return out
    state = polarization_current_trace(
        out,
        current_column=current_column,
        tau_sec=float(dynamic_fit["tau_sec"]),
    )
    polarization_v = float(dynamic_fit["r_polarization_ohm"]) * state
    out["battery_polarization_v"] = polarization_v
    out["battery_voltage_v_pred"] = out["battery_voltage_v_pred_static"] - polarization_v
    return out
\end{lstlisting}

\subsection{L706 関数 resolve\_fit\_plan}
            \begin{itemize}[leftmargin=1.6em]
              \item 行範囲: L706--L747
              \item 所属: module直下
              \item 定義: \path!resolve_fit_plan(options: dict | None, *, quality: str) -> Dict[str, Any]!
              \item docstring: 明示されていない。
              \item このブロックが直接呼ぶ主な関数・メソッド: dict, float, get, isinstance, items, lower, str, strip
              \item 戻り値の要点: plan
              \item 主なローカル代入名: key\_map, opt\_key, options, out\_key, plan, quality\_norm
              \item 読み取る主なインスタンス属性: 特になし。
              \item 更新する主なインスタンス属性: 特になし。
              \item 明示的に送出する例外: 明示的raiseはない。
              \item 制御構造の規模: 条件分岐 2、ループ 1、try 0
            \end{itemize}
            \paragraph{この定義を読むためのPython構文}
            \begin{itemize}[leftmargin=1.6em]
            \item `->`以後は戻り値の型注釈であり、実行時の値を自動変換する命令ではない。
            \end{itemize}
            \paragraph{上から順の処理}
            \begin{enumerate}[leftmargin=1.8em]
            \item options に options if isinstance(options, dict) else \{\} の結果を代入する。
\item quality\_norm に str(quality or options.get('fit\_quality', 'standard')).strip().lower() の結果を代入する。
\item 条件 quality\_norm not in FIT\_QUALITY\_PRESETS を判定し、真なら内部処理を行う。
\item   quality\_norm に 'standard' の結果を代入する。
\item plan に dict(FIT\_QUALITY\_PRESETS[quality\_norm]) の結果を代入する。
\item key\_map に \{'battery\_restart\_count': 'battery\_restart\_count', 'battery\_maxiter': 'battery\_maxiter', 'motion\_restart\_count': 'motion\_restart\_count', 'motion\_maxiter': 'motion\_maxiter', 'joint\_restart\_count': 'joint\_restart\_count', 'joint\_random\_start\_count': 'joint\_random\_start\_count', 'joint\_local\_topk': 'joint\_local\_topk', 'joint\_maxiter': 'joint\_maxiter', 'fit\_stride': 'fit\_stride', 'allow\_map\_shape\_fit': 'allow\_map\_shape\_fit', 'post\_refine\_enabled': 'post\_refine\_enabled'\} の結果を代入する。
\item key\_map.items() を順に走査し、各要素を (out\_key, opt\_key) に入れて処理する。
\item   条件 opt\_key in options を判定し、真なら内部処理を行う。
\item     plan[out\_key] に options[opt\_key] の結果を代入する。
\item plan['panel\_deployment\_stopped\_speed\_kmh'] に float(options.get('panel\_deployment\_stopped\_speed\_kmh', 2.0)) の結果を代入する。
\item plan['panel\_deployment\_min\_dwell\_sec'] に float(options.get('panel\_deployment\_min\_dwell\_sec', 300.0)) の結果を代入する。
\item plan['panel\_deployment\_max\_sample\_gap\_sec'] に float(options.get('panel\_deployment\_max\_sample\_gap\_sec', 60.0)) の結果を代入する。
\item plan['panel\_control\_stop\_tolerance\_km'] に float(options.get('panel\_control\_stop\_tolerance\_km', 1.0)) の結果を代入する。
\item plan['quality'] に quality\_norm の結果を代入する。
\item plan['terminal\_anchor\_role'] に str(options.get('terminal\_anchor\_role', 'independent\_consensus')) の結果を代入する。
\item plan['sensor\_filter'] に options.get('sensor\_filter', \{\}) の結果を代入する。
\item plan['acceleration\_observation'] に options.get('acceleration\_observation', \{\}) の結果を代入する。
\item plan['grade\_observation'] に options.get('grade\_observation', \{\}) の結果を代入する。
\item plan を返す。
            \end{enumerate}
            \paragraph{代表コード断片}
            \begin{lstlisting}[language=Python]
def resolve_fit_plan(options: dict | None, *, quality: str) -> Dict[str, Any]:
    options = options if isinstance(options, dict) else {}
    quality_norm = str(quality or options.get("fit_quality", "standard")).strip().lower()
    if quality_norm not in FIT_QUALITY_PRESETS:
        quality_norm = "standard"
    plan = dict(FIT_QUALITY_PRESETS[quality_norm])
    key_map = {
        "battery_restart_count": "battery_restart_count",
        "battery_maxiter": "battery_maxiter",
        "motion_restart_count": "motion_restart_count",
        "motion_maxiter": "motion_maxiter",
        "joint_restart_count": "joint_restart_count",
        "joint_random_start_count": "joint_random_start_count",
        "joint_local_topk": "joint_local_topk",
        "joint_maxiter": "joint_maxiter",
        "fit_stride": "fit_stride",
        "allow_map_shape_fit": "allow_map_shape_fit",
        "post_refine_enabled": "post_refine_enabled",
    }
    for out_key, opt_key in key_map.items():
        if opt_key in options:
            plan[out_key] = options[opt_key]
    plan["panel_deployment_stopped_speed_kmh"] = float(
        options.get("panel_deployment_stopped_speed_kmh", 2.0)
    )
    plan["panel_deployment_min_dwell_sec"] = float(
        options.get("panel_deployment_min_dwell_sec", 300.0)
    )
    plan["panel_deployment_max_sample_gap_sec"] = float(
        options.get("panel_deployment_max_sample_gap_sec", 60.0)
    )
    plan["panel_control_stop_tolerance_km"] = float(
        options.get("panel_control_stop_tolerance_km", 1.0)
    )
    plan["quality"] = quality_norm
...
\end{lstlisting}

\subsection{L750 関数 resolve\_identification\_output\_layout}
            \begin{itemize}[leftmargin=1.6em]
              \item 行範囲: L750--L774
              \item 所属: module直下
              \item 定義: \path!resolve_identification_output_layout(package_dir: Path, profile_cfg: dict, *, output_tag_override: str | None = None) -> Dict[str, Path | str]!
              \item docstring: 明示されていない。
              \item このブロックが直接呼ぶ主な関数・メソッド: get, resolve\_relative, str, strip, sub
              \item 戻り値の要点: \{'tag': tag, 'output\_root': output\_root, 'run\_root': run\_root, 'report\_root': report\_root, 'grounded\_maps': run\_root / 'grounded\_base\_maps', 'adopted\_maps': run\_root / 'adopted\_maps'\}
              \item 主なローカル代入名: identification\_cfg, output\_root, raw\_tag, report\_root, run\_root, tag
              \item 読み取る主なインスタンス属性: 特になし。
              \item 更新する主なインスタンス属性: 特になし。
              \item 明示的に送出する例外: 明示的raiseはない。
              \item 制御構造の規模: 条件分岐 1、ループ 0、try 0
            \end{itemize}
            \paragraph{この定義を読むためのPython構文}
            \begin{itemize}[leftmargin=1.6em]
            \item `->`以後は戻り値の型注釈であり、実行時の値を自動変換する命令ではない。
            \end{itemize}
            \paragraph{上から順の処理}
            \begin{enumerate}[leftmargin=1.8em]
            \item identification\_cfg に profile\_cfg.get('identification', \{\}) or \{\} の結果を代入する。
\item output\_root に resolve\_relative(package\_dir, str(identification\_cfg.get('output\_dir', 'outputs/identification') or 'outputs/identification')) の結果を代入する。
\item raw\_tag に output\_tag\_override の結果を代入する。
\item 条件 raw\_tag is None を判定し、真なら内部処理を行う。
\item   raw\_tag に str(identification\_cfg.get('output\_tag', '') or '') の結果を代入する。
\item tag に re.sub('[\textasciicircum{}A-Za-z0-9\_.-]+', '\_', str(raw\_tag).strip()).strip('.\_') の結果を代入する。
\item run\_root に output\_root / 'runs' / tag if tag else output\_root の結果を代入する。
\item report\_root に run\_root / 'reports' if tag else package\_dir / 'outputs' / 'reports' の結果を代入する。
\item \{'tag': tag, 'output\_root': output\_root, 'run\_root': run\_root, 'report\_root': report\_root, 'grounded\_maps': run\_root / 'grounded\_base\_maps', 'adopted\_maps': run\_root / 'adopted\_maps'\} を返す。
            \end{enumerate}
            \paragraph{代表コード断片}
            \begin{lstlisting}[language=Python]
def resolve_identification_output_layout(
    package_dir: Path,
    profile_cfg: dict,
    *,
    output_tag_override: str | None = None,
) -> Dict[str, Path | str]:
    identification_cfg = profile_cfg.get("identification", {}) or {}
    output_root = resolve_relative(
        package_dir,
        str(identification_cfg.get("output_dir", "outputs/identification") or "outputs/identification"),
    )
    raw_tag = output_tag_override
    if raw_tag is None:
        raw_tag = str(identification_cfg.get("output_tag", "") or "")
    tag = re.sub(r"[^A-Za-z0-9_.-]+", "_", str(raw_tag).strip()).strip("._")
    run_root = output_root / "runs" / tag if tag else output_root
    report_root = run_root / "reports" if tag else package_dir / "outputs" / "reports"
    return {
        "tag": tag,
        "output_root": output_root,
        "run_root": run_root,
        "report_root": report_root,
        "grounded_maps": run_root / "grounded_base_maps",
        "adopted_maps": run_root / "adopted_maps",
    }
\end{lstlisting}

\subsection{L777 関数 identification\_profile\_output\_path}
            \begin{itemize}[leftmargin=1.6em]
              \item 行範囲: L777--L786
              \item 所属: module直下
              \item 定義: \path|identification_profile_output_path(canonical_profile: Path, run_output_dir: Path, *, output_tag: str, adopt_profile: bool) -> Path|
              \item docstring: 明示されていない。
              \item このブロックが直接呼ぶ主な関数・メソッド: Path, str, strip
              \item 戻り値の要点: Path(canonical\_profile) / Path(run\_output\_dir) / 'profile\_candidate.yaml'
              \item 主なローカル代入名: 特になし。
              \item 読み取る主なインスタンス属性: 特になし。
              \item 更新する主なインスタンス属性: 特になし。
              \item 明示的に送出する例外: 明示的raiseはない。
              \item 制御構造の規模: 条件分岐 1、ループ 0、try 0
            \end{itemize}
            \paragraph{この定義を読むためのPython構文}
            \begin{itemize}[leftmargin=1.6em]
            \item `->`以後は戻り値の型注釈であり、実行時の値を自動変換する命令ではない。
            \end{itemize}
            \paragraph{上から順の処理}
            \begin{enumerate}[leftmargin=1.8em]
            \item 条件 str(output\_tag).strip() and (not adopt\_profile) を判定し、真なら内部処理を行う。
\item   Path(run\_output\_dir) / 'profile\_candidate.yaml' を返す。
\item Path(canonical\_profile) を返す。
            \end{enumerate}
            \paragraph{代表コード断片}
            \begin{lstlisting}[language=Python]
def identification_profile_output_path(
    canonical_profile: Path,
    run_output_dir: Path,
    *,
    output_tag: str,
    adopt_profile: bool,
) -> Path:
    if str(output_tag).strip() and not adopt_profile:
        return Path(run_output_dir) / "profile_candidate.yaml"
    return Path(canonical_profile)
\end{lstlisting}

\subsection{L789 関数 load\_ocv\_df}
            \begin{itemize}[leftmargin=1.6em]
              \item 行範囲: L789--L794
              \item 所属: module直下
              \item 定義: \path|load_ocv_df(profile_cfg: dict, profile_yaml: Path) -> pd.DataFrame|
              \item docstring: 明示されていない。
              \item このブロックが直接呼ぶ主な関数・メソッド: FileNotFoundError, get, read\_csv, resolve\_relative, str, strip
              \item 戻り値の要点: pd.read\_csv(ocv\_path)
              \item 主なローカル代入名: ocv\_path, raw
              \item 読み取る主なインスタンス属性: 特になし。
              \item 更新する主なインスタンス属性: 特になし。
              \item 明示的に送出する例外: FileNotFoundError('profile.paths.ocv\_soc\_map is required')
              \item 制御構造の規模: 条件分岐 1、ループ 0、try 0
            \end{itemize}
            \paragraph{この定義を読むためのPython構文}
            \begin{itemize}[leftmargin=1.6em]
            \item `->`以後は戻り値の型注釈であり、実行時の値を自動変換する命令ではない。
            \end{itemize}
            \paragraph{上から順の処理}
            \begin{enumerate}[leftmargin=1.8em]
            \item raw に str((profile\_cfg.get('paths', \{\}) or \{\}).get('ocv\_soc\_map', '') or '').strip() の結果を代入する。
\item 条件 not raw を判定し、真なら内部処理を行う。
\item   FileNotFoundError('profile.paths.ocv\_soc\_map is required') を送出する。
\item ocv\_path に resolve\_relative(profile\_yaml.parent, raw) の結果を代入する。
\item pd.read\_csv(ocv\_path) を返す。
            \end{enumerate}
            \paragraph{代表コード断片}
            \begin{lstlisting}[language=Python]
def load_ocv_df(profile_cfg: dict, profile_yaml: Path) -> pd.DataFrame:
    raw = str((profile_cfg.get("paths", {}) or {}).get("ocv_soc_map", "") or "").strip()
    if not raw:
        raise FileNotFoundError("profile.paths.ocv_soc_map is required")
    ocv_path = resolve_relative(profile_yaml.parent, raw)
    return pd.read_csv(ocv_path)
\end{lstlisting}

\subsection{L797 関数 build\_source\_map\_assets}
            \begin{itemize}[leftmargin=1.6em]
              \item 行範囲: L797--L818
              \item 所属: module直下
              \item 定義: \path|build_source_map_assets(profile_cfg: dict, profile_yaml: Path) -> Dict[str, Path]|
              \item docstring: 明示されていない。
              \item このブロックが直接呼ぶ主な関数・メソッド: FileNotFoundError, get, rel, resolve\_relative, str, strip
              \item 戻り値の要点: \{'drive\_eff\_map': rel('drive\_eff\_map'), 'drive\_map\_eco': rel('drive\_map\_eco', paths.get('drive\_eff\_map', '')), 'drive\_map\_power': rel('drive\_map\_power', paths.get('drive\_eff\_map', '')), 'regen\_eff\_map': rel('regen\_eff\_map'), 'regen\_map\_eco': rel('regen\_map\_eco', paths.get('regen\_eff\_map', '')), 'regen\_map\_power': rel('regen\_map\_power', paths.get('regen\_eff\_map', '')), 'rint\_map': rel('rint\_map'), 'panel\_eff\_map': rel('panel\_eff\_map'), 'mppt\_eff\_map': rel('mppt\_eff\_map'), 'ocv\_soc\_map': rel('ocv\_soc\_map')\} / resolve\_relative(base\_dir, raw)
              \item 主なローカル代入名: base\_dir, paths, raw
              \item 読み取る主なインスタンス属性: 特になし。
              \item 更新する主なインスタンス属性: 特になし。
              \item 明示的に送出する例外: FileNotFoundError(f'profile.paths.\{key\} is required')
              \item 制御構造の規模: 条件分岐 1、ループ 0、try 0
            \end{itemize}
            \paragraph{この定義を読むためのPython構文}
            \begin{itemize}[leftmargin=1.6em]
            \item `->`以後は戻り値の型注釈であり、実行時の値を自動変換する命令ではない。
            \end{itemize}
            \paragraph{上から順の処理}
            \begin{enumerate}[leftmargin=1.8em]
            \item base\_dir に profile\_yaml.parent の結果を代入する。
\item paths に profile\_cfg.get('paths', \{\}) or \{\} の結果を代入する。
\item 関数 rel を定義する。
\item \{'drive\_eff\_map': rel('drive\_eff\_map'), 'drive\_map\_eco': rel('drive\_map\_eco', paths.get('drive\_eff\_map', '')), 'drive\_map\_power': rel('drive\_map\_power', paths.get('drive\_eff\_map', '')), 'regen\_eff\_map': rel('regen\_eff\_map'), 'regen\_map\_eco': rel('regen\_map\_eco', paths.get('regen\_eff\_map', '')), 'regen\_map\_power': rel('regen\_map\_power', paths.get('regen\_eff\_map', '')), 'rint\_map': rel('rint\_map'), 'panel\_eff\_map': rel('panel\_eff\_map'), 'mppt\_eff\_map': rel('mppt\_eff\_map'), 'ocv\_soc\_map': rel('ocv\_soc\_map')\} を返す。
            \end{enumerate}
            \paragraph{代表コード断片}
            \begin{lstlisting}[language=Python]
def build_source_map_assets(profile_cfg: dict, profile_yaml: Path) -> Dict[str, Path]:
    base_dir = profile_yaml.parent
    paths = profile_cfg.get("paths", {}) or {}

    def rel(key: str, fallback: str | None = None) -> Path:
        raw = str(paths.get(key, fallback or "") or "").strip()
        if not raw:
            raise FileNotFoundError(f"profile.paths.{key} is required")
        return resolve_relative(base_dir, raw)

    return {
        "drive_eff_map": rel("drive_eff_map"),
        "drive_map_eco": rel("drive_map_eco", paths.get("drive_eff_map", "")),
        "drive_map_power": rel("drive_map_power", paths.get("drive_eff_map", "")),
        "regen_eff_map": rel("regen_eff_map"),
        "regen_map_eco": rel("regen_map_eco", paths.get("regen_eff_map", "")),
        "regen_map_power": rel("regen_map_power", paths.get("regen_eff_map", "")),
        "rint_map": rel("rint_map"),
        "panel_eff_map": rel("panel_eff_map"),
        "mppt_eff_map": rel("mppt_eff_map"),
        "ocv_soc_map": rel("ocv_soc_map"),
    }
\end{lstlisting}

\subsection{L801 関数 build\_source\_map\_assets.rel}
            \begin{itemize}[leftmargin=1.6em]
              \item 行範囲: L801--L805
              \item 所属: \path|build_source_map_assets|
              \item 定義: \path!rel(key: str, fallback: str | None = None) -> Path!
              \item docstring: 明示されていない。
              \item このブロックが直接呼ぶ主な関数・メソッド: FileNotFoundError, get, resolve\_relative, str, strip
              \item 戻り値の要点: resolve\_relative(base\_dir, raw)
              \item 主なローカル代入名: raw
              \item 読み取る主なインスタンス属性: 特になし。
              \item 更新する主なインスタンス属性: 特になし。
              \item 明示的に送出する例外: FileNotFoundError(f'profile.paths.\{key\} is required')
              \item 制御構造の規模: 条件分岐 1、ループ 0、try 0
            \end{itemize}
            \paragraph{この定義を読むためのPython構文}
            \begin{itemize}[leftmargin=1.6em]
            \item `->`以後は戻り値の型注釈であり、実行時の値を自動変換する命令ではない。
            \end{itemize}
            \paragraph{上から順の処理}
            \begin{enumerate}[leftmargin=1.8em]
            \item raw に str(paths.get(key, fallback or '') or '').strip() の結果を代入する。
\item 条件 not raw を判定し、真なら内部処理を行う。
\item   FileNotFoundError(f'profile.paths.\{key\} is required') を送出する。
\item resolve\_relative(base\_dir, raw) を返す。
            \end{enumerate}
            \paragraph{代表コード断片}
            \begin{lstlisting}[language=Python]
    def rel(key: str, fallback: str | None = None) -> Path:
        raw = str(paths.get(key, fallback or "") or "").strip()
        if not raw:
            raise FileNotFoundError(f"profile.paths.{key} is required")
        return resolve_relative(base_dir, raw)
\end{lstlisting}

\subsection{L821 関数 apply\_actual\_event\_annotations}
            \begin{itemize}[leftmargin=1.6em]
              \item 行範囲: L821--L860
              \item 所属: module直下
              \item 定義: \path!apply_actual_event_annotations(work: pd.DataFrame, payload: dict | None) -> pd.DataFrame!
              \item docstring: 明示されていない。
              \item このブロックが直接呼ぶ主な関数・メソッド: Timestamp, \_append\_reason\_column, any, astype, bool, copy, fillna, get, hasattr, isinstance, str, strip
              \item 戻り値の要点: out / work / work
              \item 主なローカル代入名: end\_utc, event, events, fit\_flags, key, mask, out, raw\_end, raw\_start, reason, start\_utc, time\_utc, touched, tz\_name
              \item 読み取る主なインスタンス属性: 特になし。
              \item 更新する主なインスタンス属性: 特になし。
              \item 明示的に送出する例外: 明示的raiseはない。
              \item 制御構造の規模: 条件分岐 9、ループ 3、try 1
            \end{itemize}
            \paragraph{この定義を読むためのPython構文}
            \begin{itemize}[leftmargin=1.6em]
            \item `->`以後は戻り値の型注釈であり、実行時の値を自動変換する命令ではない。
\item try文は例外が起き得る処理と、異常時または終了時の経路を分ける。
            \end{itemize}
            \paragraph{上から順の処理}
            \begin{enumerate}[leftmargin=1.8em]
            \item 条件 not isinstance(payload, dict) を判定し、真なら内部処理を行う。
\item   work を返す。
\item events に payload.get('events', payload) の結果を代入する。
\item 条件 not isinstance(events, list) or not events を判定し、真なら内部処理を行う。
\item   work を返す。
\item out に work.copy() の結果を代入する。
\item ('exclude\_power\_fit', 'exclude\_voltage\_fit', 'exclude\_weather\_fit') を順に走査し、各要素を key に入れて処理する。
\item   条件 key not in out.columns を判定し、真なら内部処理を行う。
\item     out[key] に False の結果を代入する。
\item   out[key] に out[key].fillna(False).astype(bool) の結果を代入する。
\item 条件 'exclude\_reason' not in out.columns を判定し、真なら内部処理を行う。
\item   out['exclude\_reason'] に '' の結果を代入する。
\item time\_utc に pd.to\_datetime(out['time\_utc'], format='mixed', utc=True, errors='coerce') の結果を代入する。
\item events を順に走査し、各要素を event に入れて処理する。
\item   条件 not isinstance(event, dict) を判定し、真なら内部処理を行う。
\item     Continue 文を実行する。
\item   raw\_start に str(event.get('start\_local', '') or '').strip() の結果を代入する。
\item   raw\_end に str(event.get('end\_local', '') or '').strip() の結果を代入する。
\item   条件 not raw\_start or not raw\_end を判定し、真なら内部処理を行う。
\item     Continue 文を実行する。
\item   tz\_name に str(event.get('timezone', TIMEZONE\_LOCAL.key if hasattr(TIMEZONE\_LOCAL, 'key') else 'Australia/Darwin') or 'Australia/Darwin') の結果を代入する。
\item   例外処理を伴う try ブロックを実行する。
\item     start\_utc に pd.Timestamp(raw\_start).tz\_localize(tz\_name).tz\_convert('UTC') の結果を代入する。
\item     end\_utc に pd.Timestamp(raw\_end).tz\_localize(tz\_name).tz\_convert('UTC') の結果を代入する。
\item     Exceptionを捕捉した場合:
\item     Continue 文を実行する。
\item   mask に (time\_utc >= start\_utc) \& (time\_utc <= end\_utc) の結果を代入する。
\item   条件 not mask.any() を判定し、真なら内部処理を行う。
\item     Continue 文を実行する。
\item   fit\_flags に event.get('fit\_flags', \{\}) if isinstance(event.get('fit\_flags', \{\}), dict) else \{\} の結果を代入する。
\item   reason に str(event.get('label', event.get('kind', 'actual\_event')) or 'actual\_event') の結果を代入する。
\item   touched に False の結果を代入する。
\item   ('exclude\_power\_fit', 'exclude\_voltage\_fit', 'exclude\_weather\_fit') を順に走査し、各要素を key に入れて処理する。
\item     条件 bool(fit\_flags.get(key, False)) を判定し、真なら内部処理を行う。
\item       out.loc[mask, key] に True の結果を代入する。
\item       touched に True の結果を代入する。
\item   条件 touched を判定し、真なら内部処理を行う。
\item     \_append\_reason\_column(...) を実行する。
\item out を返す。
            \end{enumerate}
            \paragraph{代表コード断片}
            \begin{lstlisting}[language=Python]
def apply_actual_event_annotations(work: pd.DataFrame, payload: dict | None) -> pd.DataFrame:
    if not isinstance(payload, dict):
        return work
    events = payload.get("events", payload)
    if not isinstance(events, list) or not events:
        return work
    out = work.copy()
    for key in ("exclude_power_fit", "exclude_voltage_fit", "exclude_weather_fit"):
        if key not in out.columns:
            out[key] = False
        out[key] = out[key].fillna(False).astype(bool)
    if "exclude_reason" not in out.columns:
        out["exclude_reason"] = ""
    time_utc = pd.to_datetime(out["time_utc"], format="mixed", utc=True, errors="coerce")
    for event in events:
        if not isinstance(event, dict):
            continue
        raw_start = str(event.get("start_local", "") or "").strip()
        raw_end = str(event.get("end_local", "") or "").strip()
        if not raw_start or not raw_end:
            continue
        tz_name = str(event.get("timezone", TIMEZONE_LOCAL.key if hasattr(TIMEZONE_LOCAL, "key") else "Australia/Darwin") or "Australia/Darwin")
        try:
            start_utc = pd.Timestamp(raw_start).tz_localize(tz_name).tz_convert("UTC")
            end_utc = pd.Timestamp(raw_end).tz_localize(tz_name).tz_convert("UTC")
        except Exception:
            continue
        mask = (time_utc >= start_utc) & (time_utc <= end_utc)
        if not mask.any():
            continue
        fit_flags = event.get("fit_flags", {}) if isinstance(event.get("fit_flags", {}), dict) else {}
        reason = str(event.get("label", event.get("kind", "actual_event")) or "actual_event")
        touched = False
        for key in ("exclude_power_fit", "exclude_voltage_fit", "exclude_weather_fit"):
            if bool(fit_flags.get(key, False)):
...
\end{lstlisting}

\subsection{L863 関数 truncate\_at\_retire\_event}
            \begin{itemize}[leftmargin=1.6em]
              \item 行範囲: L863--L889
              \item 所属: module直下
              \item 定義: \path!truncate_at_retire_event(work: pd.DataFrame, payload: dict | None) -> pd.DataFrame!
              \item docstring: End historical replay at the first authoritative retirement timestamp.
              \item このブロックが直接呼ぶ主な関数・メソッド: Timestamp, ValueError, append, copy, get, isinstance, isoformat, lower, min, reset\_index, str, strip
              \item 戻り値の要点: retained.reset\_index(drop=True) / work / work / work
              \item 主なローカル代入名: cutoff, cutoffs, event, events, raw\_start, retained, time\_utc, tz\_name
              \item 読み取る主なインスタンス属性: 特になし。
              \item 更新する主なインスタンス属性: 特になし。
              \item 明示的に送出する例外: ValueError(f'retire-event cutoff \{cutoff.isoformat()\} removed the complete replay')
              \item 制御構造の規模: 条件分岐 6、ループ 1、try 1
            \end{itemize}
            \paragraph{この定義を読むためのPython構文}
            \begin{itemize}[leftmargin=1.6em]
            \item `->`以後は戻り値の型注釈であり、実行時の値を自動変換する命令ではない。
\item try文は例外が起き得る処理と、異常時または終了時の経路を分ける。
            \end{itemize}
            \paragraph{上から順の処理}
            \begin{enumerate}[leftmargin=1.8em]
            \item 条件 not isinstance(payload, dict) を判定し、真なら内部処理を行う。
\item   work を返す。
\item events に payload.get('events', payload) の結果を代入する。
\item 条件 not isinstance(events, list) を判定し、真なら内部処理を行う。
\item   work を返す。
\item cutoffs に [] を代入する。
\item events を順に走査し、各要素を event に入れて処理する。
\item   条件 not isinstance(event, dict) or str(event.get('kind', '')).strip().lower() != 'retire\_event' を判定し、真なら内部処理を行う。
\item     Continue 文を実行する。
\item   raw\_start に str(event.get('start\_local', '') or '').strip() の結果を代入する。
\item   条件 not raw\_start を判定し、真なら内部処理を行う。
\item     Continue 文を実行する。
\item   tz\_name に str(event.get('timezone', 'Australia/Darwin') or 'Australia/Darwin') の結果を代入する。
\item   例外処理を伴う try ブロックを実行する。
\item     cutoffs.append(...) を実行する。
\item     Exceptionを捕捉した場合:
\item     Continue 文を実行する。
\item 条件 not cutoffs を判定し、真なら内部処理を行う。
\item   work を返す。
\item cutoff に min(cutoffs) の結果を代入する。
\item time\_utc に pd.to\_datetime(work['time\_utc'], format='mixed', utc=True, errors='coerce') の結果を代入する。
\item retained に work.loc[time\_utc <= cutoff].copy() の結果を代入する。
\item 条件 retained.empty を判定し、真なら内部処理を行う。
\item   ValueError(f'retire-event cutoff \{cutoff.isoformat()\} removed the complete replay') を送出する。
\item retained.reset\_index(drop=True) を返す。
            \end{enumerate}
            \paragraph{代表コード断片}
            \begin{lstlisting}[language=Python]
def truncate_at_retire_event(work: pd.DataFrame, payload: dict | None) -> pd.DataFrame:
    """End historical replay at the first authoritative retirement timestamp."""
    if not isinstance(payload, dict):
        return work
    events = payload.get("events", payload)
    if not isinstance(events, list):
        return work
    cutoffs: list[pd.Timestamp] = []
    for event in events:
        if not isinstance(event, dict) or str(event.get("kind", "")).strip().lower() != "retire_event":
            continue
        raw_start = str(event.get("start_local", "") or "").strip()
        if not raw_start:
            continue
        tz_name = str(event.get("timezone", "Australia/Darwin") or "Australia/Darwin")
        try:
            cutoffs.append(pd.Timestamp(raw_start).tz_localize(tz_name).tz_convert("UTC"))
        except Exception:
            continue
    if not cutoffs:
        return work
    cutoff = min(cutoffs)
    time_utc = pd.to_datetime(work["time_utc"], format="mixed", utc=True, errors="coerce")
    retained = work.loc[time_utc <= cutoff].copy()
    if retained.empty:
        raise ValueError(f"retire-event cutoff {cutoff.isoformat()} removed the complete replay")
    return retained.reset_index(drop=True)
\end{lstlisting}

\subsection{L892 関数 normalize\_generic\_log}
            \begin{itemize}[leftmargin=1.6em]
              \item 行範囲: L892--L991
              \item 所属: module直下
              \item 定義: \path!normalize_generic_log(log_csv: Path, *, actual_event_payload: dict | None = None, base_model = None, options: dict | None = None) -> pd.DataFrame!
              \item docstring: 明示されていない。
              \item このブロックが直接呼ぶ主な関数・メソッド: ValueError, apply\_actual\_event\_annotations, apply\_sensor\_quality\_annotations, astype, clip, copy, diff, dropna, extend, fillna, float, items
              \item 戻り値の要点: work
              \item 主なローカル代入名: accel, column, df, dt, key, numeric\_cols, required\_defaults, segment\_start, speed\_ms, time\_gap, value, work
              \item 読み取る主なインスタンス属性: 特になし。
              \item 更新する主なインスタンス属性: 特になし。
              \item 明示的に送出する例外: ValueError('normalized replay log must contain time\_utc')
              \item 制御構造の規模: 条件分岐 12、ループ 3、try 0
            \end{itemize}
            \paragraph{この定義を読むためのPython構文}
            \begin{itemize}[leftmargin=1.6em]
            \item `->`以後は戻り値の型注釈であり、実行時の値を自動変換する命令ではない。
\item 内包表記は、反復・条件・値生成を一つの式で記述してcollectionを作る。
            \end{itemize}
            \paragraph{上から順の処理}
            \begin{enumerate}[leftmargin=1.8em]
            \item df に pd.read\_csv(log\_csv, low\_memory=False) の結果を代入する。
\item 条件 'time\_utc' not in df.columns を判定し、真なら内部処理を行う。
\item   ValueError('normalized replay log must contain time\_utc') を送出する。
\item work に df.copy() の結果を代入する。
\item work['time\_utc'] に pd.to\_datetime(work['time\_utc'], format='mixed', utc=True) の結果を代入する。
\item work に work.sort\_values('time\_utc').reset\_index(drop=True) の結果を代入する。
\item required\_defaults に \{'s\_km': np.nan, 'speed\_kmh': 0.0, 'slope\_pct': 0.0, 'route\_heading\_deg': 0.0, 'headwind\_archive\_ms': 0.0, 'GHI\_archive': 0.0, 'Tamb\_archive\_C': 25.0, 'solar\_power\_w\_obs': 0.0, 'battery\_power\_w\_obs': 0.0, 'battery\_current\_a': 0.0, 'battery\_voltage\_v': np.nan, 'lat': np.nan, 'lon': np.nan, 'alt\_m': 0.0\} の結果を代入する。
\item required\_defaults.items() を順に走査し、各要素を (key, value) に入れて処理する。
\item   条件 key not in work.columns を判定し、真なら内部処理を行う。
\item     work[key] に value の結果を代入する。
\item 条件 'dt\_sec' not in work.columns を判定し、真なら内部処理を行う。
\item   dt に work['time\_utc'].diff().dt.total\_seconds().fillna(0.0) の結果を代入する。
\item   条件 len(dt) >= 2 and dt.iloc[0] <= 0.0 を判定し、真なら内部処理を行う。
\item     dt.iloc[0] に float(np.nanmedian(dt.iloc[1:].replace(0.0, np.nan).dropna())) if len(dt.iloc[1:].dropna()) else 5.0 の結果を代入する。
\item   work['dt\_sec'] に dt.clip(lower=0.0).fillna(5.0) の結果を代入する。
\item 条件 'time\_local' not in work.columns を判定し、真なら内部処理を行う。
\item   work['time\_local'] に work['time\_utc'].dt.tz\_convert(TIMEZONE\_LOCAL).astype(str) の結果を代入する。
\item 条件 'day' not in work.columns を判定し、真なら内部処理を行う。
\item   work['day'] に (work['time\_utc'].dt.normalize() - work['time\_utc'].dt.normalize().min()).dt.days + 1 の結果を代入する。
\item ('exclude\_power\_fit', 'exclude\_voltage\_fit', 'exclude\_weather\_fit') を順に走査し、各要素を key に入れて処理する。
\item   条件 key not in work.columns を判定し、真なら内部処理を行う。
\item     work[key] に False の結果を代入する。
\item   work[key] に work[key].fillna(False).astype(bool) の結果を代入する。
\item 条件 'exclude\_reason' not in work.columns を判定し、真なら内部処理を行う。
\item   work['exclude\_reason'] に '' の結果を代入する。
\item 条件 'GHI\_effective' not in work.columns を判定し、真なら内部処理を行う。
\item   work['GHI\_effective'] に work['GHI\_archive'] の結果を代入する。
\item 条件 'Tcell\_effective\_C' not in work.columns を判定し、真なら内部処理を行う。
\item   work['Tcell\_effective\_C'] に work['Tamb\_archive\_C'] + 0.03 * pd.to\_numeric(work['GHI\_effective'], errors='coerce').fillna(0.0) の結果を代入する。
\item 条件 'headwind\_effective\_ms' not in work.columns を判定し、真なら内部処理を行う。
\item   work['headwind\_effective\_ms'] に work['headwind\_archive\_ms'] の結果を代入する。
\item numeric\_cols に ['s\_km', 'speed\_kmh', 'slope\_pct', 'route\_heading\_deg', 'headwind\_archive\_ms', 'headwind\_effective\_ms', 'GHI\_archive', 'GHI\_effective', 'Tamb\_archive\_C', 'Tcell\_effective\_C', 'solar\_power\_w\_obs', 'battery\_power\_w\_obs', 'battery\_current\_a', 'battery\_voltage\_v', 'dt\_sec', 'lat', 'lon', 'alt\_m'] の結果を代入する。
\item numeric\_cols.extend(...) を実行する。
\item numeric\_cols を順に走査し、各要素を key に入れて処理する。
\item   work[key] に pd.to\_numeric(work[key], errors='coerce') の結果を代入する。
\item speed\_ms に work['speed\_kmh'].fillna(0.0) / 3.6 の結果を代入する。
\item dt に work['dt\_sec'].replace(0.0, np.nan) の結果を代入する。
\item accel に speed\_ms.diff() / dt の結果を代入する。
\item segment\_start に work['day'].ne(work['day'].shift(1)) の結果を代入する。
\item time\_gap に work['time\_utc'].diff().dt.total\_seconds().fillna(0.0) > 7200.0 の結果を代入する。
\item accel.loc[segment\_start | time\_gap] に 0.0 の結果を代入する。
\item work['accel\_ms2'] に accel.replace([np.inf, -np.inf], np.nan).rolling(5, center=True, min\_periods=1).median().fillna(0.0).clip(lower=-1.5, upper=1.5) の結果を代入する。
\item work に apply\_actual\_event\_annotations(work, actual\_event\_payload) の結果を代入する。
\item work に truncate\_at\_retire\_event(work, actual\_event\_payload) の結果を代入する。
\item 条件 base\_model is not None を判定し、真なら内部処理を行う。
\item   work に apply\_sensor\_quality\_annotations(work, base\_model=base\_model, options=options) の結果を代入する。
\item work を返す。
            \end{enumerate}
            \paragraph{代表コード断片}
            \begin{lstlisting}[language=Python]
def normalize_generic_log(
    log_csv: Path,
    *,
    actual_event_payload: dict | None = None,
    base_model=None,
    options: dict | None = None,
) -> pd.DataFrame:
    df = pd.read_csv(log_csv, low_memory=False)
    if "time_utc" not in df.columns:
        raise ValueError("normalized replay log must contain time_utc")
    work = df.copy()
    work["time_utc"] = pd.to_datetime(work["time_utc"], format="mixed", utc=True)
    work = work.sort_values("time_utc").reset_index(drop=True)

    required_defaults = {
        "s_km": np.nan,
        "speed_kmh": 0.0,
        "slope_pct": 0.0,
        "route_heading_deg": 0.0,
        "headwind_archive_ms": 0.0,
        "GHI_archive": 0.0,
        "Tamb_archive_C": 25.0,
        "solar_power_w_obs": 0.0,
        "battery_power_w_obs": 0.0,
        "battery_current_a": 0.0,
        "battery_voltage_v": np.nan,
        "lat": np.nan,
        "lon": np.nan,
        "alt_m": 0.0,
    }
    for key, value in required_defaults.items():
        if key not in work.columns:
            work[key] = value

    if "dt_sec" not in work.columns:
...
\end{lstlisting}

\subsection{L994 関数 build\_terminal\_anchor}
            \begin{itemize}[leftmargin=1.6em]
              \item 行範囲: L994--L1026
              \item 所属: module直下
              \item 定義: \path|build_terminal_anchor(log_df: pd.DataFrame, ocv_df: pd.DataFrame, base_model, anchor_km: float) -> Dict[str, float]|
              \item docstring: 明示されていない。
              \item このブロックが直接呼ぶ主な関数・メソッド: abs, asarray, astype, copy, fillna, float, infer\_soc\_from\_loaded\_state, int, isfinite, itertuples, len, median
              \item 戻り値の要点: \{'count': int(len(window)), 's\_km': float(window['s\_km'].median()), 'time\_utc': pd.to\_datetime(window['time\_utc'].iloc[-1], utc=True).strftime('\%Y-\%m-\%dT\%H:\%M:\%SZ'), 'voltage\_v': float(window['battery\_voltage\_v'].median()), 'current\_a': float(window['battery\_current\_a'].median()), 'temp\_c': float(window['Tamb\_archive\_C'].median()), 'soc\_target': float(np.nanmedian(np.asarray(soc\_targets, dtype=float)))\}
              \item 主なローカル代入名: row, soc\_targets, valid, window
              \item 読み取る主なインスタンス属性: 特になし。
              \item 更新する主なインスタンス属性: 特になし。
              \item 明示的に送出する例外: 明示的raiseはない。
              \item 制御構造の規模: 条件分岐 1、ループ 0、try 0
            \end{itemize}
            \paragraph{この定義を読むためのPython構文}
            \begin{itemize}[leftmargin=1.6em]
            \item `->`以後は戻り値の型注釈であり、実行時の値を自動変換する命令ではない。
\item 内包表記は、反復・条件・値生成を一つの式で記述してcollectionを作る。
            \end{itemize}
            \paragraph{上から順の処理}
            \begin{enumerate}[leftmargin=1.8em]
            \item valid に np.isfinite(log\_df['battery\_voltage\_v']) \& np.isfinite(log\_df['battery\_current\_a']) \& np.isfinite(log\_df['Tamb\_archive\_C']) \& np.isfinite(log\_df['s\_km']) \& (log\_df['battery\_voltage\_v'] >= INVALID\_PACK\_VOLTAGE\_MIN\_V) \& \textasciitilde{}log\_df['exclude\_voltage\_fit'].fillna(False).astype(bool) の結果を代入する。
\item window に log\_df.loc[valid \& (np.abs(log\_df['s\_km'] - float(anchor\_km)) <= 1.5)].copy() の結果を代入する。
\item 条件 window.empty を判定し、真なら内部処理を行う。
\item   window に log\_df.loc[valid].tail(6).copy() の結果を代入する。
\item   上の条件が偽の場合:
\item   window に window.tail(min(6, len(window))).copy() の結果を代入する。
\item soc\_targets に [infer\_soc\_from\_loaded\_state(float(row.battery\_voltage\_v), float(row.battery\_current\_a), float(row.Tamb\_archive\_C), ocv\_df, base\_model) for row in window.itertuples(index=False)] の結果を代入する。
\item \{'count': int(len(window)), 's\_km': float(window['s\_km'].median()), 'time\_utc': pd.to\_datetime(window['time\_utc'].iloc[-1], utc=True).strftime('\%Y-\%m-\%dT\%H:\%M:\%SZ'), 'voltage\_v': float(window['battery\_voltage\_v'].median()), 'current\_a': float(window['battery\_current\_a'].median()), 'temp\_c': float(window['Tamb\_archive\_C'].median()), 'soc\_target': float(np.nanmedian(np.asarray(soc\_targets, dtype=float)))\} を返す。
            \end{enumerate}
            \paragraph{代表コード断片}
            \begin{lstlisting}[language=Python]
def build_terminal_anchor(log_df: pd.DataFrame, ocv_df: pd.DataFrame, base_model, anchor_km: float) -> Dict[str, float]:
    valid = (
        np.isfinite(log_df["battery_voltage_v"])
        & np.isfinite(log_df["battery_current_a"])
        & np.isfinite(log_df["Tamb_archive_C"])
        & np.isfinite(log_df["s_km"])
        & (log_df["battery_voltage_v"] >= INVALID_PACK_VOLTAGE_MIN_V)
        & (~log_df["exclude_voltage_fit"].fillna(False).astype(bool))
    )
    window = log_df.loc[valid & (np.abs(log_df["s_km"] - float(anchor_km)) <= 1.5)].copy()
    if window.empty:
        window = log_df.loc[valid].tail(6).copy()
    else:
        window = window.tail(min(6, len(window))).copy()
    soc_targets = [
        infer_soc_from_loaded_state(
            float(row.battery_voltage_v),
            float(row.battery_current_a),
            float(row.Tamb_archive_C),
            ocv_df,
            base_model,
        )
        for row in window.itertuples(index=False)
    ]
    return {
        "count": int(len(window)),
        "s_km": float(window["s_km"].median()),
        "time_utc": pd.to_datetime(window["time_utc"].iloc[-1], utc=True).strftime("%Y-%m-%dT%H:%M:%SZ"),
        "voltage_v": float(window["battery_voltage_v"].median()),
        "current_a": float(window["battery_current_a"].median()),
        "temp_c": float(window["Tamb_archive_C"].median()),
        "soc_target": float(np.nanmedian(np.asarray(soc_targets, dtype=float))),
    }
\end{lstlisting}

\subsection{L1029 関数 terminal\_metrics}
            \begin{itemize}[leftmargin=1.6em]
              \item 行範囲: L1029--L1075
              \item 所属: module直下
              \item 定義: \path!terminal_metrics(replay_df: pd.DataFrame, ocv_df: pd.DataFrame, base_model, batt: BatteryFitResult, anchor_km: float, *, terminal_anchor: dict | None = None) -> Dict[str, float]!
              \item docstring: 明示されていない。
              \item このブロックが直接呼ぶ主な関数・メソッド: abs, asarray, astype, copy, fillna, float, get, infer\_soc\_from\_loaded\_state, isfinite, isinstance, itertuples, len
              \item 戻り値の要点: \{'retire\_anchor\_s\_km': float(window['s\_km'].median()), 'retire\_anchor\_voltage\_obs\_v': float(window['battery\_voltage\_v\_obs'].median()), 'retire\_anchor\_voltage\_pred\_v': float(window['battery\_voltage\_v\_pred'].median()), 'retire\_anchor\_soc\_obs': soc\_obs\_value, 'retire\_anchor\_soc\_pred': float(window['soc\_pred'].median()), 'retire\_anchor\_soc\_error': float(window['soc\_pred'].median() - soc\_obs\_value)\}
              \item 主なローカル代入名: row, soc\_obs, soc\_obs\_value, terminal\_anchor, valid, window
              \item 読み取る主なインスタンス属性: 特になし。
              \item 更新する主なインスタンス属性: 特になし。
              \item 明示的に送出する例外: 明示的raiseはない。
              \item 制御構造の規模: 条件分岐 2、ループ 0、try 0
            \end{itemize}
            \paragraph{この定義を読むためのPython構文}
            \begin{itemize}[leftmargin=1.6em]
            \item `->`以後は戻り値の型注釈であり、実行時の値を自動変換する命令ではない。
\item 内包表記は、反復・条件・値生成を一つの式で記述してcollectionを作る。
            \end{itemize}
            \paragraph{上から順の処理}
            \begin{enumerate}[leftmargin=1.8em]
            \item valid に np.isfinite(replay\_df['battery\_voltage\_v\_obs']) \& np.isfinite(replay\_df['battery\_current\_a\_obs']) \& np.isfinite(replay\_df['Tamb\_C']) \& np.isfinite(replay\_df['s\_km']) \& (replay\_df['battery\_voltage\_v\_obs'] >= INVALID\_PACK\_VOLTAGE\_MIN\_V) \& \textasciitilde{}replay\_df['exclude\_voltage\_fit'].fillna(False).astype(bool) の結果を代入する。
\item window に replay\_df.loc[valid \& (np.abs(replay\_df['s\_km'] - float(anchor\_km)) <= 1.5)].copy() の結果を代入する。
\item 条件 window.empty を判定し、真なら内部処理を行う。
\item   window に replay\_df.loc[valid].tail(6).copy() の結果を代入する。
\item   上の条件が偽の場合:
\item   window に window.tail(min(6, len(window))).copy() の結果を代入する。
\item terminal\_anchor に terminal\_anchor if isinstance(terminal\_anchor, dict) else \{\} の結果を代入する。
\item 条件 'soc\_target' in terminal\_anchor and terminal\_anchor.get('soc\_target') not in (None, '') を判定し、真なら内部処理を行う。
\item   soc\_obs\_value に float(terminal\_anchor['soc\_target']) の結果を代入する。
\item   上の条件が偽の場合:
\item   soc\_obs に [infer\_soc\_from\_loaded\_state(float(row.battery\_voltage\_v\_obs), float(row.battery\_current\_a\_obs), float(row.Tamb\_C), ocv\_df, base\_model, rint\_scale=float(batt.rint\_scale), r\_line\_ohm=float(batt.r\_line\_ohm)) for row in window.itertuples(index=False)] の結果を代入する。
\item   soc\_obs\_value に float(np.nanmedian(np.asarray(soc\_obs, dtype=float))) の結果を代入する。
\item \{'retire\_anchor\_s\_km': float(window['s\_km'].median()), 'retire\_anchor\_voltage\_obs\_v': float(window['battery\_voltage\_v\_obs'].median()), 'retire\_anchor\_voltage\_pred\_v': float(window['battery\_voltage\_v\_pred'].median()), 'retire\_anchor\_soc\_obs': soc\_obs\_value, 'retire\_anchor\_soc\_pred': float(window['soc\_pred'].median()), 'retire\_anchor\_soc\_error': float(window['soc\_pred'].median() - soc\_obs\_value)\} を返す。
            \end{enumerate}
            \paragraph{代表コード断片}
            \begin{lstlisting}[language=Python]
def terminal_metrics(
    replay_df: pd.DataFrame,
    ocv_df: pd.DataFrame,
    base_model,
    batt: BatteryFitResult,
    anchor_km: float,
    *,
    terminal_anchor: dict | None = None,
) -> Dict[str, float]:
    valid = (
        np.isfinite(replay_df["battery_voltage_v_obs"])
        & np.isfinite(replay_df["battery_current_a_obs"])
        & np.isfinite(replay_df["Tamb_C"])
        & np.isfinite(replay_df["s_km"])
        & (replay_df["battery_voltage_v_obs"] >= INVALID_PACK_VOLTAGE_MIN_V)
        & (~replay_df["exclude_voltage_fit"].fillna(False).astype(bool))
    )
    window = replay_df.loc[valid & (np.abs(replay_df["s_km"] - float(anchor_km)) <= 1.5)].copy()
    if window.empty:
        window = replay_df.loc[valid].tail(6).copy()
    else:
        window = window.tail(min(6, len(window))).copy()
    terminal_anchor = terminal_anchor if isinstance(terminal_anchor, dict) else {}
    if "soc_target" in terminal_anchor and terminal_anchor.get("soc_target") not in (None, ""):
        soc_obs_value = float(terminal_anchor["soc_target"])
    else:
        soc_obs = [
            infer_soc_from_loaded_state(
                float(row.battery_voltage_v_obs),
                float(row.battery_current_a_obs),
                float(row.Tamb_C),
                ocv_df,
                base_model,
                rint_scale=float(batt.rint_scale),
                r_line_ohm=float(batt.r_line_ohm),
...
\end{lstlisting}

\subsection{L1078 関数 replay\_day\_metrics}
            \begin{itemize}[leftmargin=1.6em]
              \item 行範囲: L1078--L1096
              \item 所属: module直下
              \item 定義: \path|replay_day_metrics(replay_df: pd.DataFrame) -> list[dict]|
              \item docstring: 明示されていない。
              \item このブロックが直接呼ぶ主な関数・メソッド: append, astype, fillna, float, groupby, int, len, max, mean, notna, sqrt, sum
              \item 戻り値の要点: rows
              \item 主なローカル代入名: day, group, power\_mask, power\_resid, rows, volt\_mask, volt\_resid
              \item 読み取る主なインスタンス属性: 特になし。
              \item 更新する主なインスタンス属性: 特になし。
              \item 明示的に送出する例外: 明示的raiseはない。
              \item 制御構造の規模: 条件分岐 0、ループ 1、try 0
            \end{itemize}
            \paragraph{この定義を読むためのPython構文}
            \begin{itemize}[leftmargin=1.6em]
            \item `->`以後は戻り値の型注釈であり、実行時の値を自動変換する命令ではない。
            \end{itemize}
            \paragraph{上から順の処理}
            \begin{enumerate}[leftmargin=1.8em]
            \item rows に [] を代入する。
\item replay\_df.groupby('day', dropna=False) を順に走査し、各要素を (day, group) に入れて処理する。
\item   power\_mask に \textasciitilde{}group['exclude\_power\_fit'] の結果を代入する。
\item   volt\_mask に \textasciitilde{}group['exclude\_voltage\_fit'] の結果を代入する。
\item   power\_resid に group.loc[power\_mask, 'battery\_power\_w\_obs'] - group.loc[power\_mask, 'battery\_power\_w\_pred'] の結果を代入する。
\item   volt\_resid に group.loc[volt\_mask, 'battery\_voltage\_v\_obs'] - group.loc[volt\_mask, 'battery\_voltage\_v\_pred'] の結果を代入する。
\item   rows.append(...) を実行する。
\item rows を返す。
            \end{enumerate}
            \paragraph{代表コード断片}
            \begin{lstlisting}[language=Python]
def replay_day_metrics(replay_df: pd.DataFrame) -> list[dict]:
    rows: list[dict] = []
    for day, group in replay_df.groupby("day", dropna=False):
        power_mask = ~group["exclude_power_fit"]
        volt_mask = ~group["exclude_voltage_fit"]
        power_resid = group.loc[power_mask, "battery_power_w_obs"] - group.loc[power_mask, "battery_power_w_pred"]
        volt_resid = group.loc[volt_mask, "battery_voltage_v_obs"] - group.loc[volt_mask, "battery_voltage_v_pred"]
        rows.append(
            {
                "day": int(day) if pd.notna(day) else -1,
                "distance_end_km": float(pd.to_numeric(group["s_km"], errors="coerce").max()),
                "final_soc_pred": float(pd.to_numeric(group["soc_pred"], errors="coerce").iloc[-1]),
                "power_rmse_clean_w": float(np.sqrt(np.mean(power_resid.to_numpy(dtype=float) ** 2))) if len(power_resid) else float("nan"),
                "voltage_rmse_clean_v": float(np.sqrt(np.mean(volt_resid.to_numpy(dtype=float) ** 2))) if len(volt_resid) else float("nan"),
                "excluded_power_points": int(pd.to_numeric(group["exclude_power_fit"], errors="coerce").fillna(False).astype(bool).sum()),
                "excluded_voltage_points": int(pd.to_numeric(group["exclude_voltage_fit"], errors="coerce").fillna(False).astype(bool).sum()),
            }
        )
    return rows
\end{lstlisting}

\subsection{L1099 関数 identification\_selection\_score}
            \begin{itemize}[leftmargin=1.6em]
              \item 行範囲: L1099--L1132
              \item 所属: module直下
              \item 定義: \path|identification_selection_score(metrics: Dict[str, float]) -> float|
              \item docstring: 明示されていない。
              \item このブロックが直接呼ぶ主な関数・メソッド: abs, float, get
              \item 戻り値の要点: power\_term + energy\_term + independent\_pv\_term + vehicle\_voltage\_term + battery\_voltage\_term + battery\_terminal\_term + vehicle\_terminal\_term + end\_to\_end\_terminal\_term + fit\_window\_term
              \item 主なローカル代入名: battery\_terminal\_term, battery\_voltage\_term, end\_to\_end\_terminal\_term, energy\_rmse\_25km, energy\_term, fit\_window\_term, independent\_pv\_term, power\_rmse, power\_term, robust\_power\_rmse, vehicle\_terminal\_term, vehicle\_voltage\_term
              \item 読み取る主なインスタンス属性: 特になし。
              \item 更新する主なインスタンス属性: 特になし。
              \item 明示的に送出する例外: 明示的raiseはない。
              \item 制御構造の規模: 条件分岐 0、ループ 0、try 0
            \end{itemize}
            \paragraph{この定義を読むためのPython構文}
            \begin{itemize}[leftmargin=1.6em]
            \item `->`以後は戻り値の型注釈であり、実行時の値を自動変換する命令ではない。
            \end{itemize}
            \paragraph{上から順の処理}
            \begin{enumerate}[leftmargin=1.8em]
            \item power\_rmse に float(metrics.get('power\_rmse\_clean\_w', float('inf'))) の結果を代入する。
\item robust\_power\_rmse に float(metrics.get('power\_residual\_mean\_120s\_rmse\_w', power\_rmse)) の結果を代入する。
\item energy\_rmse\_25km に float(metrics.get('energy\_error\_25km\_rmse\_wh', float('inf'))) の結果を代入する。
\item power\_term に 0.6 * power\_rmse + 0.4 * robust\_power\_rmse の結果を代入する。
\item energy\_term に 0.5 * energy\_rmse\_25km の結果を代入する。
\item independent\_pv\_term に 0.25 * float(metrics.get('end\_to\_end\_moving\_pv\_rmse\_w', float('inf'))) の結果を代入する。
\item vehicle\_voltage\_term に 8.0 * float(metrics.get('voltage\_rmse\_clean\_v', float('inf'))) の結果を代入する。
\item battery\_voltage\_term に 4.0 * float(metrics.get('battery\_conditional\_voltage\_rmse\_clean\_v', float('inf'))) の結果を代入する。
\item battery\_terminal\_term に 80.0 * abs(float(metrics.get('battery\_conditional\_retire\_anchor\_soc\_error', float('inf')))) の結果を代入する。
\item vehicle\_terminal\_term に 1000.0 * abs(float(metrics.get('retire\_anchor\_soc\_error', float('inf')))) の結果を代入する。
\item end\_to\_end\_terminal\_term に 300.0 * abs(float(metrics.get('end\_to\_end\_retire\_anchor\_soc\_error', float('inf')))) の結果を代入する。
\item fit\_window\_term に 0.35 * float(metrics.get('power\_rmse\_fit\_window\_w', power\_term)) の結果を代入する。
\item power\_term + energy\_term + independent\_pv\_term + vehicle\_voltage\_term + battery\_voltage\_term + battery\_terminal\_term + vehicle\_terminal\_term + end\_to\_end\_terminal\_term + fit\_window\_term を返す。
            \end{enumerate}
            \paragraph{代表コード断片}
            \begin{lstlisting}[language=Python]
def identification_selection_score(metrics: Dict[str, float]) -> float:
    power_rmse = float(metrics.get("power_rmse_clean_w", float("inf")))
    robust_power_rmse = float(metrics.get("power_residual_mean_120s_rmse_w", power_rmse))
    energy_rmse_25km = float(metrics.get("energy_error_25km_rmse_wh", float("inf")))
    power_term = 0.60 * power_rmse + 0.40 * robust_power_rmse
    energy_term = 0.50 * energy_rmse_25km
    independent_pv_term = 0.25 * float(
        metrics.get("end_to_end_moving_pv_rmse_w", float("inf"))
    )
    vehicle_voltage_term = 8.0 * float(metrics.get("voltage_rmse_clean_v", float("inf")))
    battery_voltage_term = 4.0 * float(
        metrics.get("battery_conditional_voltage_rmse_clean_v", float("inf"))
    )
    battery_terminal_term = 80.0 * abs(
        float(metrics.get("battery_conditional_retire_anchor_soc_error", float("inf")))
    )
    vehicle_terminal_term = 1000.0 * abs(
        float(metrics.get("retire_anchor_soc_error", float("inf")))
    )
    end_to_end_terminal_term = 300.0 * abs(
        float(metrics.get("end_to_end_retire_anchor_soc_error", float("inf")))
    )
    fit_window_term = 0.35 * float(metrics.get("power_rmse_fit_window_w", power_term))
    return (
        power_term
        + energy_term
        + independent_pv_term
        + vehicle_voltage_term
        + battery_voltage_term
        + battery_terminal_term
        + vehicle_terminal_term
        + end_to_end_terminal_term
        + fit_window_term
    )
\end{lstlisting}

\subsection{L1135 関数 condition\_vehicle\_fit\_on\_measured\_pv}
            \begin{itemize}[leftmargin=1.6em]
              \item 行範囲: L1135--L1152
              \item 所属: module直下
              \item 定義: \path|condition_vehicle_fit_on_measured_pv(frame: pd.DataFrame) -> pd.DataFrame|
              \item docstring: Use measured array power when identifying vehicle-side coefficients.

Forecast/PV-map error is validated separately.  Feeding predicted PV into the
vehicle fit otherwise lets a cloudy-day irradiance error masquerade as CdA,
rolling resistance, or drivetrain-efficiency error.
              \item このブロックが直接呼ぶ主な関数・メソッド: clip, copy, get, notna, to\_numeric, where
              \item 戻り値の要点: out / out
              \item 主なローカル代入名: measured, out, predicted, usable
              \item 読み取る主なインスタンス属性: 特になし。
              \item 更新する主なインスタンス属性: 特になし。
              \item 明示的に送出する例外: 明示的raiseはない。
              \item 制御構造の規模: 条件分岐 1、ループ 0、try 0
            \end{itemize}
            \paragraph{この定義を読むためのPython構文}
            \begin{itemize}[leftmargin=1.6em]
            \item `->`以後は戻り値の型注釈であり、実行時の値を自動変換する命令ではない。
            \end{itemize}
            \paragraph{上から順の処理}
            \begin{enumerate}[leftmargin=1.8em]
            \item out に frame.copy() の結果を代入する。
\item predicted に pd.to\_numeric(out.get('solar\_power\_w\_model'), errors='coerce') の結果を代入する。
\item measured に pd.to\_numeric(out.get('solar\_power\_w\_obs'), errors='coerce') の結果を代入する。
\item 条件 predicted is None or measured is None を判定し、真なら内部処理を行う。
\item   out を返す。
\item measured に measured.clip(lower=0.0) の結果を代入する。
\item usable に measured.notna() の結果を代入する。
\item out['solar\_power\_w\_forecast\_model'] に predicted の結果を代入する。
\item out['solar\_power\_w\_model'] に predicted.where(\textasciitilde{}usable, measured) の結果を代入する。
\item out['vehicle\_fit\_solar\_source'] に np.where(usable, 'measured', 'forecast\_fallback') の結果を代入する。
\item out を返す。
            \end{enumerate}
            \paragraph{代表コード断片}
            \begin{lstlisting}[language=Python]
def condition_vehicle_fit_on_measured_pv(frame: pd.DataFrame) -> pd.DataFrame:
    """Use measured array power when identifying vehicle-side coefficients.

    Forecast/PV-map error is validated separately.  Feeding predicted PV into the
    vehicle fit otherwise lets a cloudy-day irradiance error masquerade as CdA,
    rolling resistance, or drivetrain-efficiency error.
    """
    out = frame.copy()
    predicted = pd.to_numeric(out.get("solar_power_w_model"), errors="coerce")
    measured = pd.to_numeric(out.get("solar_power_w_obs"), errors="coerce")
    if predicted is None or measured is None:
        return out
    measured = measured.clip(lower=0.0)
    usable = measured.notna()
    out["solar_power_w_forecast_model"] = predicted
    out["solar_power_w_model"] = predicted.where(~usable, measured)
    out["vehicle_fit_solar_source"] = np.where(usable, "measured", "forecast_fallback")
    return out
\end{lstlisting}

\subsection{L1155 関数 calibrate\_solar\_measurement\_to\_pack}
            \begin{itemize}[leftmargin=1.6em]
              \item 行範囲: L1155--L1268
              \item 所属: module直下
              \item 定義: \path|calibrate_solar_measurement_to_pack(frame: pd.DataFrame, *, known_aux_power_w: float, stopped_speed_kmh: float = 1.0, minimum_solar_power_w: float = 50.0, minimum_samples: int = 500, gain_bounds: tuple[float, float] = (0.7, 1.05)) -> tuple[pd.DataFrame, Dict[str, Any]]|
              \item docstring: Calibrate the ZP solar channel from stationary DC-bus power balance.

At zero wheel speed, the independently measured channels satisfy
``P\_batt = P\_aux - gain * P\_solar\_raw``.  The known 21 W auxiliary load
anchors the intercept, so the fitted gain cannot absorb vehicle drag or
forecast irradiance error.
              \item このブロックが直接呼ぶ主な関数・メソッド: Series, abs, array, asarray, astype, bool, copy, fillna, float, get, groupby, int
              \item 戻り値の要点: (out, result)
              \item 主なローカル代入名: accepted, battery, candidate\_gain, daily\_gains, daily\_std, daily\_values, day\_fit, day\_value, diagnostic, excluded, fixed, free\_intercept, gain, group, gx, gy, mask, out, raw, raw\_column
              \item 読み取る主なインスタンス属性: 特になし。
              \item 更新する主なインスタンス属性: 特になし。
              \item 明示的に送出する例外: 明示的raiseはない。
              \item 制御構造の規模: 条件分岐 5、ループ 1、try 0
            \end{itemize}
            \paragraph{この定義を読むためのPython構文}
            \begin{itemize}[leftmargin=1.6em]
            \item `->`以後は戻り値の型注釈であり、実行時の値を自動変換する命令ではない。
\item lambdaは名前を付けずに短い関数オブジェクトを作る構文である。
            \end{itemize}
            \paragraph{上から順の処理}
            \begin{enumerate}[leftmargin=1.8em]
            \item out に frame.copy() の結果を代入する。
\item raw\_column に 'solar\_power\_w\_obs\_raw' if 'solar\_power\_w\_obs\_raw' in out.columns else 'solar\_power\_w\_obs' の結果を代入する。
\item raw に pd.to\_numeric(out.get(raw\_column), errors='coerce') の結果を代入する。
\item battery に pd.to\_numeric(out.get('battery\_power\_w\_obs'), errors='coerce') の結果を代入する。
\item speed に pd.to\_numeric(out.get('speed\_kmh'), errors='coerce') の結果を代入する。
\item excluded に pd.Series(False, index=out.index) の結果を代入する。
\item 条件 'exclude\_power\_fit' in out.columns を判定し、真なら内部処理を行う。
\item   excluded に out['exclude\_power\_fit'].fillna(False).astype(bool) の結果を代入する。
\item mask に \textasciitilde{}excluded \& raw.notna() \& battery.notna() \& speed.notna() \& (speed <= float(stopped\_speed\_kmh)) \& (raw >= float(minimum\_solar\_power\_w)) の結果を代入する。
\item x に raw.loc[mask].to\_numpy(dtype=float) の結果を代入する。
\item y に battery.loc[mask].to\_numpy(dtype=float) の結果を代入する。
\item result に \{'method': 'stationary DC-bus balance P\_batt=P\_aux-gain*P\_solar\_raw with robust Huber loss', 'known\_aux\_power\_w': float(known\_aux\_power\_w), 'sample\_count': int(len(x)), 'minimum\_samples': int(minimum\_samples), 'gain\_bounds': [float(gain\_bounds[0]), float(gain\_bounds[1])], 'accepted': False, 'gain\_to\_pack': 1.0, 'raw\_column': raw\_column, 'corrected\_column': 'solar\_power\_w\_obs'\} を代入する。
\item gain に 1.0 の結果を代入する。
\item 条件 len(x) >= int(minimum\_samples) を判定し、真なら内部処理を行う。
\item   fixed に least\_squares(lambda theta: y - (float(known\_aux\_power\_w) - float(theta[0]) * x), x0=np.array([0.93], dtype=float), bounds=(np.array([gain\_bounds[0]]), np.array([gain\_bounds[1]])), loss='huber', f\_scale=20.0) の結果を代入する。
\item   diagnostic に least\_squares(lambda theta: y - (float(theta[0]) - float(theta[1]) * x), x0=np.array([float(known\_aux\_power\_w), float(fixed.x[0])], dtype=float), bounds=(np.array([0.0, gain\_bounds[0]], dtype=float), np.array([100.0, gain\_bounds[1]], dtype=float)), loss='huber', f\_scale=20.0) の結果を代入する。
\item   candidate\_gain に float(fixed.x[0]) の結果を代入する。
\item   residual に y - (float(known\_aux\_power\_w) - candidate\_gain * x) の結果を代入する。
\item   daily\_gains に \{\} を代入する。
\item   条件 'day' in out.columns を判定し、真なら内部処理を行う。
\item     out.loc[mask].groupby('day', dropna=False) を順に走査し、各要素を (day\_value, group) に入れて処理する。
\item       gx に pd.to\_numeric(group[raw\_column], errors='coerce').to\_numpy(dtype=float) の結果を代入する。
\item       gy に pd.to\_numeric(group['battery\_power\_w\_obs'], errors='coerce').to\_numpy(dtype=float) の結果を代入する。
\item       条件 len(gx) < 50 を判定し、真なら内部処理を行う。
\item       day\_fit に least\_squares(lambda theta: gy - (float(known\_aux\_power\_w) - float(theta[0]) * gx), x0=np.array([candidate\_gain], dtype=float), bounds=(np.array([gain\_bounds[0]]), np.array([gain\_bounds[1]])), loss='huber', f\_scale=20.0) の結果を代入する。
\item       daily\_gains[str(day\_value)] に float(day\_fit.x[0]) の結果を代入する。
\item   daily\_values に np.asarray(list(daily\_gains.values()), dtype=float) の結果を代入する。
\item   daily\_std に float(np.std(daily\_values, ddof=1)) if len(daily\_values) >= 2 else float('nan') の結果を代入する。
\item   free\_intercept に float(diagnostic.x[0]) の結果を代入する。
\item   accepted に bool(np.isfinite(candidate\_gain) and float(gain\_bounds[0]) < candidate\_gain < float(gain\_bounds[1]) and (abs(free\_intercept - float(known\_aux\_power\_w)) <= 5.0) and (not np.isfinite(daily\_std) or daily\_std <= 0.02)) の結果を代入する。
\item   result.update(...) を実行する。
\item   条件 accepted を判定し、真なら内部処理を行う。
\item     gain に candidate\_gain の結果を代入する。
\item out['solar\_power\_w\_obs\_raw'] に raw の結果を代入する。
\item out['solar\_measurement\_gain\_to\_pack'] に float(gain) の結果を代入する。
\item out['solar\_power\_w\_obs'] に raw * float(gain) の結果を代入する。
\item (out, result) を返す。
            \end{enumerate}
            \paragraph{代表コード断片}
            \begin{lstlisting}[language=Python]
def calibrate_solar_measurement_to_pack(
    frame: pd.DataFrame,
    *,
    known_aux_power_w: float,
    stopped_speed_kmh: float = 1.0,
    minimum_solar_power_w: float = 50.0,
    minimum_samples: int = 500,
    gain_bounds: tuple[float, float] = (0.70, 1.05),
) -> tuple[pd.DataFrame, Dict[str, Any]]:
    """Calibrate the ZP solar channel from stationary DC-bus power balance.

    At zero wheel speed, the independently measured channels satisfy
    ``P_batt = P_aux - gain * P_solar_raw``.  The known 21 W auxiliary load
    anchors the intercept, so the fitted gain cannot absorb vehicle drag or
    forecast irradiance error.
    """
    out = frame.copy()
    raw_column = (
        "solar_power_w_obs_raw"
        if "solar_power_w_obs_raw" in out.columns
        else "solar_power_w_obs"
    )
    raw = pd.to_numeric(out.get(raw_column), errors="coerce")
    battery = pd.to_numeric(out.get("battery_power_w_obs"), errors="coerce")
    speed = pd.to_numeric(out.get("speed_kmh"), errors="coerce")
    excluded = pd.Series(False, index=out.index)
    if "exclude_power_fit" in out.columns:
        excluded = out["exclude_power_fit"].fillna(False).astype(bool)
    mask = (
        (~excluded)
        & raw.notna()
        & battery.notna()
        & speed.notna()
        & (speed <= float(stopped_speed_kmh))
        & (raw >= float(minimum_solar_power_w))
...
\end{lstlisting}

\subsection{L1271 関数 \_shift\_acceleration\_within\_segments}
            \begin{itemize}[leftmargin=1.6em]
              \item 行範囲: L1271--L1280
              \item 所属: module直下
              \item 定義: \path!_shift_acceleration_within_segments(frame: pd.DataFrame, sample_shift: int, acceleration: pd.Series | None = None) -> pd.Series!
              \item docstring: Shift a derived acceleration trace without leaking across race days or log gaps.
              \item このブロックが直接呼ぶ主な関数・メソッド: Series, cumsum, groupby, int, replay\_segment\_start\_mask, shift, to\_numeric
              \item 戻り値の要点: acceleration.groupby(segment\_id).shift(int(sample\_shift))
              \item 主なローカル代入名: acceleration, segment\_id
              \item 読み取る主なインスタンス属性: 特になし。
              \item 更新する主なインスタンス属性: 特になし。
              \item 明示的に送出する例外: 明示的raiseはない。
              \item 制御構造の規模: 条件分岐 1、ループ 0、try 0
            \end{itemize}
            \paragraph{この定義を読むためのPython構文}
            \begin{itemize}[leftmargin=1.6em]
            \item `->`以後は戻り値の型注釈であり、実行時の値を自動変換する命令ではない。
            \end{itemize}
            \paragraph{上から順の処理}
            \begin{enumerate}[leftmargin=1.8em]
            \item 条件 acceleration is None を判定し、真なら内部処理を行う。
\item   acceleration に pd.to\_numeric(frame['accel\_ms2'], errors='coerce') の結果を代入する。
\item segment\_id に pd.Series(replay\_segment\_start\_mask(frame), index=frame.index).cumsum() の結果を代入する。
\item acceleration.groupby(segment\_id).shift(int(sample\_shift)) を返す。
            \end{enumerate}
            \paragraph{代表コード断片}
            \begin{lstlisting}[language=Python]
def _shift_acceleration_within_segments(
    frame: pd.DataFrame,
    sample_shift: int,
    acceleration: pd.Series | None = None,
) -> pd.Series:
    """Shift a derived acceleration trace without leaking across race days or log gaps."""
    if acceleration is None:
        acceleration = pd.to_numeric(frame["accel_ms2"], errors="coerce")
    segment_id = pd.Series(replay_segment_start_mask(frame), index=frame.index).cumsum()
    return acceleration.groupby(segment_id).shift(int(sample_shift))
\end{lstlisting}

\subsection{L1283 関数 \_acceleration\_trace\_from\_speed}
            \begin{itemize}[leftmargin=1.6em]
              \item 行範囲: L1283--L1318
              \item 所属: module直下
              \item 定義: \path|_acceleration_trace_from_speed(frame: pd.DataFrame, *, method: str, window_samples: int) -> pd.Series|
              \item docstring: 明示されていない。
              \item このブロックが直接呼ぶ主な関数・メソッド: Series, ValueError, clip, cumsum, diff, fillna, get, groupby, int, len, lower, max
              \item 戻り値の要点: smoothed.fillna(0.0).clip(lower=-1.5, upper=1.5)
              \item 主なローカル代入名: \_, dt\_sec, group, raw, segment\_id, segment\_start, smoothed, speed\_ms, usable\_window, window
              \item 読み取る主なインスタンス属性: 特になし。
              \item 更新する主なインスタンス属性: 特になし。
              \item 明示的に送出する例外: ValueError(f'unsupported acceleration filter method: \{method\}')
              \item 制御構造の規模: 条件分岐 4、ループ 1、try 0
            \end{itemize}
            \paragraph{この定義を読むためのPython構文}
            \begin{itemize}[leftmargin=1.6em]
            \item `->`以後は戻り値の型注釈であり、実行時の値を自動変換する命令ではない。
            \end{itemize}
            \paragraph{上から順の処理}
            \begin{enumerate}[leftmargin=1.8em]
            \item speed\_ms に pd.to\_numeric(frame['speed\_kmh'], errors='coerce').fillna(0.0) / 3.6 の結果を代入する。
\item dt\_sec に pd.to\_numeric(frame.get('dt\_sec'), errors='coerce').replace(0.0, np.nan) の結果を代入する。
\item raw に speed\_ms.diff() / dt\_sec の結果を代入する。
\item segment\_start に pd.Series(replay\_segment\_start\_mask(frame), index=frame.index) の結果を代入する。
\item raw.loc[segment\_start] に 0.0 の結果を代入する。
\item raw に raw.replace([np.inf, -np.inf], np.nan).fillna(0.0) の結果を代入する。
\item segment\_id に segment\_start.cumsum() の結果を代入する。
\item window に max(1, int(window\_samples)) の結果を代入する。
\item 条件 window \% 2 == 0 を判定し、真なら内部処理を行う。
\item   window を Add で更新する。
\item smoothed に pd.Series(index=frame.index, dtype=float) の結果を代入する。
\item raw.groupby(segment\_id) を順に走査し、各要素を (\_, group) に入れて処理する。
\item   条件 str(method).lower() == 'savgol' を判定し、真なら内部処理を行う。
\item     usable\_window に min(window, len(group) if len(group) \% 2 else len(group) - 1) の結果を代入する。
\item     条件 usable\_window >= 5 を判定し、真なら内部処理を行う。
\item       smoothed.loc[group.index] に savgol\_filter(group.to\_numpy(dtype=float), usable\_window, min(2, usable\_window - 1), mode='interp') の結果を代入する。
\item       上の条件が偽の場合:
\item       smoothed.loc[group.index] に group の結果を代入する。
\item     上の条件が偽の場合:
\item     条件 str(method).lower() == 'median' を判定し、真なら内部処理を行う。
\item       smoothed.loc[group.index] に group.rolling(window, center=True, min\_periods=1).median() の結果を代入する。
\item       上の条件が偽の場合:
\item       ValueError(f'unsupported acceleration filter method: \{method\}') を送出する。
\item smoothed.fillna(0.0).clip(lower=-1.5, upper=1.5) を返す。
            \end{enumerate}
            \paragraph{代表コード断片}
            \begin{lstlisting}[language=Python]
def _acceleration_trace_from_speed(
    frame: pd.DataFrame,
    *,
    method: str,
    window_samples: int,
) -> pd.Series:
    speed_ms = pd.to_numeric(frame["speed_kmh"], errors="coerce").fillna(0.0) / 3.6
    dt_sec = pd.to_numeric(frame.get("dt_sec"), errors="coerce").replace(0.0, np.nan)
    raw = speed_ms.diff() / dt_sec
    segment_start = pd.Series(replay_segment_start_mask(frame), index=frame.index)
    raw.loc[segment_start] = 0.0
    raw = raw.replace([np.inf, -np.inf], np.nan).fillna(0.0)
    segment_id = segment_start.cumsum()
    window = max(1, int(window_samples))
    if window % 2 == 0:
        window += 1
    smoothed = pd.Series(index=frame.index, dtype=float)
    for _, group in raw.groupby(segment_id):
        if str(method).lower() == "savgol":
            usable_window = min(window, len(group) if len(group) % 2 else len(group) - 1)
            if usable_window >= 5:
                smoothed.loc[group.index] = savgol_filter(
                    group.to_numpy(dtype=float),
                    usable_window,
                    min(2, usable_window - 1),
                    mode="interp",
                )
            else:
                smoothed.loc[group.index] = group
        elif str(method).lower() == "median":
            smoothed.loc[group.index] = group.rolling(
                window, center=True, min_periods=1
            ).median()
        else:
            raise ValueError(f"unsupported acceleration filter method: {method}")
...
\end{lstlisting}

\subsection{L1321 関数 \_bilinear\_interp\_array}
            \begin{itemize}[leftmargin=1.6em]
              \item 行範囲: L1321--L1338
              \item 所属: module直下
              \item 定義: \path|_bilinear_interp_array(x_grid, y_grid, values, x, y) -> np.ndarray|
              \item docstring: 明示されていない。
              \item このブロックが直接呼ぶ主な関数・メソッド: asarray, clip, divide, len, searchsorted, zeros\_like
              \item 戻り値の要点: (1.0 - wx) * (1.0 - wy) * values[ix, iy] + wx * (1.0 - wy) * values[ix + 1, iy] + (1.0 - wx) * wy * values[ix, iy + 1] + wx * wy * values[ix + 1, iy + 1]
              \item 主なローカル代入名: ix, iy, values, wx, wy, x, x0, x1, x\_grid, y, y0, y1, y\_grid
              \item 読み取る主なインスタンス属性: 特になし。
              \item 更新する主なインスタンス属性: 特になし。
              \item 明示的に送出する例外: 明示的raiseはない。
              \item 制御構造の規模: 条件分岐 0、ループ 0、try 0
            \end{itemize}
            \paragraph{この定義を読むためのPython構文}
            \begin{itemize}[leftmargin=1.6em]
            \item `->`以後は戻り値の型注釈であり、実行時の値を自動変換する命令ではない。
            \end{itemize}
            \paragraph{上から順の処理}
            \begin{enumerate}[leftmargin=1.8em]
            \item x\_grid に np.asarray(x\_grid, dtype=float) の結果を代入する。
\item y\_grid に np.asarray(y\_grid, dtype=float) の結果を代入する。
\item values に np.asarray(values, dtype=float) の結果を代入する。
\item x に np.clip(np.asarray(x, dtype=float), x\_grid[0], x\_grid[-1]) の結果を代入する。
\item y に np.clip(np.asarray(y, dtype=float), y\_grid[0], y\_grid[-1]) の結果を代入する。
\item ix に np.clip(np.searchsorted(x\_grid, x) - 1, 0, len(x\_grid) - 2) の結果を代入する。
\item iy に np.clip(np.searchsorted(y\_grid, y) - 1, 0, len(y\_grid) - 2) の結果を代入する。
\item (x0, x1) に (x\_grid[ix], x\_grid[ix + 1]) の結果を代入する。
\item (y0, y1) に (y\_grid[iy], y\_grid[iy + 1]) の結果を代入する。
\item wx に np.divide(x - x0, x1 - x0, out=np.zeros\_like(x), where=x1 != x0) の結果を代入する。
\item wy に np.divide(y - y0, y1 - y0, out=np.zeros\_like(y), where=y1 != y0) の結果を代入する。
\item (1.0 - wx) * (1.0 - wy) * values[ix, iy] + wx * (1.0 - wy) * values[ix + 1, iy] + (1.0 - wx) * wy * values[ix, iy + 1] + wx * wy * values[ix + 1, iy + 1] を返す。
            \end{enumerate}
            \paragraph{代表コード断片}
            \begin{lstlisting}[language=Python]
def _bilinear_interp_array(x_grid, y_grid, values, x, y) -> np.ndarray:
    x_grid = np.asarray(x_grid, dtype=float)
    y_grid = np.asarray(y_grid, dtype=float)
    values = np.asarray(values, dtype=float)
    x = np.clip(np.asarray(x, dtype=float), x_grid[0], x_grid[-1])
    y = np.clip(np.asarray(y, dtype=float), y_grid[0], y_grid[-1])
    ix = np.clip(np.searchsorted(x_grid, x) - 1, 0, len(x_grid) - 2)
    iy = np.clip(np.searchsorted(y_grid, y) - 1, 0, len(y_grid) - 2)
    x0, x1 = x_grid[ix], x_grid[ix + 1]
    y0, y1 = y_grid[iy], y_grid[iy + 1]
    wx = np.divide(x - x0, x1 - x0, out=np.zeros_like(x), where=x1 != x0)
    wy = np.divide(y - y0, y1 - y0, out=np.zeros_like(y), where=y1 != y0)
    return (
        (1.0 - wx) * (1.0 - wy) * values[ix, iy]
        + wx * (1.0 - wy) * values[ix + 1, iy]
        + (1.0 - wx) * wy * values[ix, iy + 1]
        + wx * wy * values[ix + 1, iy + 1]
    )
\end{lstlisting}

\subsection{L1341 関数 \_map\_efficiency\_array}
            \begin{itemize}[leftmargin=1.6em]
              \item 行範囲: L1341--L1361
              \item 所属: module直下
              \item 定義: \path|_map_efficiency_array(base_model, maps, speed_ms, torque_nm, *, regen: bool) -> np.ndarray|
              \item docstring: 明示されていない。
              \item このブロックが直接呼ぶ主な関数・メソッド: \_bilinear\_interp\_array, clip, empty, float, full, get, len, lower, str
              \item 戻り値の要点: np.clip(result, 0.4 if regen else 0.55, 0.95 if regen else 0.99)
              \item 主なローカル代入名: eco\_max, grid, key, mode, result, selected, selected\_power, use\_power
              \item 読み取る主なインスタンス属性: 特になし。
              \item 更新する主なインスタンス属性: 特になし。
              \item 明示的に送出する例外: 明示的raiseはない。
              \item 制御構造の規模: 条件分岐 1、ループ 1、try 0
            \end{itemize}
            \paragraph{この定義を読むためのPython構文}
            \begin{itemize}[leftmargin=1.6em]
            \item `->`以後は戻り値の型注釈であり、実行時の値を自動変換する命令ではない。
            \end{itemize}
            \paragraph{上から順の処理}
            \begin{enumerate}[leftmargin=1.8em]
            \item mode に str(base\_model.drive\_mode or 'default').lower() の結果を代入する。
\item 条件 mode in \{'eco', 'power'\} を判定し、真なら内部処理を行う。
\item   selected\_power に np.full(len(speed\_ms), mode == 'power', dtype=bool) の結果を代入する。
\item   上の条件が偽の場合:
\item   eco\_max に base\_model.tau\_max.get('eco', base\_model.tau\_max.get('default', 0.0)) の結果を代入する。
\item   selected\_power に torque\_nm > float(eco\_max) + float(base\_model.drive\_mode\_tau\_margin) の結果を代入する。
\item result に np.empty(len(speed\_ms), dtype=float) の結果を代入する。
\item (False, True) を順に走査し、各要素を use\_power に入れて処理する。
\item   selected に selected\_power == use\_power の結果を代入する。
\item   key に 'power' if use\_power else 'eco' の結果を代入する。
\item   grid に maps.get(key, maps['default']) の結果を代入する。
\item   result[selected] に \_bilinear\_interp\_array(*grid, speed\_ms[selected], torque\_nm[selected]) の結果を代入する。
\item np.clip(result, 0.4 if regen else 0.55, 0.95 if regen else 0.99) を返す。
            \end{enumerate}
            \paragraph{代表コード断片}
            \begin{lstlisting}[language=Python]
def _map_efficiency_array(base_model, maps, speed_ms, torque_nm, *, regen: bool) -> np.ndarray:
    mode = str(base_model.drive_mode or "default").lower()
    if mode in {"eco", "power"}:
        selected_power = np.full(len(speed_ms), mode == "power", dtype=bool)
    else:
        eco_max = base_model.tau_max.get("eco", base_model.tau_max.get("default", 0.0))
        selected_power = torque_nm > (
            float(eco_max) + float(base_model.drive_mode_tau_margin)
        )
    result = np.empty(len(speed_ms), dtype=float)
    for use_power in (False, True):
        selected = selected_power == use_power
        key = "power" if use_power else "eco"
        grid = maps.get(key, maps["default"])
        result[selected] = _bilinear_interp_array(
            *grid, speed_ms[selected], torque_nm[selected]
        )
    # The candidate scale is applied by the caller. Inheriting the scale from
    # the input profile here would apply old_scale * candidate_scale during
    # fitting, while the generated profile contains candidate_scale only.
    return np.clip(result, 0.40 if regen else 0.55, 0.95 if regen else 0.99)
\end{lstlisting}

\subsection{L1364 関数 \_motion\_predictions\_for\_acceleration}
            \begin{itemize}[leftmargin=1.6em]
              \item 行範囲: L1364--L1428
              \item 所属: module直下
              \item 定義: \path|_motion_predictions_for_acceleration(frame: pd.DataFrame, acceleration: pd.Series, base_model, mot: MotionFitResult) -> np.ndarray|
              \item docstring: Vectorized equivalent of motion\_power\_prediction for filter/lag search.
              \item このブロックが直接呼ぶ主な関数・メソッド: Series, \_map\_efficiency\_array, abs, arctan, clip, cos, empty, fillna, float, full, len, lower
              \item 戻り値の要点: electrical + float(mot.p\_aux\_w) - solar
              \item 主なローカル代入名: air\_density, altitude, ambient, ambient\_source, drive\_eff, electrical, elevation, elevation\_source, force, headwind, normal\_force, omega\_wheel, p\_mech, positive, pressure, pressure\_ratio, regen\_eff, relative\_speed, slope, solar
              \item 読み取る主なインスタンス属性: 特になし。
              \item 更新する主なインスタンス属性: 特になし。
              \item 明示的に送出する例外: 明示的raiseはない。
              \item 制御構造の規模: 条件分岐 1、ループ 0、try 0
            \end{itemize}
            \paragraph{この定義を読むためのPython構文}
            \begin{itemize}[leftmargin=1.6em]
            \item `->`以後は戻り値の型注釈であり、実行時の値を自動変換する命令ではない。
            \end{itemize}
            \paragraph{上から順の処理}
            \begin{enumerate}[leftmargin=1.8em]
            \item speed\_ms に pd.to\_numeric(frame['speed\_kmh'], errors='coerce').fillna(0.0).to\_numpy() / 3.6 の結果を代入する。
\item slope に pd.to\_numeric(frame['slope\_pct'], errors='coerce').fillna(0.0).to\_numpy() の結果を代入する。
\item headwind に pd.to\_numeric(frame['headwind\_archive\_ms'], errors='coerce').fillna(0.0).to\_numpy() * float(mot.headwind\_gain) の結果を代入する。
\item solar に pd.to\_numeric(frame['solar\_power\_w\_obs'], errors='coerce').fillna(0.0).to\_numpy() の結果を代入する。
\item theta に np.arctan(slope * float(mot.grade\_scale) / 100.0) の結果を代入する。
\item relative\_speed に np.maximum(0.0, speed\_ms + headwind) の結果を代入する。
\item ambient\_source に frame['Tamb\_archive\_C'] if 'Tamb\_archive\_C' in frame else pd.Series(25.0, index=frame.index) の結果を代入する。
\item elevation\_source に frame['alt\_m'] if 'alt\_m' in frame else pd.Series(0.0, index=frame.index) の結果を代入する。
\item ambient に pd.to\_numeric(ambient\_source, errors='coerce').fillna(25.0).to\_numpy() の結果を代入する。
\item elevation に pd.to\_numeric(elevation\_source, errors='coerce').fillna(0.0).to\_numpy() の結果を代入する。
\item 条件 str(base\_model.p.air\_density\_mode or 'constant').lower() == 'ideal\_gas\_altitude' を判定し、真なら内部処理を行う。
\item   altitude に np.clip(elevation, -500.0, 11000.0) の結果を代入する。
\item   pressure\_ratio に np.maximum(0.05, 1.0 - 0.0065 * altitude / 288.15) ** 5.255877 の結果を代入する。
\item   pressure に float(base\_model.p.air\_density\_reference\_pressure\_pa) * pressure\_ratio の結果を代入する。
\item   air\_density に pressure / (287.05 * np.maximum(180.0, ambient + 273.15)) の結果を代入する。
\item   上の条件が偽の場合:
\item   air\_density に np.full(len(frame), float(base\_model.p.rho), dtype=float) の結果を代入する。
\item normal\_force に base\_model.p.m * base\_model.p.g * np.cos(theta) の結果を代入する。
\item force に 0.5 * air\_density * float(mot.cda) * relative\_speed ** 2 + float(mot.crr) * normal\_force + base\_model.p.m * base\_model.p.g * np.sin(theta) の結果を代入する。
\item p\_mech に force * speed\_ms + base\_model.p.m * acceleration.fillna(0.0).to\_numpy(dtype=float) * speed\_ms の結果を代入する。
\item omega\_wheel に speed\_ms / float(base\_model.p.wheel\_radius) の結果を代入する。
\item torque に np.abs(p\_mech) / (omega\_wheel + 0.001) / float(base\_model.p.gear\_ratio) の結果を代入する。
\item drive\_eff に \_map\_efficiency\_array(base\_model, base\_model.maps\_drive, speed\_ms, torque, regen=False) の結果を代入する。
\item drive\_eff に np.clip(drive\_eff * float(mot.drive\_eff\_scale), 0.55, 0.99) の結果を代入する。
\item regen\_eff に \_map\_efficiency\_array(base\_model, base\_model.maps\_regen, speed\_ms, torque, regen=True) の結果を代入する。
\item regen\_eff に np.clip(regen\_eff * float(mot.drive\_eff\_scale), 0.4, 0.95) の結果を代入する。
\item positive に p\_mech >= 0.0 の結果を代入する。
\item electrical に np.empty(len(frame), dtype=float) の結果を代入する。
\item electrical[positive] に p\_mech[positive] / np.maximum(drive\_eff[positive] * base\_model.p.gear\_eta * base\_model.p.inverter\_eta, 1e-06) の結果を代入する。
\item electrical[\textasciitilde{}positive] に -float(np.clip(mot.regen\_utilization, 0.0, 1.0)) * regen\_eff[\textasciitilde{}positive] * base\_model.p.gear\_eta * base\_model.p.inverter\_eta * np.abs(p\_mech[\textasciitilde{}positive]) の結果を代入する。
\item electrical + float(mot.p\_aux\_w) - solar を返す。
            \end{enumerate}
            \paragraph{代表コード断片}
            \begin{lstlisting}[language=Python]
def _motion_predictions_for_acceleration(
    frame: pd.DataFrame,
    acceleration: pd.Series,
    base_model,
    mot: MotionFitResult,
) -> np.ndarray:
    """Vectorized equivalent of motion_power_prediction for filter/lag search."""
    speed_ms = (
        pd.to_numeric(frame["speed_kmh"], errors="coerce").fillna(0.0).to_numpy()
        / 3.6
    )
    slope = pd.to_numeric(frame["slope_pct"], errors="coerce").fillna(0.0).to_numpy()
    headwind = (
        pd.to_numeric(frame["headwind_archive_ms"], errors="coerce").fillna(0.0).to_numpy()
        * float(mot.headwind_gain)
    )
    solar = (
        pd.to_numeric(frame["solar_power_w_obs"], errors="coerce").fillna(0.0).to_numpy()
    )
    theta = np.arctan((slope * float(mot.grade_scale)) / 100.0)
    relative_speed = np.maximum(0.0, speed_ms + headwind)
    ambient_source = frame["Tamb_archive_C"] if "Tamb_archive_C" in frame else pd.Series(25.0, index=frame.index)
    elevation_source = frame["alt_m"] if "alt_m" in frame else pd.Series(0.0, index=frame.index)
    ambient = pd.to_numeric(ambient_source, errors="coerce").fillna(25.0).to_numpy()
    elevation = pd.to_numeric(elevation_source, errors="coerce").fillna(0.0).to_numpy()
    if str(base_model.p.air_density_mode or "constant").lower() == "ideal_gas_altitude":
        altitude = np.clip(elevation, -500.0, 11000.0)
        pressure_ratio = np.maximum(0.05, 1.0 - 0.0065 * altitude / 288.15) ** 5.255877
        pressure = float(base_model.p.air_density_reference_pressure_pa) * pressure_ratio
        air_density = pressure / (287.05 * np.maximum(180.0, ambient + 273.15))
    else:
        air_density = np.full(len(frame), float(base_model.p.rho), dtype=float)
    normal_force = base_model.p.m * base_model.p.g * np.cos(theta)
    force = (
        0.5 * air_density * float(mot.cda) * relative_speed**2
...
\end{lstlisting}

\subsection{L1431 関数 fit\_acceleration\_timestamp\_alignment}
            \begin{itemize}[leftmargin=1.6em]
              \item 行範囲: L1431--L1628
              \item 所属: module直下
              \item 定義: \path!fit_acceleration_timestamp_alignment(frame: pd.DataFrame, base_model, mot: MotionFitResult, options: dict | None = None) -> tuple[pd.DataFrame, Dict[str, Any]]!
              \item docstring: Identify the filter and timestamp offset of GPS-derived acceleration.

Vehicle mass remains fixed. Candidate observation filters and offsets are
fitted on all but the last race day and adopted only when the held-out last
day does not regress. This models quantized/asynchronous GNSS observations;
it is not a tunable vehicle force or a live-command filter.
              \item このブロックが直接呼ぶ主な関数・メソッド: \_acceleration\_trace\_from\_speed, \_motion\_predictions\_for\_acceleration, \_shift\_acceleration\_within\_segments, abs, any, append, astype, bool, clip, copy, fillna, float
              \item 戻り値の要点: (out, \{'enabled': True, 'adopted': adopted, 'method': 'fixed-mass GNSS acceleration filter/lag selection with last-race-day holdout', 'sample\_period\_sec': sample\_period\_sec, 'selected\_filter\_method': selected\_method, 'selected\_filter\_window\_samples': selected\_window, 'selected\_filter\_window\_sec': selected\_window * sample\_period\_sec, 'selected\_lag\_sec': selected\_lag, 'lag\_search\_min\_sec': lag\_search\_min, 'lag\_search\_max\_sec': lag\_search\_max, 'lag\_search\_boundary\_hit': lag\_boundary\_hit, 'holdout\_day': holdout\_day, 'training\_sample\_count': int(train\_mask.sum()), 'validation\_sample\_count': int(validation\_mask.sum()), 'baseline\_training\_rmse\_w': float(baseline['training\_rmse\_w']), 'selected\_training\_rmse\_w': float(selected['training\_rmse\_w']), 'baseline\_validation\_rmse\_w': float(baseline['validation\_rmse\_w']), 'selected\_validation\_rmse\_w': float(selected['validation\_rmse\_w']), 'training\_improvement\_w': float(float(baseline['training\_rmse\_w']) - float(selected['training\_rmse\_w'])), 'validation\_rmse\_ratio': validation\_ratio if adopted else 1.0, 'candidates': records\}) / (frame.copy(), \{'enabled': enabled, 'adopted': False, 'reason': 'disabled\_or\_missing\_columns'\}) / (out, \{'enabled': True, 'adopted': False, 'reason': 'insufficient\_training\_samples', 'training\_sample\_count': int(train\_mask.sum()), 'validation\_sample\_count': int(validation\_mask.sum())\}) / float(np.sqrt(np.mean(np.square(observed[use] - predicted[use])))) if use.any() else float('nan')
              \item 主なローカル代入名: actual\_lag\_sec, adopted, aligned, base\_valid, baseline, best, candidate\_filters, candidate\_lag\_sec, cfg, configured\_holdout\_day, day\_values, dt\_sec, enabled, enough\_validation, filter\_method, filter\_window, filtered, finite\_days, finite\_records, holdout\_day
              \item 読み取る主なインスタンス属性: 特になし。
              \item 更新する主なインスタンス属性: 特になし。
              \item 明示的に送出する例外: 明示的raiseはない。
              \item 制御構造の規模: 条件分岐 5、ループ 3、try 0
            \end{itemize}
            \paragraph{この定義を読むためのPython構文}
            \begin{itemize}[leftmargin=1.6em]
            \item `->`以後は戻り値の型注釈であり、実行時の値を自動変換する命令ではない。
\item lambdaは名前を付けずに短い関数オブジェクトを作る構文である。
\item 内包表記は、反復・条件・値生成を一つの式で記述してcollectionを作る。
            \end{itemize}
            \paragraph{上から順の処理}
            \begin{enumerate}[leftmargin=1.8em]
            \item cfg に options if isinstance(options, dict) else \{\} の結果を代入する。
\item enabled に bool(cfg.get('enabled', True)) の結果を代入する。
\item required に \{'time\_utc', 'day', 'speed\_kmh', 'slope\_pct', 'headwind\_archive\_ms', 'solar\_power\_w\_obs', 'battery\_power\_w\_obs', 'exclude\_power\_fit', 'accel\_ms2'\} の結果を代入する。
\item 条件 not enabled or not required.issubset(frame.columns) を判定し、真なら内部処理を行う。
\item   (frame.copy(), \{'enabled': enabled, 'adopted': False, 'reason': 'disabled\_or\_missing\_columns'\}) を返す。
\item out に frame.copy() の結果を代入する。
\item out['accel\_ms2\_previous'] に pd.to\_numeric(out['accel\_ms2'], errors='coerce') の結果を代入する。
\item out['accel\_ms2\_raw'] に \_acceleration\_trace\_from\_speed(out, method='median', window\_samples=1) の結果を代入する。
\item dt\_sec に pd.to\_numeric(out.get('dt\_sec'), errors='coerce') の結果を代入する。
\item usable\_dt に dt\_sec[np.isfinite(dt\_sec) \& (dt\_sec > 0.0) \& (dt\_sec <= 60.0)] の結果を代入する。
\item sample\_period\_sec に float(np.median(usable\_dt)) if len(usable\_dt) else 5.0 の結果を代入する。
\item raw\_candidates に cfg.get('candidate\_lag\_sec', [-30.0, -20.0, -15.0, -10.0, -5.0, 0.0, 5.0, 10.0, 15.0, 20.0, 30.0]) の結果を代入する。
\item candidate\_lag\_sec に sorted(set((float(value) for value in raw\_candidates if np.isfinite(float(value)))) | \{0.0\}) の結果を代入する。
\item raw\_filters に cfg.get('candidate\_filters', [\{'method': 'median', 'window\_samples': 5\}, \{'method': 'median', 'window\_samples': 9\}, \{'method': 'median', 'window\_samples': 15\}, \{'method': 'savgol', 'window\_samples': 11\}, \{'method': 'savgol', 'window\_samples': 15\}, \{'method': 'savgol', 'window\_samples': 21\}]) の結果を代入する。
\item candidate\_filters に [] を代入する。
\item raw\_filters を順に走査し、各要素を item に入れて処理する。
\item   条件 isinstance(item, dict) を判定し、真なら内部処理を行う。
\item     method に str(item.get('method', 'median')).strip().lower() の結果を代入する。
\item     window に int(item.get('window\_samples', 5)) の結果を代入する。
\item     上の条件が偽の場合:
\item     条件 isinstance(item, (list, tuple)) and len(item) >= 2 を判定し、真なら内部処理を行う。
\item       (method, window) に (str(item[0]).strip().lower(), int(item[1])) の結果を代入する。
\item       上の条件が偽の場合:
\item       Continue 文を実行する。
\item   条件 method in \{'median', 'savgol'\} を判定し、真なら内部処理を行う。
\item     candidate\_filters.append(...) を実行する。
\item candidate\_filters に list(dict.fromkeys(candidate\_filters + [('median', 5)])) の結果を代入する。
\item day\_values に pd.to\_numeric(out['day'], errors='coerce') の結果を代入する。
\item finite\_days に sorted(set((int(value) for value in day\_values[np.isfinite(day\_values)]))) の結果を代入する。
\item configured\_holdout\_day に int(cfg.get('holdout\_day', 0) or 0) の結果を代入する。
\item holdout\_day に configured\_holdout\_day if configured\_holdout\_day > 0 else int(finite\_days[-1] if finite\_days else 0) の結果を代入する。
\item base\_valid に \textasciitilde{}out['exclude\_power\_fit'].fillna(True).astype(bool) \& (pd.to\_numeric(out['speed\_kmh'], errors='coerce') >= float(cfg.get('min\_speed\_kmh', 12.0))) \& np.isfinite(pd.to\_numeric(out['battery\_power\_w\_obs'], errors='coerce')) \& np.isfinite(pd.to\_numeric(out['solar\_power\_w\_obs'], errors='coerce')) の結果を代入する。
\item train\_mask に base\_valid \& (day\_values != holdout\_day) の結果を代入する。
\item validation\_mask に base\_valid \& (day\_values == holdout\_day) の結果を代入する。
\item 条件 int(train\_mask.sum()) < int(cfg.get('minimum\_train\_samples', 500)) を判定し、真なら内部処理を行う。
\item   (out, \{'enabled': True, 'adopted': False, 'reason': 'insufficient\_training\_samples', 'training\_sample\_count': int(train\_mask.sum()), 'validation\_sample\_count': int(validation\_mask.sum())\}) を返す。
\item observed に pd.to\_numeric(out['battery\_power\_w\_obs'], errors='coerce').to\_numpy(dtype=float) の結果を代入する。
\item records に [] を代入する。
\item traces に \{\} を代入する。
\item candidate\_filters を順に走査し、各要素を (filter\_method, filter\_window) に入れて処理する。
\item   filtered に \_acceleration\_trace\_from\_speed(out, method=filter\_method, window\_samples=filter\_window) の結果を代入する。
\item   candidate\_lag\_sec を順に走査し、各要素を lag\_sec に入れて処理する。
\item     sample\_shift に int(round(float(lag\_sec) / sample\_period\_sec)) の結果を代入する。
\item     actual\_lag\_sec に float(sample\_shift * sample\_period\_sec) の結果を代入する。
\item     aligned に \_shift\_acceleration\_within\_segments(out, sample\_shift, acceleration=filtered) の結果を代入する。
\item     predicted に \_motion\_predictions\_for\_acceleration(out, aligned, base\_model, mot) の結果を代入する。
\item     関数 rmse を定義する。
\item     records.append(...) を実行する。
\item     traces[filter\_method, int(filter\_window), actual\_lag\_sec] に aligned の結果を代入する。
\item baseline に next((row for row in records if row['filter\_method'] == 'median' and int(row['filter\_window\_samples']) == 5 and (abs(float(row['lag\_sec'])) < 0.5 * sample\_period\_sec)), min(records, key=lambda row: abs(float(row['lag\_sec'])))) の結果を代入する。
\item finite\_records に [row for row in records if np.isfinite(float(row['training\_rmse\_w']))] の結果を代入する。
\item best に min(finite\_records, key=lambda row: float(row['training\_rmse\_w'])) の結果を代入する。
\item train\_improvement に float(baseline['training\_rmse\_w']) - float(best['training\_rmse\_w']) の結果を代入する。
\item validation\_ratio に float(best['validation\_rmse\_w']) / max(float(baseline['validation\_rmse\_w']), 1e-09) if np.isfinite(float(best['validation\_rmse\_w'])) and np.isfinite(float(baseline['validation\_rmse\_w'])) else float('nan') の結果を代入する。
\item enough\_validation に int(validation\_mask.sum()) >= int(cfg.get('minimum\_validation\_samples', 100)) の結果を代入する。
\item adopted に bool(train\_improvement >= float(cfg.get('minimum\_training\_improvement\_w', 2.0)) and (not enough\_validation or (np.isfinite(validation\_ratio) and validation\_ratio <= float(cfg.get('maximum\_validation\_rmse\_ratio', 1.02))))) の結果を代入する。
\item selected に best if adopted else baseline の結果を代入する。
\item selected\_lag に float(selected['lag\_sec']) の結果を代入する。
\item selected\_method に str(selected['filter\_method']) の結果を代入する。
\item selected\_window に int(selected['filter\_window\_samples']) の結果を代入する。
\item lag\_search\_min に float(min(candidate\_lag\_sec)) の結果を代入する。
\item lag\_search\_max に float(max(candidate\_lag\_sec)) の結果を代入する。
\item lag\_boundary\_hit に bool(math.isclose(selected\_lag, lag\_search\_min, abs\_tol=0.5 * sample\_period\_sec) or math.isclose(selected\_lag, lag\_search\_max, abs\_tol=0.5 * sample\_period\_sec)) の結果を代入する。
\item out['accel\_ms2'] に traces[selected\_method, selected\_window, selected\_lag].fillna(0.0).clip(lower=-1.5, upper=1.5) の結果を代入する。
\item out['acceleration\_timestamp\_lag\_sec'] に selected\_lag の結果を代入する。
\item out['acceleration\_filter\_method'] に selected\_method の結果を代入する。
\item out['acceleration\_filter\_window\_samples'] に selected\_window の結果を代入する。
\item (out, \{'enabled': True, 'adopted': adopted, 'method': 'fixed-mass GNSS acceleration filter/lag selection with last-race-day holdout', 'sample\_period\_sec': sample\_period\_sec, 'selected\_filter\_method': selected\_method, 'selected\_filter\_window\_samples': selected\_window, 'selected\_filter\_window\_sec': selected\_window * sample\_period\_sec, 'selected\_lag\_sec': selected\_lag, 'lag\_search\_min\_sec': lag\_search\_min, 'lag\_search\_max\_sec': lag\_search\_max, 'lag\_search\_boundary\_hit': lag\_boundary\_hit, 'holdout\_day': holdout\_day, 'training\_sample\_count': int(train\_mask.sum()), 'validation\_sample\_count': int(validation\_mask.sum()), 'baseline\_training\_rmse\_w': float(baseline['training\_rmse\_w']), 'selected\_training\_rmse\_w': float(selected['training\_rmse\_w']), 'baseline\_validation\_rmse\_w': float(baseline['validation\_rmse\_w']), 'selected\_validation\_rmse\_w': float(selected['validation\_rmse\_w']), 'training\_improvement\_w': float(float(baseline['training\_rmse\_w']) - float(selected['training\_rmse\_w'])), 'validation\_rmse\_ratio': validation\_ratio if adopted else 1.0, 'candidates': records\}) を返す。
            \end{enumerate}
            \paragraph{代表コード断片}
            \begin{lstlisting}[language=Python]
def fit_acceleration_timestamp_alignment(
    frame: pd.DataFrame,
    base_model,
    mot: MotionFitResult,
    options: dict | None = None,
) -> tuple[pd.DataFrame, Dict[str, Any]]:
    """Identify the filter and timestamp offset of GPS-derived acceleration.

    Vehicle mass remains fixed. Candidate observation filters and offsets are
    fitted on all but the last race day and adopted only when the held-out last
    day does not regress. This models quantized/asynchronous GNSS observations;
    it is not a tunable vehicle force or a live-command filter.
    """
    cfg = options if isinstance(options, dict) else {}
    enabled = bool(cfg.get("enabled", True))
    required = {
        "time_utc",
        "day",
        "speed_kmh",
        "slope_pct",
        "headwind_archive_ms",
        "solar_power_w_obs",
        "battery_power_w_obs",
        "exclude_power_fit",
        "accel_ms2",
    }
    if not enabled or not required.issubset(frame.columns):
        return frame.copy(), {
            "enabled": enabled,
            "adopted": False,
            "reason": "disabled_or_missing_columns",
        }

    out = frame.copy()
    out["accel_ms2_previous"] = pd.to_numeric(out["accel_ms2"], errors="coerce")
...
\end{lstlisting}

\subsection{L1631 関数 \_grade\_from\_smoothed\_elevation}
            \begin{itemize}[leftmargin=1.6em]
              \item 行範囲: L1631--L1650
              \item 所属: module直下
              \item 定義: \path|_grade_from_smoothed_elevation(distance_km: np.ndarray, elevation_m: np.ndarray, smoothing_window_km: float) -> np.ndarray|
              \item docstring: 明示されていない。
              \item このブロックが直接呼ぶ主な関数・メソッド: diff, float, gradient, int, len, max, median, min, round, savgol\_filter
              \item 戻り値の要点: np.gradient(smoothed, distance\_km) * 0.1 / np.gradient(elevation\_m, distance\_km) * 0.1
              \item 主なローカル代入名: samples, smoothed, spacing
              \item 読み取る主なインスタンス属性: 特になし。
              \item 更新する主なインスタンス属性: 特になし。
              \item 明示的に送出する例外: 明示的raiseはない。
              \item 制御構造の規模: 条件分岐 2、ループ 0、try 0
            \end{itemize}
            \paragraph{この定義を読むためのPython構文}
            \begin{itemize}[leftmargin=1.6em]
            \item `->`以後は戻り値の型注釈であり、実行時の値を自動変換する命令ではない。
            \end{itemize}
            \paragraph{上から順の処理}
            \begin{enumerate}[leftmargin=1.8em]
            \item spacing に float(np.median(np.diff(distance\_km))) の結果を代入する。
\item samples に max(3, int(round(float(smoothing\_window\_km) / max(spacing, 1e-09)))) の結果を代入する。
\item 条件 samples \% 2 == 0 を判定し、真なら内部処理を行う。
\item   samples を Add で更新する。
\item samples に min(samples, len(elevation\_m) if len(elevation\_m) \% 2 else len(elevation\_m) - 1) の結果を代入する。
\item 条件 samples < 3 を判定し、真なら内部処理を行う。
\item   np.gradient(elevation\_m, distance\_km) * 0.1 を返す。
\item smoothed に savgol\_filter(elevation\_m, samples, min(2, samples - 1), mode='interp') の結果を代入する。
\item np.gradient(smoothed, distance\_km) * 0.1 を返す。
            \end{enumerate}
            \paragraph{代表コード断片}
            \begin{lstlisting}[language=Python]
def _grade_from_smoothed_elevation(
    distance_km: np.ndarray,
    elevation_m: np.ndarray,
    smoothing_window_km: float,
) -> np.ndarray:
    spacing = float(np.median(np.diff(distance_km)))
    samples = max(3, int(round(float(smoothing_window_km) / max(spacing, 1.0e-9))))
    if samples % 2 == 0:
        samples += 1
    samples = min(samples, len(elevation_m) if len(elevation_m) % 2 else len(elevation_m) - 1)
    if samples < 3:
        return np.gradient(elevation_m, distance_km) * 0.1
    smoothed = savgol_filter(
        elevation_m,
        samples,
        min(2, samples - 1),
        mode="interp",
    )
    # elevation [m] / distance [km] * 0.1 converts to percent grade.
    return np.gradient(smoothed, distance_km) * 0.1
\end{lstlisting}

\subsection{L1653 関数 fit\_grade\_observation\_alignment}
            \begin{itemize}[leftmargin=1.6em]
              \item 行範囲: L1653--L1883
              \item 所属: module直下
              \item 定義: \path!fit_grade_observation_alignment(frame: pd.DataFrame, route_profile_csv: Path, output_csv: Path, base_model, mot: MotionFitResult, options: dict | None = None) -> tuple[pd.DataFrame, Dict[str, Any], Path | None]!
              \item docstring: Cross-validate a DEM smoothing length and distance alignment.

This stage calibrates the route observation, not vehicle mass or resistance.
The selected route stores the unscaled DEM grade. ``grade\_scale`` remains a
separately fitted model coefficient in the following motion fit.
              \item このブロックが直接呼ぶ主な関数・メソッド: DataFrame, MotionFitResult, \_grade\_from\_smoothed\_elevation, \_motion\_predictions\_for\_acceleration, all, append, arange, astype, bool, copy, diff, difference
              \item 戻り値の要点: (out, \{'enabled': True, 'adopted': True, 'reason': 'training\_improvement\_and\_last\_day\_holdout\_passed', 'method': 'Savitzky-Golay DEM elevation differentiation with last-day holdout', 'source\_route\_profile\_csv': str(route\_profile\_csv), 'adopted\_route\_profile\_csv': str(output\_csv), 'holdout\_day': holdout\_day, 'training\_sample\_count': int(train\_mask.sum()), 'validation\_sample\_count': int(validation\_mask.sum()), 'selected\_smoothing\_window\_km': selected\_window, 'smoothing\_search\_max\_km': smoothing\_search\_max, 'smoothing\_search\_boundary\_hit': smoothing\_boundary\_hit, 'selected\_distance\_offset\_km': selected\_offset, 'selected\_provisional\_grade\_scale': float(best['provisional\_grade\_scale']), 'baseline\_training\_rmse\_w': baseline\_training\_rmse, 'selected\_training\_rmse\_w': float(best['training\_rmse\_w']), 'training\_improvement\_w': training\_improvement, 'baseline\_validation\_rmse\_w': baseline\_validation\_rmse, 'selected\_validation\_rmse\_w': float(best['validation\_rmse\_w']), 'validation\_rmse\_ratio': validation\_ratio, 'candidate\_count': len(records), 'top\_candidates': top\_candidates\}, output\_csv) / (frame.copy(), \{'enabled': bool(cfg.get('enabled', True)), 'adopted': False, 'reason': 'disabled\_or\_route\_profile\_missing'\}, None) / (frame.copy(), \{'enabled': True, 'adopted': False, 'reason': 'route\_profile\_requires\_distance\_and\_elevation', 'route\_profile\_csv': str(route\_profile\_csv)\}, None) / (frame.copy(), \{'enabled': True, 'adopted': False, 'reason': 'insufficient\_monotonic\_route\_elevation\_samples', 'sample\_count': int(len(source))\}, None)
              \item 主なローカル代入名: acceleration, adopted, adopted\_route, aligned\_route\_slope, baseline\_predicted, baseline\_training\_rmse, baseline\_validation\_rmse, best, candidate\_frame, candidate\_mot, cfg, common, distance, distance\_column, elevation, elevation\_column, grade\_scale, grade\_scales, holdout\_day, key
              \item 読み取る主なインスタンス属性: 特になし。
              \item 更新する主なインスタンス属性: 特になし。
              \item 明示的に送出する例外: 明示的raiseはない。
              \item 制御構造の規模: 条件分岐 6、ループ 3、try 0
            \end{itemize}
            \paragraph{この定義を読むためのPython構文}
            \begin{itemize}[leftmargin=1.6em]
            \item `->`以後は戻り値の型注釈であり、実行時の値を自動変換する命令ではない。
\item lambdaは名前を付けずに短い関数オブジェクトを作る構文である。
\item 内包表記は、反復・条件・値生成を一つの式で記述してcollectionを作る。
            \end{itemize}
            \paragraph{上から順の処理}
            \begin{enumerate}[leftmargin=1.8em]
            \item cfg に options if isinstance(options, dict) else \{\} の結果を代入する。
\item 条件 not bool(cfg.get('enabled', True)) or not route\_profile\_csv.is\_file() を判定し、真なら内部処理を行う。
\item   (frame.copy(), \{'enabled': bool(cfg.get('enabled', True)), 'adopted': False, 'reason': 'disabled\_or\_route\_profile\_missing'\}, None) を返す。
\item route に pd.read\_csv(route\_profile\_csv) の結果を代入する。
\item distance\_column に next((key for key in ('dist\_km', 's\_km', 'distance\_km') if key in route.columns), None) の結果を代入する。
\item elevation\_column に next((key for key in ('elev\_m', 'elev\_dem\_m', 'alt\_m', 'altitude\_m') if key in route.columns), None) の結果を代入する。
\item 条件 distance\_column is None or elevation\_column is None を判定し、真なら内部処理を行う。
\item   (frame.copy(), \{'enabled': True, 'adopted': False, 'reason': 'route\_profile\_requires\_distance\_and\_elevation', 'route\_profile\_csv': str(route\_profile\_csv)\}, None) を返す。
\item source に pd.DataFrame(\{'dist\_km': pd.to\_numeric(route[distance\_column], errors='coerce'), 'elev\_m': pd.to\_numeric(route[elevation\_column], errors='coerce')\}).dropna() の結果を代入する。
\item source に source.groupby('dist\_km', as\_index=False).mean().sort\_values('dist\_km') の結果を代入する。
\item 条件 len(source) < 21 or not np.all(np.diff(source['dist\_km'].to\_numpy()) > 0.0) を判定し、真なら内部処理を行う。
\item   (frame.copy(), \{'enabled': True, 'adopted': False, 'reason': 'insufficient\_monotonic\_route\_elevation\_samples', 'sample\_count': int(len(source))\}, None) を返す。
\item required に \{'day', 's\_km', 'speed\_kmh', 'accel\_ms2', 'headwind\_archive\_ms', 'solar\_power\_w\_obs', 'battery\_power\_w\_obs', 'exclude\_power\_fit', 'slope\_pct'\} の結果を代入する。
\item 条件 not required.issubset(frame.columns) を判定し、真なら内部処理を行う。
\item   (frame.copy(), \{'enabled': True, 'adopted': False, 'reason': 'replay\_missing\_grade\_validation\_columns', 'missing\_columns': sorted(required.difference(frame.columns))\}, None) を返す。
\item out に frame.copy() の結果を代入する。
\item observed に pd.to\_numeric(out['battery\_power\_w\_obs'], errors='coerce').to\_numpy(dtype=float) の結果を代入する。
\item common に \textasciitilde{}out['exclude\_power\_fit'].fillna(False).astype(bool) \& pd.to\_numeric(out['speed\_kmh'], errors='coerce').ge(12.0) \& np.isfinite(observed) \& np.isfinite(pd.to\_numeric(out['s\_km'], errors='coerce')) の結果を代入する。
\item holdout\_day に int(cfg.get('holdout\_day', pd.to\_numeric(out['day'], errors='coerce').max())) の結果を代入する。
\item train\_mask に (common \& pd.to\_numeric(out['day'], errors='coerce').ne(holdout\_day)).to\_numpy(dtype=bool) の結果を代入する。
\item validation\_mask に (common \& pd.to\_numeric(out['day'], errors='coerce').eq(holdout\_day)).to\_numpy(dtype=bool) の結果を代入する。
\item min\_train に int(cfg.get('minimum\_train\_samples', 500)) の結果を代入する。
\item min\_validation に int(cfg.get('minimum\_validation\_samples', 100)) の結果を代入する。
\item 条件 int(train\_mask.sum()) < min\_train or int(validation\_mask.sum()) < min\_validation を判定し、真なら内部処理を行う。
\item   (out, \{'enabled': True, 'adopted': False, 'reason': 'insufficient\_train\_or\_validation\_samples', 'training\_sample\_count': int(train\_mask.sum()), 'validation\_sample\_count': int(validation\_mask.sum())\}, None) を返す。
\item acceleration に pd.to\_numeric(out['accel\_ms2'], errors='coerce').fillna(0.0) の結果を代入する。
\item 関数 rmse を定義する。
\item baseline\_predicted に \_motion\_predictions\_for\_acceleration(out, acceleration, base\_model, mot) の結果を代入する。
\item baseline\_training\_rmse に rmse(baseline\_predicted, train\_mask) の結果を代入する。
\item baseline\_validation\_rmse に rmse(baseline\_predicted, validation\_mask) の結果を代入する。
\item windows に [float(value) for value in cfg.get('smoothing\_windows\_km', [0.5, 1.1, 2.1, 3.1, 5.1])] の結果を代入する。
\item offsets に [float(value) for value in cfg.get('distance\_offsets\_km', np.round(np.arange(-0.5, 0.5001, 0.1), 6).tolist())] の結果を代入する。
\item grade\_scales に [float(value) for value in cfg.get('grade\_scales', np.round(np.arange(0.5, 1.0001, 0.05), 6).tolist())] の結果を代入する。
\item distance に source['dist\_km'].to\_numpy(dtype=float) の結果を代入する。
\item elevation に source['elev\_m'].to\_numpy(dtype=float) の結果を代入する。
\item replay\_distance に pd.to\_numeric(out['s\_km'], errors='coerce').to\_numpy(dtype=float) の結果を代入する。
\item records に [] を代入する。
\item route\_slopes に \{\} を代入する。
\item windows を順に走査し、各要素を window\_km に入れて処理する。
\item   source\_slope に \_grade\_from\_smoothed\_elevation(distance, elevation, window\_km) の結果を代入する。
\item   offsets を順に走査し、各要素を offset\_km に入れて処理する。
\item     aligned\_route\_slope に np.interp(distance + offset\_km, distance, source\_slope, left=source\_slope[0], right=source\_slope[-1]) の結果を代入する。
\item     route\_slopes[window\_km, offset\_km] に aligned\_route\_slope の結果を代入する。
\item     candidate\_frame に out.copy() の結果を代入する。
\item     candidate\_frame['slope\_pct'] に np.interp(replay\_distance, distance, aligned\_route\_slope, left=aligned\_route\_slope[0], right=aligned\_route\_slope[-1]) の結果を代入する。
\item     grade\_scales を順に走査し、各要素を grade\_scale に入れて処理する。
\item       candidate\_mot に MotionFitResult(**\{**mot.\_\_dict\_\_, 'grade\_scale': float(grade\_scale)\}) の結果を代入する。
\item       predicted に \_motion\_predictions\_for\_acceleration(candidate\_frame, acceleration, base\_model, candidate\_mot) の結果を代入する。
\item       records.append(...) を実行する。
\item best に min(records, key=lambda row: float(row['training\_rmse\_w'])) の結果を代入する。
\item training\_improvement に baseline\_training\_rmse - float(best['training\_rmse\_w']) の結果を代入する。
\item validation\_ratio に float(best['validation\_rmse\_w']) / max(baseline\_validation\_rmse, 1e-09) の結果を代入する。
\item adopted に bool(training\_improvement >= float(cfg.get('minimum\_training\_improvement\_w', 5.0)) and validation\_ratio <= float(cfg.get('maximum\_validation\_rmse\_ratio', 1.02))) の結果を代入する。
\item 条件 not adopted を判定し、真なら内部処理を行う。
\item   (out, \{'enabled': True, 'adopted': False, 'reason': 'training\_gain\_or\_holdout\_gate\_failed', 'method': 'Savitzky-Golay DEM elevation differentiation with last-day holdout', 'baseline\_training\_rmse\_w': baseline\_training\_rmse, 'selected\_training\_rmse\_w': float(best['training\_rmse\_w']), 'baseline\_validation\_rmse\_w': baseline\_validation\_rmse, 'selected\_validation\_rmse\_w': float(best['validation\_rmse\_w']), 'validation\_rmse\_ratio': validation\_ratio, 'candidate\_count': len(records)\}, None) を返す。
\item selected\_window に float(best['smoothing\_window\_km']) の結果を代入する。
\item selected\_offset に float(best['distance\_offset\_km']) の結果を代入する。
\item smoothing\_search\_max に float(max(windows)) の結果を代入する。
\item smoothing\_boundary\_hit に math.isclose(selected\_window, smoothing\_search\_max, rel\_tol=0.0, abs\_tol=1e-09) の結果を代入する。
\item selected\_route\_slope に route\_slopes[selected\_window, selected\_offset] の結果を代入する。
\item out['slope\_pct\_previous'] に pd.to\_numeric(out['slope\_pct'], errors='coerce') の結果を代入する。
\item out['slope\_pct'] に np.interp(replay\_distance, distance, selected\_route\_slope, left=selected\_route\_slope[0], right=selected\_route\_slope[-1]) の結果を代入する。
\item adopted\_route に pd.DataFrame(\{'dist\_km': distance, 'elev\_m': elevation, 'slope\_pct': selected\_route\_slope, 'headwind\_ms': np.zeros(len(distance), dtype=float)\}) の結果を代入する。
\item ensure\_dir(...) を実行する。
\item adopted\_route.to\_csv(...) を実行する。
\item top\_candidates に sorted(records, key=lambda row: float(row['training\_rmse\_w']))[:25] の結果を代入する。
\item (out, \{'enabled': True, 'adopted': True, 'reason': 'training\_improvement\_and\_last\_day\_holdout\_passed', 'method': 'Savitzky-Golay DEM elevation differentiation with last-day holdout', 'source\_route\_profile\_csv': str(route\_profile\_csv), 'adopted\_route\_profile\_csv': str(output\_csv), 'holdout\_day': holdout\_day, 'training\_sample\_count': int(train\_mask.sum()), 'validation\_sample\_count': int(validation\_mask.sum()), 'selected\_smoothing\_window\_km': selected\_window, 'smoothing\_search\_max\_km': smoothing\_search\_max, 'smoothing\_search\_boundary\_hit': smoothing\_boundary\_hit, 'selected\_distance\_offset\_km': selected\_offset, 'selected\_provisional\_grade\_scale': float(best['provisional\_grade\_scale']), 'baseline\_training\_rmse\_w': baseline\_training\_rmse, 'selected\_training\_rmse\_w': float(best['training\_rmse\_w']), 'training\_improvement\_w': training\_improvement, 'baseline\_validation\_rmse\_w': baseline\_validation\_rmse, 'selected\_validation\_rmse\_w': float(best['validation\_rmse\_w']), 'validation\_rmse\_ratio': validation\_ratio, 'candidate\_count': len(records), 'top\_candidates': top\_candidates\}, output\_csv) を返す。
            \end{enumerate}
            \paragraph{代表コード断片}
            \begin{lstlisting}[language=Python]
def fit_grade_observation_alignment(
    frame: pd.DataFrame,
    route_profile_csv: Path,
    output_csv: Path,
    base_model,
    mot: MotionFitResult,
    options: dict | None = None,
) -> tuple[pd.DataFrame, Dict[str, Any], Path | None]:
    """Cross-validate a DEM smoothing length and distance alignment.

    This stage calibrates the route observation, not vehicle mass or resistance.
    The selected route stores the unscaled DEM grade. ``grade_scale`` remains a
    separately fitted model coefficient in the following motion fit.
    """
    cfg = options if isinstance(options, dict) else {}
    if not bool(cfg.get("enabled", True)) or not route_profile_csv.is_file():
        return frame.copy(), {
            "enabled": bool(cfg.get("enabled", True)),
            "adopted": False,
            "reason": "disabled_or_route_profile_missing",
        }, None

    route = pd.read_csv(route_profile_csv)
    distance_column = next(
        (key for key in ("dist_km", "s_km", "distance_km") if key in route.columns),
        None,
    )
    elevation_column = next(
        (key for key in ("elev_m", "elev_dem_m", "alt_m", "altitude_m") if key in route.columns),
        None,
    )
    if distance_column is None or elevation_column is None:
        return frame.copy(), {
            "enabled": True,
            "adopted": False,
...
\end{lstlisting}

\subsection{L1750 関数 fit\_grade\_observation\_alignment.rmse}
            \begin{itemize}[leftmargin=1.6em]
              \item 行範囲: L1750--L1752
              \item 所属: \path|fit_grade_observation_alignment|
              \item 定義: \path|rmse(predicted: np.ndarray, mask: np.ndarray) -> float|
              \item docstring: 明示されていない。
              \item このブロックが直接呼ぶ主な関数・メソッド: float, isfinite, mean, sqrt, square
              \item 戻り値の要点: float(np.sqrt(np.mean(np.square(observed[usable] - predicted[usable]))))
              \item 主なローカル代入名: usable
              \item 読み取る主なインスタンス属性: 特になし。
              \item 更新する主なインスタンス属性: 特になし。
              \item 明示的に送出する例外: 明示的raiseはない。
              \item 制御構造の規模: 条件分岐 0、ループ 0、try 0
            \end{itemize}
            \paragraph{この定義を読むためのPython構文}
            \begin{itemize}[leftmargin=1.6em]
            \item `->`以後は戻り値の型注釈であり、実行時の値を自動変換する命令ではない。
            \end{itemize}
            \paragraph{上から順の処理}
            \begin{enumerate}[leftmargin=1.8em]
            \item usable に mask \& np.isfinite(predicted) \& np.isfinite(observed) の結果を代入する。
\item float(np.sqrt(np.mean(np.square(observed[usable] - predicted[usable])))) を返す。
            \end{enumerate}
            \paragraph{代表コード断片}
            \begin{lstlisting}[language=Python]
    def rmse(predicted: np.ndarray, mask: np.ndarray) -> float:
        usable = mask & np.isfinite(predicted) & np.isfinite(observed)
        return float(np.sqrt(np.mean(np.square(observed[usable] - predicted[usable]))))
\end{lstlisting}

\subsection{L1886 関数 pv\_leave\_one\_day\_out\_validation}
            \begin{itemize}[leftmargin=1.6em]
              \item 行範囲: L1886--L1961
              \item 所属: module直下
              \item 定義: \path!pv_leave_one_day_out_validation(frame: pd.DataFrame, model, *, panel_deployment_options: Dict[str, float] | None = None) -> Dict[str, float]!
              \item docstring: Validate the PV chain on days excluded from all PV-scalar fitting.
              \item このブロックが直接呼ぶ主な関数・メソッド: ValueError, asarray, astype, attach\_archive\_pv\_model, copy, dict, difference, drop, drop\_duplicates, dropna, extend, fillna
              \item 戻り値の要点: \{'pv\_lodo\_fold\_count': int(fold\_count), 'pv\_lodo\_moving\_rmse\_w': rmse(moving\_residuals), 'pv\_lodo\_moving\_sample\_count': int(len(moving\_residuals)), 'pv\_lodo\_deployed\_stop\_rmse\_w': rmse(deployed\_residuals), 'pv\_lodo\_deployed\_stop\_sample\_count': int(len(deployed\_residuals))\} / float(np.sqrt(np.mean(np.square(array)))) if array.size else float('nan')
              \item 主なローカル代入名: array, day\_value, deployed, deployed\_residuals, deployment, fold\_count, fold\_fit, holdout, keep, local\_time, missing, modeled, moving, moving\_residuals, observed, predicted, required, train, valid, work
              \item 読み取る主なインスタンス属性: 特になし。
              \item 更新する主なインスタンス属性: 特になし。
              \item 明示的に送出する例外: ValueError(f'PV leave-one-day-out validation is missing columns: \{missing\}')
              \item 制御構造の規模: 条件分岐 4、ループ 1、try 1
            \end{itemize}
            \paragraph{この定義を読むためのPython構文}
            \begin{itemize}[leftmargin=1.6em]
            \item `->`以後は戻り値の型注釈であり、実行時の値を自動変換する命令ではない。
\item try文は例外が起き得る処理と、異常時または終了時の経路を分ける。
            \end{itemize}
            \paragraph{上から順の処理}
            \begin{enumerate}[leftmargin=1.8em]
            \item required に \{'time\_utc', 'speed\_kmh', 'GHI\_archive', 'DNI\_archive', 'DHI\_archive', 'Tamb\_archive\_C', 'solar\_power\_w\_obs', 'exclude\_weather\_fit'\} の結果を代入する。
\item missing に sorted(required.difference(frame.columns)) の結果を代入する。
\item 条件 missing を判定し、真なら内部処理を行う。
\item   ValueError(f'PV leave-one-day-out validation is missing columns: \{missing\}') を送出する。
\item keep に sorted(required.union(\{'day', 's\_km'\}).intersection(frame.columns)) の結果を代入する。
\item work に frame.loc[:, keep].copy() の結果を代入する。
\item 条件 'day' in work.columns を判定し、真なら内部処理を行う。
\item   work['\_pv\_cv\_day'] に pd.to\_numeric(work['day'], errors='coerce') の結果を代入する。
\item   上の条件が偽の場合:
\item   local\_time に pd.to\_datetime(work['time\_utc'], format='mixed', utc=True).dt.tz\_convert(TIMEZONE\_LOCAL) の結果を代入する。
\item   work['\_pv\_cv\_day'] に local\_time.dt.date.astype(str) の結果を代入する。
\item deployment に dict(panel\_deployment\_options or \{\}) の結果を代入する。
\item moving\_residuals に [] の結果を代入する。
\item deployed\_residuals に [] の結果を代入する。
\item fold\_count に 0 の結果を代入する。
\item work['\_pv\_cv\_day'].dropna().drop\_duplicates().tolist() を順に走査し、各要素を day\_value に入れて処理する。
\item   holdout に work.loc[work['\_pv\_cv\_day'] == day\_value].drop(columns='\_pv\_cv\_day') の結果を代入する。
\item   train に work.loc[work['\_pv\_cv\_day'] != day\_value].drop(columns='\_pv\_cv\_day') の結果を代入する。
\item   条件 len(train) < 100 or len(holdout) < 20 を判定し、真なら内部処理を行う。
\item     Continue 文を実行する。
\item   例外処理を伴う try ブロックを実行する。
\item     fold\_fit に fit\_pv\_parameters(train, model, irradiance\_source='GHI\_archive', operating\_state='moving') の結果を代入する。
\item     fold\_fit に fit\_stop\_tilt\_fraction(train, model, fold\_fit, **deployment) の結果を代入する。
\item     predicted に attach\_archive\_pv\_model(holdout, model, fold\_fit, irradiance\_source='GHI\_archive', **deployment) の結果を代入する。
\item     (KeyError, ValueError, RuntimeError)を捕捉した場合:
\item     Continue 文を実行する。
\item   observed に pd.to\_numeric(predicted['solar\_power\_w\_obs'], errors='coerce') の結果を代入する。
\item   modeled に pd.to\_numeric(predicted['solar\_power\_w\_model'], errors='coerce') の結果を代入する。
\item   valid に observed.notna() \& modeled.notna() の結果を代入する。
\item   条件 'exclude\_weather\_fit' in predicted.columns を判定し、真なら内部処理を行う。
\item     valid を BitAnd で更新する。
\item   moving に valid \& pd.to\_numeric(predicted['speed\_kmh'], errors='coerce').ge(12.0) の結果を代入する。
\item   deployed に valid \& predicted['panel\_deployed\_model'].fillna(False).astype(bool) の結果を代入する。
\item   moving\_residuals.extend(...) を実行する。
\item   deployed\_residuals.extend(...) を実行する。
\item   fold\_count を Add で更新する。
\item 関数 rmse を定義する。
\item \{'pv\_lodo\_fold\_count': int(fold\_count), 'pv\_lodo\_moving\_rmse\_w': rmse(moving\_residuals), 'pv\_lodo\_moving\_sample\_count': int(len(moving\_residuals)), 'pv\_lodo\_deployed\_stop\_rmse\_w': rmse(deployed\_residuals), 'pv\_lodo\_deployed\_stop\_sample\_count': int(len(deployed\_residuals))\} を返す。
            \end{enumerate}
            \paragraph{代表コード断片}
            \begin{lstlisting}[language=Python]
def pv_leave_one_day_out_validation(
    frame: pd.DataFrame,
    model,
    *,
    panel_deployment_options: Dict[str, float] | None = None,
) -> Dict[str, float]:
    """Validate the PV chain on days excluded from all PV-scalar fitting."""
    required = {
        "time_utc",
        "speed_kmh",
        "GHI_archive",
        "DNI_archive",
        "DHI_archive",
        "Tamb_archive_C",
        "solar_power_w_obs",
        "exclude_weather_fit",
    }
    missing = sorted(required.difference(frame.columns))
    if missing:
        raise ValueError(f"PV leave-one-day-out validation is missing columns: {missing}")
    keep = sorted(required.union({"day", "s_km"}).intersection(frame.columns))
    work = frame.loc[:, keep].copy()
    if "day" in work.columns:
        work["_pv_cv_day"] = pd.to_numeric(work["day"], errors="coerce")
    else:
        local_time = pd.to_datetime(work["time_utc"], format="mixed", utc=True).dt.tz_convert(TIMEZONE_LOCAL)
        work["_pv_cv_day"] = local_time.dt.date.astype(str)
    deployment = dict(panel_deployment_options or {})
    moving_residuals = []
    deployed_residuals = []
    fold_count = 0
    for day_value in work["_pv_cv_day"].dropna().drop_duplicates().tolist():
        holdout = work.loc[work["_pv_cv_day"] == day_value].drop(columns="_pv_cv_day")
        train = work.loc[work["_pv_cv_day"] != day_value].drop(columns="_pv_cv_day")
        if len(train) < 100 or len(holdout) < 20:
...
\end{lstlisting}

\subsection{L1950 関数 pv\_leave\_one\_day\_out\_validation.rmse}
            \begin{itemize}[leftmargin=1.6em]
              \item 行範囲: L1950--L1953
              \item 所属: \path|pv_leave_one_day_out_validation|
              \item 定義: \path|rmse(values) -> float|
              \item docstring: 明示されていない。
              \item このブロックが直接呼ぶ主な関数・メソッド: asarray, float, isfinite, mean, sqrt, square
              \item 戻り値の要点: float(np.sqrt(np.mean(np.square(array)))) if array.size else float('nan')
              \item 主なローカル代入名: array
              \item 読み取る主なインスタンス属性: 特になし。
              \item 更新する主なインスタンス属性: 特になし。
              \item 明示的に送出する例外: 明示的raiseはない。
              \item 制御構造の規模: 条件分岐 0、ループ 0、try 0
            \end{itemize}
            \paragraph{この定義を読むためのPython構文}
            \begin{itemize}[leftmargin=1.6em]
            \item `->`以後は戻り値の型注釈であり、実行時の値を自動変換する命令ではない。
            \end{itemize}
            \paragraph{上から順の処理}
            \begin{enumerate}[leftmargin=1.8em]
            \item array に np.asarray(values, dtype=float) の結果を代入する。
\item array に array[np.isfinite(array)] の結果を代入する。
\item float(np.sqrt(np.mean(np.square(array)))) if array.size else float('nan') を返す。
            \end{enumerate}
            \paragraph{代表コード断片}
            \begin{lstlisting}[language=Python]
    def rmse(values) -> float:
        array = np.asarray(values, dtype=float)
        array = array[np.isfinite(array)]
        return float(np.sqrt(np.mean(np.square(array)))) if array.size else float("nan")
\end{lstlisting}

\subsection{L1964 関数 add\_end\_to\_end\_metrics}
            \begin{itemize}[leftmargin=1.6em]
              \item 行範囲: L1964--L1970
              \item 所属: module直下
              \item 定義: \path|add_end_to_end_metrics(primary: Dict[str, float], end_to_end: Dict[str, float]) -> Dict[str, float]|
              \item docstring: 明示されていない。
              \item このブロックが直接呼ぶ主な関数・メソッド: dict, items
              \item 戻り値の要点: out
              \item 主なローカル代入名: key, out, value
              \item 読み取る主なインスタンス属性: 特になし。
              \item 更新する主なインスタンス属性: 特になし。
              \item 明示的に送出する例外: 明示的raiseはない。
              \item 制御構造の規模: 条件分岐 0、ループ 1、try 0
            \end{itemize}
            \paragraph{この定義を読むためのPython構文}
            \begin{itemize}[leftmargin=1.6em]
            \item `->`以後は戻り値の型注釈であり、実行時の値を自動変換する命令ではない。
            \end{itemize}
            \paragraph{上から順の処理}
            \begin{enumerate}[leftmargin=1.8em]
            \item out に dict(primary) の結果を代入する。
\item end\_to\_end.items() を順に走査し、各要素を (key, value) に入れて処理する。
\item   out[f'end\_to\_end\_\{key\}'] に value の結果を代入する。
\item out['vehicle\_fit\_solar\_source'] に 'measured\_when\_available' の結果を代入する。
\item out['end\_to\_end\_solar\_source'] に 'weather\_and\_pv\_model' の結果を代入する。
\item out を返す。
            \end{enumerate}
            \paragraph{代表コード断片}
            \begin{lstlisting}[language=Python]
def add_end_to_end_metrics(primary: Dict[str, float], end_to_end: Dict[str, float]) -> Dict[str, float]:
    out = dict(primary)
    for key, value in end_to_end.items():
        out[f"end_to_end_{key}"] = value
    out["vehicle_fit_solar_source"] = "measured_when_available"
    out["end_to_end_solar_source"] = "weather_and_pv_model"
    return out
\end{lstlisting}

\subsection{L1973 関数 add\_battery\_conditional\_metrics}
            \begin{itemize}[leftmargin=1.6em]
              \item 行範囲: L1973--L1978
              \item 所属: module直下
              \item 定義: \path|add_battery_conditional_metrics(primary: Dict[str, float], conditional: Dict[str, float]) -> Dict[str, float]|
              \item docstring: 明示されていない。
              \item このブロックが直接呼ぶ主な関数・メソッド: dict, items
              \item 戻り値の要点: out
              \item 主なローカル代入名: key, out, value
              \item 読み取る主なインスタンス属性: 特になし。
              \item 更新する主なインスタンス属性: 特になし。
              \item 明示的に送出する例外: 明示的raiseはない。
              \item 制御構造の規模: 条件分岐 0、ループ 1、try 0
            \end{itemize}
            \paragraph{この定義を読むためのPython構文}
            \begin{itemize}[leftmargin=1.6em]
            \item `->`以後は戻り値の型注釈であり、実行時の値を自動変換する命令ではない。
            \end{itemize}
            \paragraph{上から順の処理}
            \begin{enumerate}[leftmargin=1.8em]
            \item out に dict(primary) の結果を代入する。
\item conditional.items() を順に走査し、各要素を (key, value) に入れて処理する。
\item   out[f'battery\_conditional\_\{key\}'] に value の結果を代入する。
\item out['battery\_conditional\_source'] に 'observed\_pack\_power\_and\_current' の結果を代入する。
\item out を返す。
            \end{enumerate}
            \paragraph{代表コード断片}
            \begin{lstlisting}[language=Python]
def add_battery_conditional_metrics(primary: Dict[str, float], conditional: Dict[str, float]) -> Dict[str, float]:
    out = dict(primary)
    for key, value in conditional.items():
        out[f"battery_conditional_{key}"] = value
    out["battery_conditional_source"] = "observed_pack_power_and_current"
    return out
\end{lstlisting}

\subsection{L1981 関数 write\_replay\_csv}
            \begin{itemize}[leftmargin=1.6em]
              \item 行範囲: L1981--L1995
              \item 所属: module直下
              \item 定義: \path|write_replay_csv(frame: pd.DataFrame, output_path: Path, *, chunk_rows: int = 5000) -> None|
              \item docstring: Write large replay tables without materializing a full string copy.
              \item このブロックが直接呼ぶ主な関数・メソッド: copy, ensure\_dir, int, len, max, range, strftime, to\_csv, to\_datetime
              \item 戻り値の要点: 戻り値が無いか、return 文が明示されていない。
              \item 主なローカル代入名: chunk, rows, start
              \item 読み取る主なインスタンス属性: 特になし。
              \item 更新する主なインスタンス属性: 特になし。
              \item 明示的に送出する例外: 明示的raiseはない。
              \item 制御構造の規模: 条件分岐 0、ループ 1、try 0
            \end{itemize}
            \paragraph{この定義を読むためのPython構文}
            \begin{itemize}[leftmargin=1.6em]
            \item `->`以後は戻り値の型注釈であり、実行時の値を自動変換する命令ではない。
            \end{itemize}
            \paragraph{上から順の処理}
            \begin{enumerate}[leftmargin=1.8em]
            \item ensure\_dir(...) を実行する。
\item rows に max(1, int(chunk\_rows)) の結果を代入する。
\item range(0, len(frame), rows) を順に走査し、各要素を start に入れて処理する。
\item   chunk に frame.iloc[start:start + rows].copy() の結果を代入する。
\item   chunk['time\_utc'] に pd.to\_datetime(chunk['time\_utc'], format='mixed', utc=True).dt.strftime('\%Y-\%m-\%dT\%H:\%M:\%SZ') の結果を代入する。
\item   chunk.to\_csv(...) を実行する。
            \end{enumerate}
            \paragraph{代表コード断片}
            \begin{lstlisting}[language=Python]
def write_replay_csv(frame: pd.DataFrame, output_path: Path, *, chunk_rows: int = 5000) -> None:
    """Write large replay tables without materializing a full string copy."""
    ensure_dir(output_path.parent)
    rows = max(1, int(chunk_rows))
    for start in range(0, len(frame), rows):
        chunk = frame.iloc[start : start + rows].copy()
        chunk["time_utc"] = pd.to_datetime(
            chunk["time_utc"], format="mixed", utc=True
        ).dt.strftime("%Y-%m-%dT%H:%M:%SZ")
        chunk.to_csv(
            output_path,
            index=False,
            mode="w" if start == 0 else "a",
            header=start == 0,
        )
\end{lstlisting}

\subsection{L1998 関数 apply\_fit\_to\_cfg}
            \begin{itemize}[leftmargin=1.6em]
              \item 行範囲: L1998--L2110
              \item 所属: module直下
              \item 定義: \path!apply_fit_to_cfg(cfg: dict, *, package_dir: Path, map_assets: Dict[str, Path], pv: PvFitResult, batt: BatteryFitResult, mot: MotionFitResult, observed_log_csv: Path, battery_dynamic_fit: Dict[str, Any] | None = None, solar_measurement_calibration: Dict[str, Any] | None = None, sync_sim_soc0: bool = False) -> dict!
              \item docstring: 明示されていない。
              \item このブロックが直接呼ぶ主な関数・メソッド: Path, ValueError, clip, dropna, float, get, items, max, min, read\_csv, relpath, replace
              \item 戻り値の要点: cfg
              \item 主なローカル代入名: autocal, dynamic, identification, key, live, model, mpc, ocv\_asset, ocv\_frame, ocv\_max\_v, ocv\_min\_v, ocv\_path, ocv\_values, path, sim, soc\_max, solar\_calibration, validation\_gate, weather
              \item 読み取る主なインスタンス属性: 特になし。
              \item 更新する主なインスタンス属性: 特になし。
              \item 明示的に送出する例外: ValueError(f'adopted OCV map has no finite ocv\_v values: \{ocv\_path\}'), ValueError(f'adopted OCV map is missing ocv\_v: \{ocv\_path\}'), ValueError(f'adopted OCV maximum \{ocv\_max\_v:.6f\} V exceeds the grounded 25S pack limit \{BATTERY\_PACK\_MAX\_CHARGE\_V:.6f\} V')
              \item 制御構造の規模: 条件分岐 5、ループ 1、try 0
            \end{itemize}
            \paragraph{この定義を読むためのPython構文}
            \begin{itemize}[leftmargin=1.6em]
            \item `->`以後は戻り値の型注釈であり、実行時の値を自動変換する命令ではない。
            \end{itemize}
            \paragraph{上から順の処理}
            \begin{enumerate}[leftmargin=1.8em]
            \item cfg.setdefault(...) を実行する。
\item map\_assets.items() を順に走査し、各要素を (key, path) に入れて処理する。
\item   cfg['paths'][key] に os.path.relpath(path, package\_dir).replace('\textbackslash{}\textbackslash{}', '/') の結果を代入する。
\item cfg['paths']['progress\_reference\_csv'] に os.path.relpath(observed\_log\_csv, package\_dir).replace('\textbackslash{}\textbackslash{}', '/') の結果を代入する。
\item model に cfg.setdefault('model', \{\}) の結果を代入する。
\item model['CdA'] に round(float(mot.cda), 6) の結果を代入する。
\item model['Crr'] に round(float(mot.crr), 6) の結果を代入する。
\item model['P\_aux'] に round(float(mot.p\_aux\_w), 3) の結果を代入する。
\item model['P\_aux\_stopped'] に round(float(mot.p\_aux\_w), 3) の結果を代入する。
\item model['P\_aux\_night'] に 0.0 の結果を代入する。
\item model.setdefault(...) を実行する。
\item model['air\_density\_mode'] に 'ideal\_gas\_altitude' の結果を代入する。
\item model.setdefault(...) を実行する。
\item solar\_calibration に solar\_measurement\_calibration or \{\} の結果を代入する。
\item model['solar\_measurement\_gain\_to\_pack'] に round(float(solar\_calibration.get('gain\_to\_pack', 1.0)), 8) の結果を代入する。
\item model['panel\_gain'] に round(float(pv.panel\_gain), 6) の結果を代入する。
\item model['grade\_scale'] に round(float(mot.grade\_scale), 6) の結果を代入する。
\item model['drive\_eff\_scale'] に round(float(mot.drive\_eff\_scale), 6) の結果を代入する。
\item model['regen\_eff\_scale'] に round(float(mot.drive\_eff\_scale), 6) の結果を代入する。
\item model['regen\_utilization'] に round(float(mot.regen\_utilization), 6) の結果を代入する。
\item model['rint\_scale'] に round(float(batt.rint\_scale), 6) の結果を代入する。
\item model['r\_line\_ohm'] に round(float(batt.r\_line\_ohm), 6) の結果を代入する。
\item model['eta\_charge'] に round(float(batt.eta\_charge), 6) の結果を代入する。
\item model['E\_nom\_Wh'] に round(float(batt.e\_nom\_wh), 3) の結果を代入する。
\item model['Q\_nom\_Ah'] に round(float(batt.e\_nom\_wh) / BATTERY\_NOMINAL\_VOLTAGE\_V, 6) の結果を代入する。
\item ocv\_asset に map\_assets.get('ocv\_soc\_map') の結果を代入する。
\item ocv\_max\_v に float('-inf') の結果を代入する。
\item 条件 ocv\_asset is not None を判定し、真なら内部処理を行う。
\item   ocv\_path に Path(ocv\_asset) の結果を代入する。
\item   ocv\_frame に pd.read\_csv(ocv\_path) の結果を代入する。
\item   条件 'ocv\_v' not in ocv\_frame.columns を判定し、真なら内部処理を行う。
\item     ValueError(f'adopted OCV map is missing ocv\_v: \{ocv\_path\}') を送出する。
\item   ocv\_values に pd.to\_numeric(ocv\_frame['ocv\_v'], errors='coerce').dropna() の結果を代入する。
\item   条件 ocv\_values.empty を判定し、真なら内部処理を行う。
\item     ValueError(f'adopted OCV map has no finite ocv\_v values: \{ocv\_path\}') を送出する。
\item   ocv\_max\_v に float(ocv\_values.max()) の結果を代入する。
\item   ocv\_min\_v に float(ocv\_values.min()) の結果を代入する。
\item   条件 ocv\_max\_v > BATTERY\_PACK\_MAX\_CHARGE\_V + 1e-06 を判定し、真なら内部処理を行う。
\item     ValueError(f'adopted OCV maximum \{ocv\_max\_v:.6f\} V exceeds the grounded 25S pack limit \{BATTERY\_PACK\_MAX\_CHARGE\_V:.6f\} V') を送出する。
\item   model['V\_min'] に round(min(float(model.get('V\_min', ocv\_min\_v)), ocv\_min\_v), 3) の結果を代入する。
\item model['V\_max'] に round(BATTERY\_PACK\_MAX\_CHARGE\_V, 3) の結果を代入する。
\item dynamic に battery\_dynamic\_fit or \{\} の結果を代入する。
\item model['r\_polarization\_ohm'] に round(float(dynamic.get('r\_polarization\_ohm', 0.0)), 6) の結果を代入する。
\item model['polarization\_tau\_sec'] に round(float(dynamic.get('tau\_sec', 60.0)), 6) の結果を代入する。
\item model['headwind\_gain'] に round(float(mot.headwind\_gain), 6) の結果を代入する。
\item identification に cfg.setdefault('identification', \{\}) の結果を代入する。
\item identification['fitted\_replay\_soc0'] に round(float(batt.soc0), 6) の結果を代入する。
\item validation\_gate に identification.setdefault('validation\_gate', \{\}) の結果を代入する。
\item validation\_gate.setdefault(...) を実行する。
\item validation\_gate.setdefault(...) を実行する。
\item validation\_gate.setdefault(...) を実行する。
\item validation\_gate.setdefault(...) を実行する。
\item validation\_gate.setdefault(...) を実行する。
\item validation\_gate.setdefault(...) を実行する。
\item validation\_gate.setdefault(...) を実行する。
\item validation\_gate.setdefault(...) を実行する。
\item validation\_gate.setdefault(...) を実行する。
\item validation\_gate.setdefault(...) を実行する。
\item validation\_gate.setdefault(...) を実行する。
\item validation\_gate.setdefault(...) を実行する。
\item validation\_gate.setdefault(...) を実行する。
\item validation\_gate.setdefault(...) を実行する。
\item validation\_gate.setdefault(...) を実行する。
\item validation\_gate.setdefault(...) を実行する。
\item validation\_gate.setdefault(...) を実行する。
\item validation\_gate.setdefault(...) を実行する。
\item validation\_gate.setdefault(...) を実行する。
\item validation\_gate.setdefault(...) を実行する。
\item sim に cfg.setdefault('simulation', \{\}) の結果を代入する。
\item soc\_max に float(model.get('soc\_max', 0.98)) の結果を代入する。
\item 条件 sync\_sim\_soc0 を判定し、真なら内部処理を行う。
\item   sim['soc0'] に round(float(np.clip(batt.soc0, 0.8, soc\_max)), 4) の結果を代入する。
\item live に cfg.setdefault('live', \{\}) の結果を代入する。
\item autocal に live.setdefault('autocal', \{\}) の結果を代入する。
\item autocal['aux\_power\_w\_init'] に round(float(mot.p\_aux\_w), 3) の結果を代入する。
\item weather に live.setdefault('weather', \{\}) の結果を代入する。
\item weather['tcell\_gain'] に round(float(pv.tcell\_gain\_c\_per\_wm2), 6) の結果を代入する。
\item mpc に cfg.setdefault('mpc', \{\}) の結果を代入する。
\item mpc['stop\_tilt\_fraction'] に round(float(np.clip(pv.stop\_tilt\_fraction, 0.0, 1.0)), 6) の結果を代入する。
\item mpc.setdefault(...) を実行する。
            \end{enumerate}
            \paragraph{代表コード断片}
            \begin{lstlisting}[language=Python]
def apply_fit_to_cfg(
    cfg: dict,
    *,
    package_dir: Path,
    map_assets: Dict[str, Path],
    pv: PvFitResult,
    batt: BatteryFitResult,
    mot: MotionFitResult,
    observed_log_csv: Path,
    battery_dynamic_fit: Dict[str, Any] | None = None,
    solar_measurement_calibration: Dict[str, Any] | None = None,
    sync_sim_soc0: bool = False,
) -> dict:
    cfg.setdefault("paths", {})
    for key, path in map_assets.items():
        cfg["paths"][key] = os.path.relpath(path, package_dir).replace("\\", "/")
    cfg["paths"]["progress_reference_csv"] = os.path.relpath(observed_log_csv, package_dir).replace("\\", "/")

    model = cfg.setdefault("model", {})
    model["CdA"] = round(float(mot.cda), 6)
    model["Crr"] = round(float(mot.crr), 6)
    model["P_aux"] = round(float(mot.p_aux_w), 3)
    model["P_aux_stopped"] = round(float(mot.p_aux_w), 3)
    model["P_aux_night"] = 0.0
    model.setdefault("aux_night_ghi_threshold_wm2", 20.0)
    model["air_density_mode"] = "ideal_gas_altitude"
    model.setdefault("air_density_reference_pressure_pa", 101325.0)
    solar_calibration = solar_measurement_calibration or {}
    model["solar_measurement_gain_to_pack"] = round(
        float(solar_calibration.get("gain_to_pack", 1.0)), 8
    )
    model["panel_gain"] = round(float(pv.panel_gain), 6)
    model["grade_scale"] = round(float(mot.grade_scale), 6)
    model["drive_eff_scale"] = round(float(mot.drive_eff_scale), 6)
    model["regen_eff_scale"] = round(float(mot.drive_eff_scale), 6)
...
\end{lstlisting}

\subsection{L2113 関数 update\_profile}
            \begin{itemize}[leftmargin=1.6em]
              \item 行範囲: L2113--L2154
              \item 所属: module直下
              \item 定義: \path!update_profile(profile_yaml: Path, cfg: dict, map_assets: Dict[str, Path], pv: PvFitResult, batt: BatteryFitResult, mot: MotionFitResult, observed_log_csv: Path, battery_dynamic_fit: Dict[str, Any] | None = None, solar_measurement_calibration: Dict[str, Any] | None = None, route_profile_asset: Path | None = None, package_dir: Path | None = None) -> None!
              \item docstring: 明示されていない。
              \item このブロックが直接呼ぶ主な関数・メソッド: Path, apply\_fit\_to\_cfg, exists, get, is\_absolute, items, list, open, relpath\_from, resolve, safe\_dump, setdefault
              \item 戻り値の要点: 戻り値が無いか、return 文が明示されていない。
              \item 主なローカル代入名: cfg, f, key, package\_dir, package\_path, path, raw
              \item 読み取る主なインスタンス属性: 特になし。
              \item 更新する主なインスタンス属性: 特になし。
              \item 明示的に送出する例外: 明示的raiseはない。
              \item 制御構造の規模: 条件分岐 5、ループ 1、try 0
            \end{itemize}
            \paragraph{この定義を読むためのPython構文}
            \begin{itemize}[leftmargin=1.6em]
            \item `->`以後は戻り値の型注釈であり、実行時の値を自動変換する命令ではない。
\item with文はcontext managerに開始・終了処理を任せ、ファイルなどの資源を確実に解放する。
            \end{itemize}
            \paragraph{上から順の処理}
            \begin{enumerate}[leftmargin=1.8em]
            \item package\_dir に profile\_yaml.parent if package\_dir is None else Path(package\_dir) の結果を代入する。
\item cfg に apply\_fit\_to\_cfg(cfg, package\_dir=package\_dir, map\_assets=map\_assets, pv=pv, batt=batt, mot=mot, battery\_dynamic\_fit=battery\_dynamic\_fit, solar\_measurement\_calibration=solar\_measurement\_calibration, observed\_log\_csv=observed\_log\_csv, sync\_sim\_soc0=False) の結果を代入する。
\item 条件 route\_profile\_asset is not None を判定し、真なら内部処理を行う。
\item   cfg.setdefault('paths', \{\})['route\_profile\_csv'] に relpath\_from(package\_dir, route\_profile\_asset) の結果を代入する。
\item 条件 profile\_yaml.parent.resolve() != package\_dir.resolve() を判定し、真なら内部処理を行う。
\item   list((cfg.get('paths', \{\}) or \{\}).items()) を順に走査し、各要素を (key, raw) に入れて処理する。
\item     条件 raw in (None, '') を判定し、真なら内部処理を行う。
\item       Continue 文を実行する。
\item     path に Path(str(raw)) の結果を代入する。
\item     条件 path.is\_absolute() を判定し、真なら内部処理を行う。
\item       Continue 文を実行する。
\item     package\_path に (package\_dir / path).resolve() の結果を代入する。
\item     条件 package\_path.exists() を判定し、真なら内部処理を行う。
\item       cfg['paths'][key] に relpath\_from(profile\_yaml.parent, package\_path) の結果を代入する。
\item with 文で profile\_yaml.open('w', encoding='utf-8', newline='\textbackslash{}n') を管理しながら処理する。
\item   yaml.safe\_dump(...) を実行する。
            \end{enumerate}
            \paragraph{代表コード断片}
            \begin{lstlisting}[language=Python]
def update_profile(
    profile_yaml: Path,
    cfg: dict,
    map_assets: Dict[str, Path],
    pv: PvFitResult,
    batt: BatteryFitResult,
    mot: MotionFitResult,
    observed_log_csv: Path,
    battery_dynamic_fit: Dict[str, Any] | None = None,
    solar_measurement_calibration: Dict[str, Any] | None = None,
    route_profile_asset: Path | None = None,
    package_dir: Path | None = None,
) -> None:
    package_dir = profile_yaml.parent if package_dir is None else Path(package_dir)
    cfg = apply_fit_to_cfg(
        cfg,
        package_dir=package_dir,
        map_assets=map_assets,
        pv=pv,
        batt=batt,
        mot=mot,
        battery_dynamic_fit=battery_dynamic_fit,
        solar_measurement_calibration=solar_measurement_calibration,
        observed_log_csv=observed_log_csv,
        sync_sim_soc0=False,
    )
    if route_profile_asset is not None:
        cfg.setdefault("paths", {})["route_profile_csv"] = relpath_from(
            package_dir, route_profile_asset
        )
    if profile_yaml.parent.resolve() != package_dir.resolve():
        for key, raw in list((cfg.get("paths", {}) or {}).items()):
            if raw in (None, ""):
                continue
            path = Path(str(raw))
...
\end{lstlisting}

\subsection{L2157 関数 update\_profile\_artifact\_references}
            \begin{itemize}[leftmargin=1.6em]
              \item 行範囲: L2157--L2177
              \item 所属: module直下
              \item 定義: \path!update_profile_artifact_references(profile_yaml: Path, package_dir: Path, *, fit_summary_yaml: Path, terminal_consistency_yaml: Path | None = None) -> None!
              \item docstring: 明示されていない。
              \item このブロックが直接呼ぶ主な関数・メソッド: is\_file, pop, read\_text, relpath\_from, safe\_dump, safe\_load, setdefault, write\_text
              \item 戻り値の要点: 戻り値が無いか、return 文が明示されていない。
              \item 主なローカル代入名: cfg, identification
              \item 読み取る主なインスタンス属性: 特になし。
              \item 更新する主なインスタンス属性: 特になし。
              \item 明示的に送出する例外: 明示的raiseはない。
              \item 制御構造の規模: 条件分岐 1、ループ 0、try 0
            \end{itemize}
            \paragraph{この定義を読むためのPython構文}
            \begin{itemize}[leftmargin=1.6em]
            \item `->`以後は戻り値の型注釈であり、実行時の値を自動変換する命令ではない。
            \end{itemize}
            \paragraph{上から順の処理}
            \begin{enumerate}[leftmargin=1.8em]
            \item cfg に yaml.safe\_load(profile\_yaml.read\_text(encoding='utf-8')) or \{\} の結果を代入する。
\item identification に cfg.setdefault('identification', \{\}) の結果を代入する。
\item identification['fit\_summary\_yaml'] に relpath\_from(profile\_yaml.parent, fit\_summary\_yaml) の結果を代入する。
\item 条件 terminal\_consistency\_yaml is not None and terminal\_consistency\_yaml.is\_file() を判定し、真なら内部処理を行う。
\item   identification['terminal\_consistency\_yaml'] に relpath\_from(profile\_yaml.parent, terminal\_consistency\_yaml) の結果を代入する。
\item   上の条件が偽の場合:
\item   identification.pop(...) を実行する。
\item profile\_yaml.write\_text(...) を実行する。
            \end{enumerate}
            \paragraph{代表コード断片}
            \begin{lstlisting}[language=Python]
def update_profile_artifact_references(
    profile_yaml: Path,
    package_dir: Path,
    *,
    fit_summary_yaml: Path,
    terminal_consistency_yaml: Path | None = None,
) -> None:
    cfg = yaml.safe_load(profile_yaml.read_text(encoding="utf-8")) or {}
    identification = cfg.setdefault("identification", {})
    identification["fit_summary_yaml"] = relpath_from(profile_yaml.parent, fit_summary_yaml)
    if terminal_consistency_yaml is not None and terminal_consistency_yaml.is_file():
        identification["terminal_consistency_yaml"] = relpath_from(
            profile_yaml.parent, terminal_consistency_yaml
        )
    else:
        identification.pop("terminal_consistency_yaml", None)
    profile_yaml.write_text(
        yaml.safe_dump(cfg, sort_keys=False, allow_unicode=True),
        encoding="utf-8",
        newline="\n",
    )
\end{lstlisting}

\subsection{L2180 関数 write\_terminal\_consistency\_from\_anchor}
            \begin{itemize}[leftmargin=1.6em]
              \item 行範囲: L2180--L2315
              \item 所属: module直下
              \item 定義: \path!write_terminal_consistency_from_anchor(output_path: Path, terminal_anchor: Dict[str, Any], *, max_spread: float, validation_metrics: Dict[str, Any] | None = None, replay_soc_error_max: float = 0.02, replay_voltage_error_max_v: float = 0.5, vehicle_soc_error_max: float = 0.02, vehicle_voltage_error_max_v: float = 0.5, end_to_end_soc_error_max: float = 0.03, end_to_end_voltage_error_max_v: float = 1.0) -> Path!
              \item docstring: Always materialize the independent terminal-evidence gate.

Detailed channel reconstruction may later enrich this file, but profile
adoption must never silently drop the gate merely because a separate
reporting command was not run.
              \item このブロックが直接呼ぶ主な関数・メソッド: abs, bool, float, get, isfinite, list, max, min, mkdir, safe\_dump, write\_text
              \item 戻り値の要点: output\_path
              \item 主なローカル代入名: cross\_channel\_gate, end\_to\_end\_gate, end\_to\_end\_soc\_error, end\_to\_end\_voltage\_error, evidence\_spread\_gate, hi, high\_precision\_gate, lo, local\_anchor\_gate, metrics, payload, replay\_gate, replay\_soc\_error, replay\_voltage\_error, replay\_voltage\_observed, replay\_voltage\_predicted, sigma, spread, target, vehicle\_gate
              \item 読み取る主なインスタンス属性: 特になし。
              \item 更新する主なインスタンス属性: 特になし。
              \item 明示的に送出する例外: 明示的raiseはない。
              \item 制御構造の規模: 条件分岐 0、ループ 0、try 0
            \end{itemize}
            \paragraph{この定義を読むためのPython構文}
            \begin{itemize}[leftmargin=1.6em]
            \item `->`以後は戻り値の型注釈であり、実行時の値を自動変換する命令ではない。
            \end{itemize}
            \paragraph{上から順の処理}
            \begin{enumerate}[leftmargin=1.8em]
            \item lo に float(terminal\_anchor.get('soc\_evidence\_min', float('nan'))) の結果を代入する。
\item hi に float(terminal\_anchor.get('soc\_evidence\_max', float('nan'))) の結果を代入する。
\item target に float(terminal\_anchor.get('soc\_target', float('nan'))) の結果を代入する。
\item sigma に float(terminal\_anchor.get('soc\_sigma', float('nan'))) の結果を代入する。
\item spread に hi - lo if np.isfinite(lo) and np.isfinite(hi) else float('nan') の結果を代入する。
\item metrics に validation\_metrics or \{\} の結果を代入する。
\item replay\_soc\_error に abs(float(metrics.get('battery\_conditional\_retire\_anchor\_soc\_error', float('nan')))) の結果を代入する。
\item replay\_voltage\_observed に float(metrics.get('battery\_conditional\_retire\_anchor\_voltage\_obs\_v', float('nan'))) の結果を代入する。
\item replay\_voltage\_predicted に float(metrics.get('battery\_conditional\_retire\_anchor\_voltage\_pred\_v', float('nan'))) の結果を代入する。
\item replay\_voltage\_error に abs(replay\_voltage\_predicted - replay\_voltage\_observed) の結果を代入する。
\item vehicle\_soc\_error に abs(float(metrics.get('retire\_anchor\_soc\_error', float('nan')))) の結果を代入する。
\item vehicle\_voltage\_error に abs(float(metrics.get('retire\_anchor\_voltage\_pred\_v', float('nan'))) - float(metrics.get('retire\_anchor\_voltage\_obs\_v', float('nan')))) の結果を代入する。
\item end\_to\_end\_soc\_error に abs(float(metrics.get('end\_to\_end\_retire\_anchor\_soc\_error', float('nan')))) の結果を代入する。
\item end\_to\_end\_voltage\_error に abs(float(metrics.get('end\_to\_end\_retire\_anchor\_voltage\_pred\_v', float('nan'))) - float(metrics.get('end\_to\_end\_retire\_anchor\_voltage\_obs\_v', float('nan')))) の結果を代入する。
\item evidence\_spread\_gate に bool(np.isfinite(spread) and spread <= float(max\_spread)) の結果を代入する。
\item local\_anchor\_gate に bool(terminal\_anchor.get('quality\_gate\_pass', False)) の結果を代入する。
\item cross\_channel\_gate に bool(terminal\_anchor.get('weak\_channel\_cross\_consistency\_gate\_pass', False)) の結果を代入する。
\item replay\_gate に bool(np.isfinite(replay\_soc\_error) and replay\_soc\_error <= float(replay\_soc\_error\_max) and np.isfinite(replay\_voltage\_error) and (replay\_voltage\_error <= float(replay\_voltage\_error\_max\_v))) の結果を代入する。
\item vehicle\_gate に bool(np.isfinite(vehicle\_soc\_error) and vehicle\_soc\_error <= float(vehicle\_soc\_error\_max) and np.isfinite(vehicle\_voltage\_error) and (vehicle\_voltage\_error <= float(vehicle\_voltage\_error\_max\_v))) の結果を代入する。
\item end\_to\_end\_gate に bool(np.isfinite(end\_to\_end\_soc\_error) and end\_to\_end\_soc\_error <= float(end\_to\_end\_soc\_error\_max) and np.isfinite(end\_to\_end\_voltage\_error) and (end\_to\_end\_voltage\_error <= float(end\_to\_end\_voltage\_error\_max\_v))) の結果を代入する。
\item high\_precision\_gate に bool(evidence\_spread\_gate and local\_anchor\_gate and cross\_channel\_gate and replay\_gate and vehicle\_gate and end\_to\_end\_gate) の結果を代入する。
\item payload に \{'source': 'terminal\_anchor evidence envelope', 'terminal\_distance\_km': float(terminal\_anchor.get('s\_km', float('nan'))), 'evidence\_interval\_min': lo, 'evidence\_interval\_max': hi, 'unweighted\_central\_estimate': target, 'spread\_percentage\_points': 100.0 * spread if np.isfinite(spread) else float('nan'), 'high\_precision\_gate\_pass': high\_precision\_gate, 'high\_precision\_checks': \{'terminal\_evidence\_spread': evidence\_spread\_gate, 'local\_terminal\_anchor\_quality': local\_anchor\_gate, 'independent\_cross\_channel\_consistency': cross\_channel\_gate, 'conditional\_replay\_terminal\_soc': bool(np.isfinite(replay\_soc\_error) and replay\_soc\_error <= float(replay\_soc\_error\_max)), 'conditional\_replay\_terminal\_voltage': bool(np.isfinite(replay\_voltage\_error) and replay\_voltage\_error <= float(replay\_voltage\_error\_max\_v)), 'vehicle\_replay\_terminal': vehicle\_gate, 'end\_to\_end\_replay\_terminal': end\_to\_end\_gate\}, 'conditional\_replay\_terminal\_soc\_error': replay\_soc\_error, 'conditional\_replay\_terminal\_soc\_error\_max': float(replay\_soc\_error\_max), 'conditional\_replay\_terminal\_voltage\_error\_v': replay\_voltage\_error, 'conditional\_replay\_terminal\_voltage\_error\_max\_v': float(replay\_voltage\_error\_max\_v), 'vehicle\_replay\_terminal\_soc\_error': vehicle\_soc\_error, 'vehicle\_replay\_terminal\_soc\_error\_max': float(vehicle\_soc\_error\_max), 'vehicle\_replay\_terminal\_voltage\_error\_v': vehicle\_voltage\_error, 'vehicle\_replay\_terminal\_voltage\_error\_max\_v': float(vehicle\_voltage\_error\_max\_v), 'end\_to\_end\_replay\_terminal\_soc\_error': end\_to\_end\_soc\_error, 'end\_to\_end\_replay\_terminal\_soc\_error\_max': float(end\_to\_end\_soc\_error\_max), 'end\_to\_end\_replay\_terminal\_voltage\_error\_v': end\_to\_end\_voltage\_error, 'end\_to\_end\_replay\_terminal\_voltage\_error\_max\_v': float(end\_to\_end\_voltage\_error\_max\_v), 'random\_effects\_soc': target, 'random\_effects\_standard\_error': sigma, 'random\_effects\_ci95\_min': max(0.0, target - 1.96 * sigma) if np.isfinite(target) and np.isfinite(sigma) else float('nan'), 'random\_effects\_ci95\_max': min(1.0, target + 1.96 * sigma) if np.isfinite(target) and np.isfinite(sigma) else float('nan'), 'method': terminal\_anchor.get('method', ''), 'source\_documents': list(terminal\_anchor.get('source\_documents', []) or []), 'fusion\_caution': 'The center summarizes the declared evidence envelope. A narrow conditional pulse interval cannot override cross-channel disagreement or replay mismatch.', 'interpretation': 'Independent evidence and the battery-only, vehicle, and end-to-end replays satisfy every configured terminal limit.' if high\_precision\_gate else 'One or more independent-evidence or replay checks failed; do not claim high-precision terminal SoC.'\} の結果を代入する。
\item output\_path.parent.mkdir(...) を実行する。
\item output\_path.write\_text(...) を実行する。
\item output\_path を返す。
            \end{enumerate}
            \paragraph{代表コード断片}
            \begin{lstlisting}[language=Python]
def write_terminal_consistency_from_anchor(
    output_path: Path,
    terminal_anchor: Dict[str, Any],
    *,
    max_spread: float,
    validation_metrics: Dict[str, Any] | None = None,
    replay_soc_error_max: float = 0.02,
    replay_voltage_error_max_v: float = 0.5,
    vehicle_soc_error_max: float = 0.02,
    vehicle_voltage_error_max_v: float = 0.5,
    end_to_end_soc_error_max: float = 0.03,
    end_to_end_voltage_error_max_v: float = 1.0,
) -> Path:
    """Always materialize the independent terminal-evidence gate.

    Detailed channel reconstruction may later enrich this file, but profile
    adoption must never silently drop the gate merely because a separate
    reporting command was not run.
    """
    lo = float(terminal_anchor.get("soc_evidence_min", float("nan")))
    hi = float(terminal_anchor.get("soc_evidence_max", float("nan")))
    target = float(terminal_anchor.get("soc_target", float("nan")))
    sigma = float(terminal_anchor.get("soc_sigma", float("nan")))
    spread = hi - lo if np.isfinite(lo) and np.isfinite(hi) else float("nan")
    metrics = validation_metrics or {}
    replay_soc_error = abs(
        float(metrics.get("battery_conditional_retire_anchor_soc_error", float("nan")))
    )
    replay_voltage_observed = float(
        metrics.get("battery_conditional_retire_anchor_voltage_obs_v", float("nan"))
    )
    replay_voltage_predicted = float(
        metrics.get("battery_conditional_retire_anchor_voltage_pred_v", float("nan"))
    )
    replay_voltage_error = abs(replay_voltage_predicted - replay_voltage_observed)
...
\end{lstlisting}

\subsection{L2318 関数 sync\_canonical\_fullsim\_profile}
            \begin{itemize}[leftmargin=1.6em]
              \item 行範囲: L2318--L2350
              \item 所属: module直下
              \item 定義: \path!sync_canonical_fullsim_profile(package_dir: Path, *, map_assets: Dict[str, Path], pv: PvFitResult, batt: BatteryFitResult, mot: MotionFitResult, observed_log_csv: Path, battery_dynamic_fit: Dict[str, Any] | None = None, solar_measurement_calibration: Dict[str, Any] | None = None) -> Path | None!
              \item docstring: 明示されていない。
              \item このブロックが直接呼ぶ主な関数・メソッド: ValueError, apply\_fit\_to\_cfg, exists, isinstance, open, safe\_dump, safe\_load
              \item 戻り値の要点: fullsim\_yaml / None
              \item 主なローカル代入名: cfg, f, fullsim\_yaml
              \item 読み取る主なインスタンス属性: 特になし。
              \item 更新する主なインスタンス属性: 特になし。
              \item 明示的に送出する例外: ValueError(f'fullsim profile must be a mapping: \{fullsim\_yaml\}')
              \item 制御構造の規模: 条件分岐 2、ループ 0、try 0
            \end{itemize}
            \paragraph{この定義を読むためのPython構文}
            \begin{itemize}[leftmargin=1.6em]
            \item `->`以後は戻り値の型注釈であり、実行時の値を自動変換する命令ではない。
\item with文はcontext managerに開始・終了処理を任せ、ファイルなどの資源を確実に解放する。
            \end{itemize}
            \paragraph{上から順の処理}
            \begin{enumerate}[leftmargin=1.8em]
            \item fullsim\_yaml に package\_dir / 'profile\_fullsim\_selflearned.yaml' の結果を代入する。
\item 条件 not fullsim\_yaml.exists() を判定し、真なら内部処理を行う。
\item   None を返す。
\item with 文で fullsim\_yaml.open('r', encoding='utf-8') を管理しながら処理する。
\item   cfg に yaml.safe\_load(f) or \{\} の結果を代入する。
\item 条件 not isinstance(cfg, dict) を判定し、真なら内部処理を行う。
\item   ValueError(f'fullsim profile must be a mapping: \{fullsim\_yaml\}') を送出する。
\item cfg に apply\_fit\_to\_cfg(cfg, package\_dir=package\_dir, map\_assets=map\_assets, pv=pv, batt=batt, mot=mot, battery\_dynamic\_fit=battery\_dynamic\_fit, solar\_measurement\_calibration=solar\_measurement\_calibration, observed\_log\_csv=observed\_log\_csv, sync\_sim\_soc0=False) の結果を代入する。
\item with 文で fullsim\_yaml.open('w', encoding='utf-8', newline='\textbackslash{}n') を管理しながら処理する。
\item   yaml.safe\_dump(...) を実行する。
\item fullsim\_yaml を返す。
            \end{enumerate}
            \paragraph{代表コード断片}
            \begin{lstlisting}[language=Python]
def sync_canonical_fullsim_profile(
    package_dir: Path,
    *,
    map_assets: Dict[str, Path],
    pv: PvFitResult,
    batt: BatteryFitResult,
    mot: MotionFitResult,
    observed_log_csv: Path,
    battery_dynamic_fit: Dict[str, Any] | None = None,
    solar_measurement_calibration: Dict[str, Any] | None = None,
) -> Path | None:
    fullsim_yaml = package_dir / "profile_fullsim_selflearned.yaml"
    if not fullsim_yaml.exists():
        return None
    with fullsim_yaml.open("r", encoding="utf-8") as f:
        cfg = yaml.safe_load(f) or {}
    if not isinstance(cfg, dict):
        raise ValueError(f"fullsim profile must be a mapping: {fullsim_yaml}")
    cfg = apply_fit_to_cfg(
        cfg,
        package_dir=package_dir,
        map_assets=map_assets,
        pv=pv,
        batt=batt,
        mot=mot,
        battery_dynamic_fit=battery_dynamic_fit,
        solar_measurement_calibration=solar_measurement_calibration,
        observed_log_csv=observed_log_csv,
        sync_sim_soc0=False,
    )
    with fullsim_yaml.open("w", encoding="utf-8", newline="\n") as f:
        yaml.safe_dump(cfg, f, sort_keys=False, allow_unicode=True)
    return fullsim_yaml
\end{lstlisting}

\subsection{L2353 関数 write\_generic\_summary}
            \begin{itemize}[leftmargin=1.6em]
              \item 行範囲: L2353--L2419
              \item 所属: module直下
              \item 定義: \path!write_generic_summary(package_dir: Path, manifest_path: Path, profile_yaml: Path, map_assets: Dict[str, Path], pv: PvFitResult, batt: BatteryFitResult, mot: MotionFitResult, metrics: Dict[str, float], terminal_anchor: Dict[str, float], stage_anchors: list[dict], map_shape_fit: Dict[str, object], post_refine: PostRefineResult, day_metrics: list[dict], battery_dynamic_fit: Dict[str, Any] | None = None, fit_plan: dict | None = None, manifest_context: dict | None = None, output_dir: Path | None = None, replay_csv: Path | None = None, battery_conditioned_replay_csv: Path | None = None, end_to_end_replay_csv: Path | None = None) -> Path!
              \item docstring: 明示されていない。
              \item このブロックが直接呼ぶ主な関数・メソッド: dict, ensure\_dir, get, items, list, open, relpath, relpath\_from, replace, safe\_dump
              \item 戻り値の要点: out\_path
              \item 主なローカル代入名: f, key, out\_path, output\_dir, path, payload
              \item 読み取る主なインスタンス属性: 特になし。
              \item 更新する主なインスタンス属性: 特になし。
              \item 明示的に送出する例外: 明示的raiseはない。
              \item 制御構造の規模: 条件分岐 0、ループ 0、try 0
            \end{itemize}
            \paragraph{この定義を読むためのPython構文}
            \begin{itemize}[leftmargin=1.6em]
            \item `->`以後は戻り値の型注釈であり、実行時の値を自動変換する命令ではない。
\item 内包表記は、反復・条件・値生成を一つの式で記述してcollectionを作る。
\item with文はcontext managerに開始・終了処理を任せ、ファイルなどの資源を確実に解放する。
            \end{itemize}
            \paragraph{上から順の処理}
            \begin{enumerate}[leftmargin=1.8em]
            \item output\_dir に output\_dir or package\_dir / 'outputs' / 'identification' の結果を代入する。
\item out\_path に output\_dir / f'\{package\_dir.name\}\_generic\_fit\_summary.yaml' の結果を代入する。
\item ensure\_dir(...) を実行する。
\item payload に \{'builder': 'generic\_replay\_mle', 'manifest\_yaml': os.path.relpath(manifest\_path, package\_dir).replace('\textbackslash{}\textbackslash{}', '/'), 'profile\_yaml': os.path.relpath(profile\_yaml, package\_dir).replace('\textbackslash{}\textbackslash{}', '/'), 'active\_maps': \{key: os.path.relpath(path, package\_dir).replace('\textbackslash{}\textbackslash{}', '/') for key, path in map\_assets.items()\}, 'pv\_fit': pv.\_\_dict\_\_, 'battery\_fit': batt.\_\_dict\_\_, 'battery\_dynamic\_fit': dict(battery\_dynamic\_fit or \{\}), 'motion\_fit': mot.\_\_dict\_\_, 'validation\_metrics': metrics, 'validation\_protocol': \{'vehicle\_conditional\_replay\_csv': relpath\_from(package\_dir, replay\_csv), 'battery\_conditioned\_replay\_csv': relpath\_from(package\_dir, battery\_conditioned\_replay\_csv), 'end\_to\_end\_replay\_csv': relpath\_from(package\_dir, end\_to\_end\_replay\_csv), 'vehicle\_fit\_solar\_source': 'measured\_when\_available', 'battery\_conditioned\_source': 'observed\_pack\_power\_and\_current', 'end\_to\_end\_solar\_source': 'independent\_GHI\_archive\_and\_moving\_PV\_model', 'restart\_soc\_anchor': 'median\_of\_valid\_stationary\_window'\}, 'terminal\_anchor': terminal\_anchor, 'stage\_anchors': stage\_anchors, 'day\_metrics': day\_metrics, 'map\_shape\_fit': map\_shape\_fit, 'post\_refine': post\_refine.\_\_dict\_\_, 'fit\_plan': dict(fit\_plan or \{\}), 'evidence\_bundle': \{'actual\_event\_yaml': relpath\_from(package\_dir, (manifest\_context or \{\}).get('actual\_event\_path')), 'counterfactual\_event\_yaml': relpath\_from(package\_dir, (manifest\_context or \{\}).get('counterfactual\_event\_path')), 'terminal\_anchor\_yaml': relpath\_from(package\_dir, (manifest\_context or \{\}).get('terminal\_anchor\_path')), 'grounded\_map\_summary\_yaml': relpath\_from(package\_dir, (manifest\_context or \{\}).get('grounded\_summary\_path')), 'source\_inventory\_json': relpath\_from(package\_dir, (manifest\_context or \{\}).get('source\_inventory\_path')), 'notes\_markdown': relpath\_from(package\_dir, (manifest\_context or \{\}).get('notes\_markdown\_path')), 'explicit\_grounded\_assets': \{key: relpath\_from(package\_dir, path) for key, path in ((manifest\_context or \{\}).get('explicit\_grounded\_assets') or \{\}).items()\}, 'external\_documents': list((manifest\_context or \{\}).get('external\_documents', []))\}\} の結果を代入する。
\item with 文で out\_path.open('w', encoding='utf-8', newline='\textbackslash{}n') を管理しながら処理する。
\item   yaml.safe\_dump(...) を実行する。
\item out\_path を返す。
            \end{enumerate}
            \paragraph{代表コード断片}
            \begin{lstlisting}[language=Python]
def write_generic_summary(
    package_dir: Path,
    manifest_path: Path,
    profile_yaml: Path,
    map_assets: Dict[str, Path],
    pv: PvFitResult,
    batt: BatteryFitResult,
    mot: MotionFitResult,
    metrics: Dict[str, float],
    terminal_anchor: Dict[str, float],
    stage_anchors: list[dict],
    map_shape_fit: Dict[str, object],
    post_refine: PostRefineResult,
    day_metrics: list[dict],
    battery_dynamic_fit: Dict[str, Any] | None = None,
    fit_plan: dict | None = None,
    manifest_context: dict | None = None,
    output_dir: Path | None = None,
    replay_csv: Path | None = None,
    battery_conditioned_replay_csv: Path | None = None,
    end_to_end_replay_csv: Path | None = None,
) -> Path:
    output_dir = output_dir or package_dir / "outputs" / "identification"
    out_path = output_dir / f"{package_dir.name}_generic_fit_summary.yaml"
    ensure_dir(out_path.parent)
    payload = {
        "builder": "generic_replay_mle",
        "manifest_yaml": os.path.relpath(manifest_path, package_dir).replace("\\", "/"),
        "profile_yaml": os.path.relpath(profile_yaml, package_dir).replace("\\", "/"),
        "active_maps": {key: os.path.relpath(path, package_dir).replace("\\", "/") for key, path in map_assets.items()},
        "pv_fit": pv.__dict__,
        "battery_fit": batt.__dict__,
        "battery_dynamic_fit": dict(battery_dynamic_fit or {}),
        "motion_fit": mot.__dict__,
        "validation_metrics": metrics,
...
\end{lstlisting}

\subsection{L2422 関数 write\_generic\_report}
            \begin{itemize}[leftmargin=1.6em]
              \item 行範囲: L2422--L3072
              \item 所属: module直下
              \item 定義: \path!write_generic_report(package_dir: Path, profile_yaml: Path, manifest_path: Path, summary_yaml: Path, observed_log_csv: Path, pv: PvFitResult, batt: BatteryFitResult, mot: MotionFitResult, metrics: Dict[str, float], post_refine: PostRefineResult, map_assets: Dict[str, Path], *, terminal_anchor: Dict[str, float] | None = None, stage_anchors: list[dict] | None = None, day_metrics: list[dict] | None = None, battery_dynamic_fit: Dict[str, Any] | None = None, fit_plan: dict | None = None, grounded_map_summary: Dict[str, object] | None = None, manifest_context: dict | None = None, terminal_consistency: Dict[str, Any] | None = None, report_dir: Path | None = None, current_maps_path: Path | None = None) -> Tuple[Path, Path]!
              \item docstring: 明示されていない。
              \item このブロックが直接呼ぶ主な関数・メソッド: bool, compile\_tex, dedent, ensure\_dir, enumerate, float, get, int, items, join, len, read\_text
              \item 戻り値の要点: (md\_path, pdf\_path)
              \item 主なローカル代入名: anchor, battery\_dynamic\_fit, current\_maps\_path, current\_maps\_rel\_package, current\_maps\_rel\_report, day\_metrics, evidence\_rows, explicit\_grounded\_assets, external\_documents, fit\_plan, grade\_observation\_fit, grounded\_yaml, idx, item, key, manifest\_context, md, md\_path, pack\_voltage\_limit\_v, path
              \item 読み取る主なインスタンス属性: 特になし。
              \item 更新する主なインスタンス属性: 特になし。
              \item 明示的に送出する例外: 明示的raiseはない。
              \item 制御構造の規模: 条件分岐 0、ループ 0、try 0
            \end{itemize}
            \paragraph{この定義を読むためのPython構文}
            \begin{itemize}[leftmargin=1.6em]
            \item `->`以後は戻り値の型注釈であり、実行時の値を自動変換する命令ではない。
\item 内包表記は、反復・条件・値生成を一つの式で記述してcollectionを作る。
            \end{itemize}
            \paragraph{上から順の処理}
            \begin{enumerate}[leftmargin=1.8em]
            \item report\_dir に report\_dir or package\_dir / 'outputs' / 'reports' の結果を代入する。
\item ensure\_dir(...) を実行する。
\item md\_path に report\_dir / f'\{package\_dir.name\}\_generic\_identification\_report.md' の結果を代入する。
\item tex\_path に report\_dir / f'\{package\_dir.name\}\_generic\_identification\_report.tex' の結果を代入する。
\item pdf\_path に tex\_path.with\_suffix('.pdf') の結果を代入する。
\item rel\_maps に \{key: os.path.relpath(path, package\_dir).replace('\textbackslash{}\textbackslash{}', '/') for key, path in map\_assets.items()\} の結果を代入する。
\item terminal\_anchor に terminal\_anchor or \{\} の結果を代入する。
\item stage\_anchors に stage\_anchors or [] の結果を代入する。
\item day\_metrics に day\_metrics or [] の結果を代入する。
\item battery\_dynamic\_fit に battery\_dynamic\_fit or \{\} の結果を代入する。
\item fit\_plan に fit\_plan or \{\} の結果を代入する。
\item solar\_calibration に fit\_plan.get('solar\_measurement\_calibration', \{\}) or \{\} の結果を代入する。
\item grade\_observation\_fit に fit\_plan.get('grade\_observation\_fit', \{\}) or \{\} の結果を代入する。
\item solar\_gain\_to\_pack に float(solar\_calibration.get('gain\_to\_pack', 1.0)) の結果を代入する。
\item solar\_calibration\_samples に int(solar\_calibration.get('sample\_count', 0) or 0) の結果を代入する。
\item solar\_calibration\_accepted に bool(solar\_calibration.get('accepted', False)) の結果を代入する。
\item solar\_calibration\_intercept\_w に float(solar\_calibration.get('free\_intercept\_w', float('nan'))) の結果を代入する。
\item solar\_calibration\_intercept\_error\_w に float(solar\_calibration.get('free\_intercept\_error\_w', float('nan'))) の結果を代入する。
\item solar\_calibration\_daily\_std に float(solar\_calibration.get('daily\_gain\_std', float('nan'))) の結果を代入する。
\item profile\_cfg に yaml.safe\_load(profile\_yaml.read\_text(encoding='utf-8-sig')) or \{\} の結果を代入する。
\item profile\_model に profile\_cfg.get('model', \{\}) or \{\} の結果を代入する。
\item pack\_voltage\_limit\_v に float(profile\_model.get('V\_max', BATTERY\_PACK\_MAX\_CHARGE\_V)) の結果を代入する。
\item grounded\_yaml に str((grounded\_map\_summary or \{\}).get('summary\_yaml', '') or '') の結果を代入する。
\item manifest\_context に manifest\_context or \{\} の結果を代入する。
\item terminal\_consistency に terminal\_consistency or \{\} の結果を代入する。
\item terminal\_checks に terminal\_consistency.get('high\_precision\_checks', \{\}) or \{\} の結果を代入する。
\item terminal\_gate\_pass に bool(terminal\_consistency.get('high\_precision\_gate\_pass', False)) の結果を代入する。
\item terminal\_evidence\_min に float(terminal\_consistency.get('evidence\_interval\_min', float('nan'))) の結果を代入する。
\item terminal\_evidence\_max に float(terminal\_consistency.get('evidence\_interval\_max', float('nan'))) の結果を代入する。
\item terminal\_evidence\_spread\_pp に float(terminal\_consistency.get('spread\_percentage\_points', float('nan'))) の結果を代入する。
\item terminal\_interpretation に str(terminal\_consistency.get('interpretation', '') or '') の結果を代入する。
\item evidence\_rows に [('actual\_event\_yaml', relpath\_from(package\_dir, manifest\_context.get('actual\_event\_path'))), ('counterfactual\_event\_yaml', relpath\_from(package\_dir, manifest\_context.get('counterfactual\_event\_path'))), ('terminal\_anchor\_yaml', relpath\_from(package\_dir, manifest\_context.get('terminal\_anchor\_path'))), ('grounded\_map\_summary\_yaml', relpath\_from(package\_dir, manifest\_context.get('grounded\_summary\_path')) or grounded\_yaml), ('source\_inventory\_json', relpath\_from(package\_dir, manifest\_context.get('source\_inventory\_path'))), ('notes\_markdown', relpath\_from(package\_dir, manifest\_context.get('notes\_markdown\_path')))] の結果を代入する。
\item explicit\_grounded\_assets に \{key: relpath\_from(package\_dir, path) for key, path in (manifest\_context.get('explicit\_grounded\_assets') or \{\}).items()\} の結果を代入する。
\item external\_documents に [str(item) for item in manifest\_context.get('external\_documents', []) if str(item).strip()] の結果を代入する。
\item current\_maps\_path に current\_maps\_path or report\_dir / 'current\_maps\_and\_coefficients.md' の結果を代入する。
\item profile\_rel に os.path.relpath(profile\_yaml, package\_dir).replace('\textbackslash{}\textbackslash{}', '/') の結果を代入する。
\item current\_maps\_rel\_package に os.path.relpath(current\_maps\_path, package\_dir).replace('\textbackslash{}\textbackslash{}', '/') の結果を代入する。
\item current\_maps\_rel\_report に os.path.relpath(current\_maps\_path, report\_dir).replace('\textbackslash{}\textbackslash{}', '/') の結果を代入する。
\item md に f"\# \{package\_dir.name\} generic identification report\textbackslash{}n\textbackslash{}n\#\# Inputs\textbackslash{}n\textbackslash{}n- profile: `\{profile\_rel\}`\textbackslash{}n- manifest: `\{os.path.relpath(manifest\_path, package\_dir).replace('\textbackslash{}\textbackslash{}', '/')\}`\textbackslash{}n- normalized replay log: `\{os.path.relpath(observed\_log\_csv, package\_dir).replace('\textbackslash{}\textbackslash{}', '/')\}`\textbackslash{}n\textbackslash{}n\#\# Adopted coefficients\textbackslash{}n\textbackslash{}n- panel\_gain: `\{pv.panel\_gain:.6f\}`\textbackslash{}n- tcell\_gain\_c\_per\_wm2: `\{pv.tcell\_gain\_c\_per\_wm2:.6f\}`\textbackslash{}n- pv\_irradiance\_source: `\{pv.irradiance\_source\}`\textbackslash{}n- pv\_operating\_state: `\{pv.operating\_state\}`\textbackslash{}n- pv\_sample\_count: `\{pv.sample\_count\}`\textbackslash{}n- stop\_tilt\_fraction: `\{pv.stop\_tilt\_fraction:.6f\}`\textbackslash{}n- stop\_solar\_rmse\_w: `\{pv.stop\_solar\_rmse\_w:.3f\}`\textbackslash{}n- solar\_measurement\_gain\_to\_pack: `\{solar\_gain\_to\_pack:.8f\}`\textbackslash{}n- solar\_measurement\_calibration\_accepted: `\{solar\_calibration\_accepted\}`\textbackslash{}n- solar\_measurement\_calibration\_samples: `\{solar\_calibration\_samples\}`\textbackslash{}n- solar\_measurement\_free\_intercept\_w: `\{solar\_calibration\_intercept\_w:.6f\}`\textbackslash{}n- solar\_measurement\_intercept\_error\_w: `\{solar\_calibration\_intercept\_error\_w:.6f\}`\textbackslash{}n- solar\_measurement\_daily\_gain\_std: `\{solar\_calibration\_daily\_std:.8f\}`\textbackslash{}n- soc0: `\{batt.soc0:.6f\}`\textbackslash{}n- E\_nom\_Wh: `\{batt.e\_nom\_wh:.3f\}`\textbackslash{}n- Q\_nom\_Ah: `\{batt.e\_nom\_wh / BATTERY\_NOMINAL\_VOLTAGE\_V:.6f\}`\textbackslash{}n- rint\_scale: `\{batt.rint\_scale:.6f\}`\textbackslash{}n- r\_line\_ohm: `\{batt.r\_line\_ohm:.6f\}`\textbackslash{}n- r\_polarization\_ohm: `\{float(battery\_dynamic\_fit.get('r\_polarization\_ohm', 0.0)):.6f\}`\textbackslash{}n- polarization\_tau\_sec: `\{float(battery\_dynamic\_fit.get('tau\_sec', 60.0)):.3f\}`\textbackslash{}n- eta\_charge: `\{batt.eta\_charge:.6f\}`\textbackslash{}n- CdA: `\{mot.cda:.6f\}`\textbackslash{}n- Crr: `\{mot.crr:.6f\}`\textbackslash{}n- P\_aux\_w: `\{mot.p\_aux\_w:.3f\}`\textbackslash{}n- grade\_scale: `\{mot.grade\_scale:.6f\}`\textbackslash{}n- drive\_eff\_scale: `\{mot.drive\_eff\_scale:.6f\}`\textbackslash{}n- regen\_utilization: `\{mot.regen\_utilization:.6f\}`\textbackslash{}n- regen\_sample\_count: `\{mot.regen\_sample\_count\}`\textbackslash{}n- regen\_fit\_rmse\_w: `\{mot.regen\_fit\_rmse\_w:.3f\}`\textbackslash{}n- headwind\_gain: `\{mot.headwind\_gain:.6f\}`\textbackslash{}n- V\_max\_v: `\{pack\_voltage\_limit\_v:.3f\}`\textbackslash{}n\textbackslash{}n\#\# Solar measurement calibration\textbackslash{}n\textbackslash{}nThe ZP solar channel is calibrated against the stationary DC-bus balance before\textbackslash{}nvehicle or PV fitting. For samples with zero traction power, the fitted model is\textbackslash{}n`P\_batt,k = P\_aux - g\_solar * P\_solar,raw,k + epsilon\_k`. The adopted gain is\textbackslash{}nthe bounded Huber M-estimator\textbackslash{}n`argmin\_(0.70 <= g\_solar <= 1.05) sum\_k rho\_H(P\_batt,k - P\_aux + g\_solar * P\_solar,raw,k)`.\textbackslash{}nThe known 21 W auxiliary load fixes the intercept; a separate free-intercept fit\textbackslash{}nis retained only as a consistency diagnostic. Daily gain spread checks whether a\textbackslash{}nsingle gain is defensible across the race.\textbackslash{}n\textbackslash{}nThe telemetry contract is deliberately one-way: the sender transmits the raw ZP\textbackslash{}n`solar\_power\_w` value, and the receiving WiFi bridge multiplies it by the active\textbackslash{}nprofile's `solar\_measurement\_gain\_to\_pack` exactly once. Corrected values must not\textbackslash{}nbe sent through that raw field, which prevents double calibration.\textbackslash{}n\textbackslash{}nThe battery terminal-voltage ceiling is a hardware constraint, not a value fitted\textbackslash{}nfrom a loaded discharge trace. For YATA the adopted `\{pack\_voltage\_limit\_v:.3f\} V`\textbackslash{}nmust equal the product limit `25 series cells * 4.35 V/cell = 108.75 V`, and the\textbackslash{}nactive OCV map must not exceed it. The package audit checks both directions.\textbackslash{}n\textbackslash{}n\#\# Validation\textbackslash{}n\textbackslash{}nVehicle coefficients are identified conditionally on measured array power so\textbackslash{}nforecast/PV error cannot be absorbed into CdA, Crr, or drivetrain efficiency.\textbackslash{}nThe unprefixed metrics below are that conditional vehicle replay.  Keys beginning\textbackslash{}nwith `end\_to\_end\_` use independent archive GHI plus the moving-PV model and\textbackslash{}ntherefore include forecast/PV error. `GHI\_effective` is target-derived diagnostic\textbackslash{}ndata and is prohibited from both fitting and end-to-end validation. Both datasets\textbackslash{}nare retained; neither is substituted for the other.\textbackslash{}n\textbackslash{}nKeys beginning with `battery\_conditional\_` reuse observed pack power/current.\textbackslash{}nThey validate the battery submodel only and cannot certify vehicle energy use or\textbackslash{}nthe 2831 km vehicle-model terminal SoC. Promotion therefore requires separate\textbackslash{}nterminal gates for the unprefixed vehicle replay and the end-to-end replay.\textbackslash{}n\textbackslash{}nThe field-analysis source states that the ZP logger start time was not recorded\textbackslash{}nand was reconstructed backward from control-stop times. Therefore 5 s samples\textbackslash{}ndo not provide a traceable pointwise synchronization certificate among power,\textbackslash{}nspeed, and DEM grade. Parameter fitting uses the configured resampling window;\textbackslash{}nthe 10 km and 25 km energy residuals are the primary route-energy checks. A low\textbackslash{}npointwise RMSE must not be manufactured by fitting a time shift to this unknown.\textbackslash{}n\textbackslash{}nThe 5 s RMSE, 120 s mean-residual RMSE, and distance-window energy RMSE answer\textbackslash{}ndifferent questions. The 5 s value includes unresolved channel synchronization\textbackslash{}nand transient noise. The 120 s value tests local mean power after aggregation.\textbackslash{}nThe 10/25 km values test the accumulated energy that drives long-horizon SoC.\textbackslash{}nNone may be relabelled as another metric, and operational promotion requires all\textbackslash{}nconfigured gates rather than choosing only the smallest number.\textbackslash{}n\textbackslash{}nDEM grade is treated as an observation model rather than a vehicle constant.\textbackslash{}nCandidate Savitzky-Golay elevation smoothing lengths and route-distance offsets\textbackslash{}nare selected on race days other than the held-out final day. The adopted route\textbackslash{}nstores the unscaled differentiated elevation; `grade\_scale` is fitted only\textbackslash{}nafter that route is locked. This prevents DEM vertical noise from being hidden\textbackslash{}ninside CdA, Crr, or drivetrain efficiency.\textbackslash{}n\textbackslash{}nRegeneration is separated into conversion efficiency and use:\textbackslash{}n`P\_reg,dc = u\_regen * eta\_reg * eta\_gear * eta\_inv * max(-P\_mech, 0)`.\textbackslash{}nThe bounded scalar `u\_regen` is fitted only where negative mechanical power is\textbackslash{}nobservable and is never folded into the motor efficiency map.\textbackslash{}n\textbackslash{}n- power\_rmse\_clean\_w: `\{metrics.get('power\_rmse\_clean\_w', float('nan')):.3f\}`\textbackslash{}n- power\_rmse\_fit\_window\_w: `\{metrics.get('power\_rmse\_fit\_window\_w', float('nan')):.3f\}`\textbackslash{}n- voltage\_rmse\_clean\_v: `\{metrics.get('voltage\_rmse\_clean\_v', float('nan')):.3f\}`\textbackslash{}n- final\_soc\_pred: `\{metrics.get('final\_soc\_pred', float('nan')):.6f\}`\textbackslash{}n- retire\_anchor\_soc\_obs: `\{metrics.get('retire\_anchor\_soc\_obs', float('nan')):.6f\}`\textbackslash{}n- retire\_anchor\_soc\_pred: `\{metrics.get('retire\_anchor\_soc\_pred', float('nan')):.6f\}`\textbackslash{}n- retire\_anchor\_soc\_error: `\{metrics.get('retire\_anchor\_soc\_error', float('nan')):.6f\}`\textbackslash{}n- battery\_conditional\_power\_rmse\_clean\_w: `\{metrics.get('battery\_conditional\_power\_rmse\_clean\_w', float('nan')):.3f\}`\textbackslash{}n- battery\_conditional\_voltage\_rmse\_clean\_v: `\{metrics.get('battery\_conditional\_voltage\_rmse\_clean\_v', float('nan')):.3f\}`\textbackslash{}n- battery\_conditional\_retire\_anchor\_soc\_error: `\{metrics.get('battery\_conditional\_retire\_anchor\_soc\_error', float('nan')):.6f\}`\textbackslash{}n- end\_to\_end\_power\_rmse\_clean\_w: `\{metrics.get('end\_to\_end\_power\_rmse\_clean\_w', float('nan')):.3f\}`\textbackslash{}n- end\_to\_end\_voltage\_rmse\_clean\_v: `\{metrics.get('end\_to\_end\_voltage\_rmse\_clean\_v', float('nan')):.3f\}`\textbackslash{}n- end\_to\_end\_final\_soc\_pred: `\{metrics.get('end\_to\_end\_final\_soc\_pred', float('nan')):.6f\}`\textbackslash{}n- end\_to\_end\_retire\_anchor\_soc\_error: `\{metrics.get('end\_to\_end\_retire\_anchor\_soc\_error', float('nan')):.6f\}`\textbackslash{}n- end\_to\_end\_moving\_pv\_rmse\_w: `\{metrics.get('end\_to\_end\_moving\_pv\_rmse\_w', float('nan')):.3f\}`\textbackslash{}n- end\_to\_end\_deployed\_stop\_pv\_rmse\_w: `\{metrics.get('end\_to\_end\_deployed\_stop\_pv\_rmse\_w', float('nan')):.3f\}`\textbackslash{}n- pv\_lodo\_moving\_rmse\_w: `\{metrics.get('pv\_lodo\_moving\_rmse\_w', float('nan')):.3f\}`\textbackslash{}n- pv\_lodo\_deployed\_stop\_rmse\_w: `\{metrics.get('pv\_lodo\_deployed\_stop\_rmse\_w', float('nan')):.3f\}`\textbackslash{}n- pv\_lodo\_fold\_count: `\{metrics.get('pv\_lodo\_fold\_count', 0)\}`\textbackslash{}n- power\_residual\_mean\_120s\_rmse\_w: `\{metrics.get('power\_residual\_mean\_120s\_rmse\_w', float('nan')):.3f\}`\textbackslash{}n- energy\_error\_10km\_rmse\_wh: `\{metrics.get('energy\_error\_10km\_rmse\_wh', float('nan')):.3f\}`\textbackslash{}n- energy\_error\_25km\_rmse\_wh: `\{metrics.get('energy\_error\_25km\_rmse\_wh', float('nan')):.3f\}`\textbackslash{}n\textbackslash{}n\#\# Terminal-SoC evidence and certification\textbackslash{}n\textbackslash{}n- evidence\_interval: `[\{terminal\_evidence\_min:.6f\}, \{terminal\_evidence\_max:.6f\}]`\textbackslash{}n- evidence\_spread\_percentage\_points: `\{terminal\_evidence\_spread\_pp:.3f\}`\textbackslash{}n- high\_precision\_gate\_pass: `\{terminal\_gate\_pass\}`\textbackslash{}n" + ('\textbackslash{}n'.join((f'- high\_precision\_check\_\{key\}: `\{bool(value)\}`' for key, value in terminal\_checks.items())) if terminal\_checks else '- high\_precision\_checks: `(not evaluated)`') + f"\textbackslash{}n- interpretation: `\{terminal\_interpretation\}`\textbackslash{}n\textbackslash{}nThe central terminal-SoC estimate is not a direct coulomb-counter measurement.\textbackslash{}nIf the evidence spread or any replay check fails, the value is retained as an\textbackslash{}nengineering anchor only; this report does not certify a high-precision actual\textbackslash{}n2831 km SoC and the candidate must not replace the operational profile.\textbackslash{}n\textbackslash{}n\#\# Method references\textbackslash{}n\textbackslash{}n- Open-Meteo Historical Weather API: https://open-meteo.com/en/docs/historical-weather-api\textbackslash{}n- Sandia PVPMC plane-of-array irradiance guidance: https://pvpmc.sandia.gov/modeling-guide/1-weather-design-inputs/plane-of-array-poa-irradiance/\textbackslash{}n- pvlib total-irradiance transposition API: https://pvlib-python.readthedocs.io/en/latest/reference/generated/pvlib.irradiance.get\_total\_irradiance.html\textbackslash{}n- Byrd et al. (1995), L-BFGS-B: https://doi.org/10.1137/0916069\textbackslash{}n- Liu et al. (2014), GPS road-grade quantification: https://doi.org/10.1016/j.atmosenv.2013.12.025\textbackslash{}n- Wang et al. (2018), synchronized link-level EV energy data: https://escholarship.org/uc/item/2bp8x04q\textbackslash{}n\textbackslash{}n\#\# Terminal anchor\textbackslash{}n\textbackslash{}n- anchor\_s\_km: `\{terminal\_anchor.get('s\_km', float('nan'))\}`\textbackslash{}n- anchor\_time\_utc: `\{terminal\_anchor.get('time\_utc', '')\}`\textbackslash{}n- anchor\_voltage\_v: `\{terminal\_anchor.get('voltage\_v', float('nan'))\}`\textbackslash{}n- anchor\_current\_a: `\{terminal\_anchor.get('current\_a', float('nan'))\}`\textbackslash{}n- anchor\_temp\_c: `\{terminal\_anchor.get('temp\_c', float('nan'))\}`\textbackslash{}n- anchor\_soc\_target: `\{terminal\_anchor.get('soc\_target', float('nan'))\}`\textbackslash{}n\textbackslash{}n\#\# Fit plan\textbackslash{}n\textbackslash{}n- quality: `\{fit\_plan.get('quality', '')\}`\textbackslash{}n- battery\_restart\_count: `\{fit\_plan.get('battery\_restart\_count', '')\}`\textbackslash{}n- battery\_maxiter: `\{fit\_plan.get('battery\_maxiter', '')\}`\textbackslash{}n- motion\_restart\_count: `\{fit\_plan.get('motion\_restart\_count', '')\}`\textbackslash{}n- motion\_maxiter: `\{fit\_plan.get('motion\_maxiter', '')\}`\textbackslash{}n- joint\_restart\_count: `\{fit\_plan.get('joint\_restart\_count', '')\}`\textbackslash{}n- joint\_random\_start\_count: `\{fit\_plan.get('joint\_random\_start\_count', '')\}`\textbackslash{}n- joint\_local\_topk: `\{fit\_plan.get('joint\_local\_topk', '')\}`\textbackslash{}n- joint\_maxiter: `\{fit\_plan.get('joint\_maxiter', '')\}`\textbackslash{}n- fit\_stride: `\{fit\_plan.get('fit\_stride', '')\}`\textbackslash{}n- allow\_map\_shape\_fit: `\{fit\_plan.get('allow\_map\_shape\_fit', '')\}`\textbackslash{}n- post\_refine\_enabled: `\{fit\_plan.get('post\_refine\_enabled', '')\}`\textbackslash{}n- terminal\_anchor\_role: `\{fit\_plan.get('terminal\_anchor\_role', '')\}`\textbackslash{}n- panel\_deployment\_stopped\_speed\_kmh: `\{fit\_plan.get('panel\_deployment\_stopped\_speed\_kmh', '')\}`\textbackslash{}n- panel\_deployment\_min\_dwell\_sec: `\{fit\_plan.get('panel\_deployment\_min\_dwell\_sec', '')\}`\textbackslash{}n- panel\_deployment\_max\_sample\_gap\_sec: `\{fit\_plan.get('panel\_deployment\_max\_sample\_gap\_sec', '')\}`\textbackslash{}n- grade\_observation\_adopted: `\{bool(grade\_observation\_fit.get('adopted', False))\}`\textbackslash{}n- grade\_smoothing\_window\_km: `\{grade\_observation\_fit.get('selected\_smoothing\_window\_km', '')\}`\textbackslash{}n- grade\_distance\_offset\_km: `\{grade\_observation\_fit.get('selected\_distance\_offset\_km', '')\}`\textbackslash{}n- grade\_training\_rmse\_w: `\{grade\_observation\_fit.get('baseline\_training\_rmse\_w', float('nan'))\}` -> `\{grade\_observation\_fit.get('selected\_training\_rmse\_w', float('nan'))\}`\textbackslash{}n- grade\_holdout\_rmse\_w: `\{grade\_observation\_fit.get('baseline\_validation\_rmse\_w', float('nan'))\}` -> `\{grade\_observation\_fit.get('selected\_validation\_rmse\_w', float('nan'))\}`\textbackslash{}n\textbackslash{}n\#\# Stage anchors\textbackslash{}n\textbackslash{}n- stage\_anchor\_count: `\{len(stage\_anchors)\}`\textbackslash{}n" + ('\textbackslash{}n'.join((f"- stage\_anchor\_\{idx:02d\}: `time=\{anchor.get('time\_utc', '')\}, s\_km=\{anchor.get('s\_km', float('nan'))\}, V=\{anchor.get('voltage\_v', float('nan'))\}, I=\{anchor.get('current\_a', float('nan'))\}, SoC=\{anchor.get('soc\_target', float('nan'))\}, dwell\_sec=\{anchor.get('dwell\_sec', float('nan'))\}`" for idx, anchor in enumerate(stage\_anchors, start=1))) if stage\_anchors else '- (none)') + f'\textbackslash{}n\textbackslash{}n\#\# Day metrics\textbackslash{}n\textbackslash{}n' + ('\textbackslash{}n'.join((f"- day \{row.get('day', '')\}: dist\_end=\{row.get('distance\_end\_km', float('nan')):.1f\} km, final\_soc=\{row.get('final\_soc\_pred', float('nan')):.4f\}, power\_rmse=\{row.get('power\_rmse\_clean\_w', float('nan')):.2f\} W, voltage\_rmse=\{row.get('voltage\_rmse\_clean\_v', float('nan')):.3f\} V, excluded\_power=\{row.get('excluded\_power\_points', 0)\}, excluded\_voltage=\{row.get('excluded\_voltage\_points', 0)\}" for row in day\_metrics)) if day\_metrics else '- (none)') + f'\textbackslash{}n\textbackslash{}n\#\# Evidence bundle\textbackslash{}n\textbackslash{}n' + '\textbackslash{}n'.join((f'- \{key\}: `\{value\}`' for key, value in evidence\_rows if value)) + ('\textbackslash{}n' + '\textbackslash{}n'.join((f'- explicit\_\{key\}: `\{value\}`' for key, value in explicit\_grounded\_assets.items())) if explicit\_grounded\_assets else '') + ('\textbackslash{}n' + '\textbackslash{}n'.join((f'- external\_document: `\{value\}`' for value in external\_documents)) if external\_documents else '') + f'\textbackslash{}n\textbackslash{}n\#\# Active maps\textbackslash{}n\textbackslash{}n' + '\textbackslash{}n'.join((f'- \{key\}: `\{value\}`' for key, value in rel\_maps.items())) + f"\textbackslash{}n\textbackslash{}n\#\# Grounded map provenance\textbackslash{}n\textbackslash{}n- grounded\_map\_summary\_yaml: `\{grounded\_yaml\}`\textbackslash{}n\textbackslash{}n\#\# Post refinement\textbackslash{}n\textbackslash{}n- accepted: `\{post\_refine.accepted\}`\textbackslash{}n- panel\_gain\_factor: `\{post\_refine.panel\_gain\_factor:.5f\}`\textbackslash{}n- cda\_factor: `\{post\_refine.cda\_factor:.5f\}`\textbackslash{}n- crr\_factor: `\{post\_refine.crr\_factor:.5f\}`\textbackslash{}n- drive\_eff\_factor: `\{post\_refine.drive\_eff\_factor:.5f\}`\textbackslash{}n- headwind\_gain\_factor: `\{post\_refine.headwind\_gain\_factor:.5f\}`\textbackslash{}n- e\_nom\_factor: `\{post\_refine.e\_nom\_factor:.5f\}`\textbackslash{}n- rint\_factor: `\{post\_refine.rint\_factor:.5f\}`\textbackslash{}n\textbackslash{}n\#\# Outputs\textbackslash{}n\textbackslash{}n- fit summary: `\{os.path.relpath(summary\_yaml, package\_dir).replace('\textbackslash{}\textbackslash{}', '/')\}`\textbackslash{}n- current maps and coefficients: `\{current\_maps\_rel\_package\}`\textbackslash{}n" の結果を代入する。
\item md\_path.write\_text(...) を実行する。
\item tex に f"\textbackslash{}n\textbackslash{}\textbackslash{}documentclass[a4paper,11pt]\{\{article\}\}\textbackslash{}n\textbackslash{}\textbackslash{}usepackage[top=18mm,bottom=22mm,left=18mm,right=18mm]\{\{geometry\}\}\textbackslash{}n\textbackslash{}\textbackslash{}usepackage\{\{fontspec\}\}\textbackslash{}n\textbackslash{}\textbackslash{}usepackage\{\{xeCJK\}\}\textbackslash{}n\textbackslash{}\textbackslash{}setmainfont\{\{Times New Roman\}\}\textbackslash{}n\textbackslash{}\textbackslash{}setCJKmainfont\{\{Yu Gothic\}\}\textbackslash{}n\textbackslash{}\textbackslash{}setCJKmonofont\{\{Yu Gothic\}\}\textbackslash{}n\textbackslash{}\textbackslash{}usepackage\{\{booktabs\}\}\textbackslash{}n\textbackslash{}\textbackslash{}usepackage\{\{longtable\}\}\textbackslash{}n\textbackslash{}\textbackslash{}usepackage\{\{array\}\}\textbackslash{}n\textbackslash{}\textbackslash{}usepackage\{\{xurl\}\}\textbackslash{}n\textbackslash{}\textbackslash{}usepackage[unicode,hidelinks]\{\{hyperref\}\}\textbackslash{}n\textbackslash{}\textbackslash{}title\{\{Generic Vehicle Identification Report\}\}\textbackslash{}n\textbackslash{}\textbackslash{}author\{\{solar\textbackslash{}\textbackslash{}\_ws0129-main\}\}\textbackslash{}n\textbackslash{}\textbackslash{}date\{\{\}\}\textbackslash{}n\textbackslash{}\textbackslash{}begin\{\{document\}\}\textbackslash{}n\textbackslash{}\textbackslash{}raggedbottom\textbackslash{}n\textbackslash{}\textbackslash{}maketitle\textbackslash{}n\textbackslash{}n\textbackslash{}\textbackslash{}section\{\{Inputs\}\}\textbackslash{}n\textbackslash{}\textbackslash{}begin\{\{itemize\}\}\textbackslash{}n  \textbackslash{}\textbackslash{}item profile: \textbackslash{}\textbackslash{}path\{\{project\_packages/\{package\_dir.name\}/profile.yaml\}\}\textbackslash{}n  \textbackslash{}\textbackslash{}item manifest: \textbackslash{}\textbackslash{}path\{\{\{os.path.relpath(manifest\_path, report\_dir).replace('\textbackslash{}\textbackslash{}', '/')\}\}\}\textbackslash{}n  \textbackslash{}\textbackslash{}item normalized replay log: \textbackslash{}\textbackslash{}path\{\{\{os.path.relpath(observed\_log\_csv, report\_dir).replace('\textbackslash{}\textbackslash{}', '/')\}\}\}\textbackslash{}n\textbackslash{}\textbackslash{}end\{\{itemize\}\}\textbackslash{}n\textbackslash{}n\textbackslash{}\textbackslash{}section\{\{Adopted coefficients\}\}\textbackslash{}n\textbackslash{}\textbackslash{}begin\{\{longtable\}\}\{\{p\{\{0.48\textbackslash{}\textbackslash{}linewidth\}\}p\{\{0.22\textbackslash{}\textbackslash{}linewidth\}\}\}\}\textbackslash{}n\textbackslash{}\textbackslash{}toprule\textbackslash{}n項目 \& 値 \textbackslash{}\textbackslash{}\textbackslash{}\textbackslash{}\textbackslash{}n\textbackslash{}\textbackslash{}midrule\textbackslash{}n\textbackslash{}\textbackslash{}endhead\textbackslash{}npanel gain \& \{pv.panel\_gain:.6f\} \textbackslash{}\textbackslash{}\textbackslash{}\textbackslash{}\textbackslash{}ntcell gain [C/(W m\$\textasciicircum{}\{\{-2\}\}\$)] \& \{pv.tcell\_gain\_c\_per\_wm2:.6f\} \textbackslash{}\textbackslash{}\textbackslash{}\textbackslash{}\textbackslash{}nstop tilt fraction [-] \& \{pv.stop\_tilt\_fraction:.6f\} \textbackslash{}\textbackslash{}\textbackslash{}\textbackslash{}\textbackslash{}nstop solar RMSE [W] \& \{pv.stop\_solar\_rmse\_w:.3f\} \textbackslash{}\textbackslash{}\textbackslash{}\textbackslash{}\textbackslash{}nsolar measurement gain to pack \& \{solar\_gain\_to\_pack:.8f\} \textbackslash{}\textbackslash{}\textbackslash{}\textbackslash{}\textbackslash{}nsolar calibration samples \& \{solar\_calibration\_samples\} \textbackslash{}\textbackslash{}\textbackslash{}\textbackslash{}\textbackslash{}nsolar free-intercept diagnostic [W] \& \{solar\_calibration\_intercept\_w:.6f\} \textbackslash{}\textbackslash{}\textbackslash{}\textbackslash{}\textbackslash{}nsolar daily gain standard deviation \& \{solar\_calibration\_daily\_std:.8f\} \textbackslash{}\textbackslash{}\textbackslash{}\textbackslash{}\textbackslash{}nsoc0 \& \{batt.soc0:.6f\} \textbackslash{}\textbackslash{}\textbackslash{}\textbackslash{}\textbackslash{}n\$E\_\{\{nom\}\}\$ [Wh] \& \{batt.e\_nom\_wh:.3f\} \textbackslash{}\textbackslash{}\textbackslash{}\textbackslash{}\textbackslash{}n\$Q\_\{\{nom\}\}\$ [Ah] \& \{batt.e\_nom\_wh / BATTERY\_NOMINAL\_VOLTAGE\_V:.6f\} \textbackslash{}\textbackslash{}\textbackslash{}\textbackslash{}\textbackslash{}n\$k\_\{\{Rint\}\}\$ \& \{batt.rint\_scale:.6f\} \textbackslash{}\textbackslash{}\textbackslash{}\textbackslash{}\textbackslash{}n\$R\_\{\{line\}\}\$ [ohm] \& \{batt.r\_line\_ohm:.6f\} \textbackslash{}\textbackslash{}\textbackslash{}\textbackslash{}\textbackslash{}n\$R\_\{\{p\}\}\$ [ohm] \& \{float(battery\_dynamic\_fit.get('r\_polarization\_ohm', 0.0)):.6f\} \textbackslash{}\textbackslash{}\textbackslash{}\textbackslash{}\textbackslash{}n\$\textbackslash{}\textbackslash{}tau\_\{\{p\}\}\$ [s] \& \{float(battery\_dynamic\_fit.get('tau\_sec', 60.0)):.3f\} \textbackslash{}\textbackslash{}\textbackslash{}\textbackslash{}\textbackslash{}n\$\textbackslash{}\textbackslash{}eta\_\{\{charge\}\}\$ \& \{batt.eta\_charge:.6f\} \textbackslash{}\textbackslash{}\textbackslash{}\textbackslash{}\textbackslash{}nCdA \& \{mot.cda:.6f\} \textbackslash{}\textbackslash{}\textbackslash{}\textbackslash{}\textbackslash{}nCrr \& \{mot.crr:.6f\} \textbackslash{}\textbackslash{}\textbackslash{}\textbackslash{}\textbackslash{}n\$P\_\{\{aux\}\}\$ [W] \& \{mot.p\_aux\_w:.3f\} \textbackslash{}\textbackslash{}\textbackslash{}\textbackslash{}\textbackslash{}ngrade scale \& \{mot.grade\_scale:.6f\} \textbackslash{}\textbackslash{}\textbackslash{}\textbackslash{}\textbackslash{}ndrive efficiency scale \& \{mot.drive\_eff\_scale:.6f\} \textbackslash{}\textbackslash{}\textbackslash{}\textbackslash{}\textbackslash{}nregen utilization \& \{mot.regen\_utilization:.6f\} \textbackslash{}\textbackslash{}\textbackslash{}\textbackslash{}\textbackslash{}nregen fit samples \& \{mot.regen\_sample\_count\} \textbackslash{}\textbackslash{}\textbackslash{}\textbackslash{}\textbackslash{}nregen subset RMSE [W] \& \{mot.regen\_fit\_rmse\_w:.3f\} \textbackslash{}\textbackslash{}\textbackslash{}\textbackslash{}\textbackslash{}nheadwind gain \& \{mot.headwind\_gain:.6f\} \textbackslash{}\textbackslash{}\textbackslash{}\textbackslash{}\textbackslash{}npack terminal voltage limit [V] \& \{pack\_voltage\_limit\_v:.3f\} \textbackslash{}\textbackslash{}\textbackslash{}\textbackslash{}\textbackslash{}n\textbackslash{}\textbackslash{}bottomrule\textbackslash{}n\textbackslash{}\textbackslash{}end\{\{longtable\}\}\textbackslash{}n\textbackslash{}n\textbackslash{}\textbackslash{}section\{\{Solar measurement calibration\}\}\textbackslash{}nBefore vehicle and PV fitting, stationary samples identify the ZP solar-power\textbackslash{}nchannel against the DC-bus balance,\textbackslash{}n\textbackslash{}\textbackslash{}[\textbackslash{}nP\_\{\{batt,k\}\}=P\_\{\{aux\}\}-g\_\{\{solar\}\}P\_\{\{solar,raw,k\}\}+\textbackslash{}\textbackslash{}varepsilon\_k.\textbackslash{}n\textbackslash{}\textbackslash{}]\textbackslash{}nWith the independently observed 21 W auxiliary load fixed, the bounded robust\textbackslash{}nestimate is\textbackslash{}n\textbackslash{}\textbackslash{}[\textbackslash{}n\textbackslash{}\textbackslash{}hat g\_\{\{solar\}\}=\textbackslash{}\textbackslash{}mathop\{\{\textbackslash{}\textbackslash{}arg\textbackslash{}\textbackslash{}min\}\}\_\{\{0.70\textbackslash{}\textbackslash{}le g\textbackslash{}\textbackslash{}le1.05\}\}\textbackslash{}n\textbackslash{}\textbackslash{}sum\_k\textbackslash{}\textbackslash{}rho\_H\textbackslash{}\textbackslash{}!\textbackslash{}\textbackslash{}left(P\_\{\{batt,k\}\}-P\_\{\{aux\}\}+gP\_\{\{solar,raw,k\}\}\textbackslash{}\textbackslash{}right).\textbackslash{}n\textbackslash{}\textbackslash{}]\textbackslash{}nThe adopted gain is \{solar\_gain\_to\_pack:.8f\} from \{solar\_calibration\_samples\}\textbackslash{}nsamples (accepted=\{solar\_calibration\_accepted\}). A free-intercept fit is reported\textbackslash{}nonly as a physical-consistency diagnostic; its intercept is\textbackslash{}n\{solar\_calibration\_intercept\_w:.6f\} W and its error relative to the fixed\textbackslash{}nauxiliary load is \{solar\_calibration\_intercept\_error\_w:.6f\} W. The day-to-day\textbackslash{}ngain standard deviation is \{solar\_calibration\_daily\_std:.8f\}.\textbackslash{}n\textbackslash{}nThe sender places the uncorrected ZP value in \textbackslash{}\textbackslash{}texttt\{\{solar\textbackslash{}\textbackslash{}\_power\textbackslash{}\textbackslash{}\_w\}\}. The\textbackslash{}nreceiving WiFi bridge applies the profile gain exactly once. A corrected sender\textbackslash{}nvalue must not be placed in that raw field because it would be calibrated twice.\textbackslash{}n\textbackslash{}nThe terminal-voltage limit is grounded in the 25-series product limit of\textbackslash{}n\$4.35\$ V/cell, or \$108.75\$ V/pack, rather than the maximum voltage seen in a\textbackslash{}nloaded discharge trace. The adopted YATA profile limit is\textbackslash{}n\{pack\_voltage\_limit\_v:.3f\} V and must equal that value; the active OCV map must\textbackslash{}nremain below it.\textbackslash{}n\textbackslash{}n\textbackslash{}\textbackslash{}section\{\{Validation\}\}\textbackslash{}nUnprefixed values are vehicle-model validation conditioned on measured array\textbackslash{}npower. The end-to-end values use independent archive GHI and the moving-PV\textbackslash{}nmodel. Target-derived effective irradiance is excluded from both fitting and\textbackslash{}nend-to-end validation. This separation prevents irradiance forecast error from\textbackslash{}nbeing fitted as vehicle resistance.\textbackslash{}n\textbackslash{}nThe field-analysis source states that the ZP logger start time was not recorded\textbackslash{}nand was reconstructed backward from control-stop times. Consequently the 5 s\textbackslash{}npower, speed and DEM-grade records do not constitute a pointwise synchronization\textbackslash{}ncertificate. The configured resampling window is used for parameter fitting,\textbackslash{}nwhile 10 km and 25 km energy residuals are the primary route-energy checks.\textbackslash{}n\textbackslash{}nThe 5 s point RMSE includes unresolved channel synchronization and transient\textbackslash{}nnoise. The 120 s mean-residual RMSE tests local mean power, while the 10 km and\textbackslash{}n25 km metrics test accumulated route energy. They are intentionally reported as\textbackslash{}nseparate quantities; no smaller aggregated number is presented as the 5 s RMSE.\textbackslash{}n\textbackslash{}nConversion efficiency and actual regeneration use are separated as\textbackslash{}n\textbackslash{}\textbackslash{}[\textbackslash{}nP\_\{\{reg,dc\}\}=u\_\{\{regen\}\}\textbackslash{}\textbackslash{}eta\_\{\{reg\}\}\textbackslash{}\textbackslash{}eta\_\{\{gear\}\}\textbackslash{}\textbackslash{}eta\_\{\{inv\}\}\textbackslash{}n\textbackslash{}\textbackslash{}max(-P\_\{\{mech\}\},0),\textbackslash{}\textbackslash{}qquad 0\textbackslash{}\textbackslash{}le u\_\{\{regen\}\}\textbackslash{}\textbackslash{}le 1.\textbackslash{}n\textbackslash{}\textbackslash{}]\textbackslash{}nOnly samples with observable negative mechanical power identify \$u\_\{\{regen\}\}\$.\textbackslash{}n\textbackslash{}n\textbackslash{}\textbackslash{}begin\{\{longtable\}\}\{\{p\{\{0.48\textbackslash{}\textbackslash{}linewidth\}\}p\{\{0.22\textbackslash{}\textbackslash{}linewidth\}\}\}\}\textbackslash{}n\textbackslash{}\textbackslash{}toprule\textbackslash{}n指標 \& 値 \textbackslash{}\textbackslash{}\textbackslash{}\textbackslash{}\textbackslash{}n\textbackslash{}\textbackslash{}midrule\textbackslash{}n\textbackslash{}\textbackslash{}endhead\textbackslash{}nvehicle replay power RMSE [W] \& \{float(metrics.get('power\_rmse\_clean\_w', float('nan'))):.3f\} \textbackslash{}\textbackslash{}\textbackslash{}\textbackslash{}\textbackslash{}nfit-window power RMSE [W] \& \{float(metrics.get('power\_rmse\_fit\_window\_w', float('nan'))):.3f\} \textbackslash{}\textbackslash{}\textbackslash{}\textbackslash{}\textbackslash{}nvehicle replay voltage RMSE [V] \& \{float(metrics.get('voltage\_rmse\_clean\_v', float('nan'))):.3f\} \textbackslash{}\textbackslash{}\textbackslash{}\textbackslash{}\textbackslash{}nvehicle replay final SoC [-] \& \{float(metrics.get('final\_soc\_pred', float('nan'))):.6f\} \textbackslash{}\textbackslash{}\textbackslash{}\textbackslash{}\textbackslash{}nvehicle terminal SoC observed [-] \& \{float(metrics.get('retire\_anchor\_soc\_obs', float('nan'))):.6f\} \textbackslash{}\textbackslash{}\textbackslash{}\textbackslash{}\textbackslash{}nvehicle terminal SoC predicted [-] \& \{float(metrics.get('retire\_anchor\_soc\_pred', float('nan'))):.6f\} \textbackslash{}\textbackslash{}\textbackslash{}\textbackslash{}\textbackslash{}nvehicle terminal SoC error [-] \& \{float(metrics.get('retire\_anchor\_soc\_error', float('nan'))):.6f\} \textbackslash{}\textbackslash{}\textbackslash{}\textbackslash{}\textbackslash{}nbattery-only power RMSE [W] \& \{float(metrics.get('battery\_conditional\_power\_rmse\_clean\_w', float('nan'))):.3f\} \textbackslash{}\textbackslash{}\textbackslash{}\textbackslash{}\textbackslash{}nbattery-only voltage RMSE [V] \& \{float(metrics.get('battery\_conditional\_voltage\_rmse\_clean\_v', float('nan'))):.3f\} \textbackslash{}\textbackslash{}\textbackslash{}\textbackslash{}\textbackslash{}nbattery-only terminal SoC error [-] \& \{float(metrics.get('battery\_conditional\_retire\_anchor\_soc\_error', float('nan'))):.6f\} \textbackslash{}\textbackslash{}\textbackslash{}\textbackslash{}\textbackslash{}nend-to-end power RMSE [W] \& \{float(metrics.get('end\_to\_end\_power\_rmse\_clean\_w', float('nan'))):.3f\} \textbackslash{}\textbackslash{}\textbackslash{}\textbackslash{}\textbackslash{}nend-to-end voltage RMSE [V] \& \{float(metrics.get('end\_to\_end\_voltage\_rmse\_clean\_v', float('nan'))):.3f\} \textbackslash{}\textbackslash{}\textbackslash{}\textbackslash{}\textbackslash{}nend-to-end final SoC [-] \& \{float(metrics.get('end\_to\_end\_final\_soc\_pred', float('nan'))):.6f\} \textbackslash{}\textbackslash{}\textbackslash{}\textbackslash{}\textbackslash{}nend-to-end terminal SoC error [-] \& \{float(metrics.get('end\_to\_end\_retire\_anchor\_soc\_error', float('nan'))):.6f\} \textbackslash{}\textbackslash{}\textbackslash{}\textbackslash{}\textbackslash{}nend-to-end moving PV RMSE [W] \& \{float(metrics.get('end\_to\_end\_moving\_pv\_rmse\_w', float('nan'))):.3f\} \textbackslash{}\textbackslash{}\textbackslash{}\textbackslash{}\textbackslash{}nend-to-end deployed-stop PV RMSE [W] \& \{float(metrics.get('end\_to\_end\_deployed\_stop\_pv\_rmse\_w', float('nan'))):.3f\} \textbackslash{}\textbackslash{}\textbackslash{}\textbackslash{}\textbackslash{}nPV leave-one-day-out moving RMSE [W] \& \{float(metrics.get('pv\_lodo\_moving\_rmse\_w', float('nan'))):.3f\} \textbackslash{}\textbackslash{}\textbackslash{}\textbackslash{}\textbackslash{}nPV leave-one-day-out deployed-stop RMSE [W] \& \{float(metrics.get('pv\_lodo\_deployed\_stop\_rmse\_w', float('nan'))):.3f\} \textbackslash{}\textbackslash{}\textbackslash{}\textbackslash{}\textbackslash{}nPV leave-one-day-out folds [-] \& \{int(metrics.get('pv\_lodo\_fold\_count', 0))\} \textbackslash{}\textbackslash{}\textbackslash{}\textbackslash{}\textbackslash{}n120 s mean-residual RMSE [W] \& \{float(metrics.get('power\_residual\_mean\_120s\_rmse\_w', float('nan'))):.3f\} \textbackslash{}\textbackslash{}\textbackslash{}\textbackslash{}\textbackslash{}n10 km energy-error RMSE [Wh] \& \{float(metrics.get('energy\_error\_10km\_rmse\_wh', float('nan'))):.3f\} \textbackslash{}\textbackslash{}\textbackslash{}\textbackslash{}\textbackslash{}n25 km energy-error RMSE [Wh] \& \{float(metrics.get('energy\_error\_25km\_rmse\_wh', float('nan'))):.3f\} \textbackslash{}\textbackslash{}\textbackslash{}\textbackslash{}\textbackslash{}n\textbackslash{}\textbackslash{}bottomrule\textbackslash{}n\textbackslash{}\textbackslash{}end\{\{longtable\}\}\textbackslash{}n\textbackslash{}n\textbackslash{}\textbackslash{}section\{\{Terminal-SoC evidence and certification\}\}\textbackslash{}nThe independent evidence interval is\textbackslash{}n\$[\{terminal\_evidence\_min:.6f\},\textbackslash{}\textbackslash{},\{terminal\_evidence\_max:.6f\}]\$, with width\textbackslash{}n\{terminal\_evidence\_spread\_pp:.3f\} percentage points. The complete\textbackslash{}nhigh-precision gate result is \textbackslash{}\textbackslash{}textbf\{\{\{terminal\_gate\_pass\}\}\}.\textbackslash{}n\textbackslash{}n\textbackslash{}\textbackslash{}begin\{\{longtable\}\}\{\{p\{\{0.58\textbackslash{}\textbackslash{}linewidth\}\}p\{\{0.18\textbackslash{}\textbackslash{}linewidth\}\}\}\}\textbackslash{}n\textbackslash{}\textbackslash{}toprule\textbackslash{}ncheck \& pass \textbackslash{}\textbackslash{}\textbackslash{}\textbackslash{}\textbackslash{}n\textbackslash{}\textbackslash{}midrule\textbackslash{}n\textbackslash{}\textbackslash{}endhead\textbackslash{}n" + '\textbackslash{}n'.join((f'\{tex\_path\_fragment(str(key))\} \& \{bool(value)\} \textbackslash{}\textbackslash{}\textbackslash{}\textbackslash{}' for key, value in terminal\_checks.items())) + f"\textbackslash{}n\textbackslash{}\textbackslash{}bottomrule\textbackslash{}n\textbackslash{}\textbackslash{}end\{\{longtable\}\}\textbackslash{}n\textbackslash{}n\{tex\_text\_fragment(terminal\_interpretation)\} The central estimate is not a\textbackslash{}ndirect coulomb-counter observation. If this gate is false, it remains an\textbackslash{}nengineering anchor only; neither a high-precision actual 2831 km SoC nor\textbackslash{}noperational model certification is claimed.\textbackslash{}n\textbackslash{}n\textbackslash{}\textbackslash{}section\{\{Method references\}\}\textbackslash{}n\textbackslash{}\textbackslash{}begin\{\{itemize\}\}\textbackslash{}n  \textbackslash{}\textbackslash{}item Open-Meteo Historical Weather API: \textbackslash{}\textbackslash{}url\{\{https://open-meteo.com/en/docs/historical-weather-api\}\}\textbackslash{}n  \textbackslash{}\textbackslash{}item Sandia PVPMC plane-of-array irradiance guidance: \textbackslash{}\textbackslash{}url\{\{https://pvpmc.sandia.gov/modeling-guide/1-weather-design-inputs/plane-of-array-poa-irradiance/\}\}\textbackslash{}n  \textbackslash{}\textbackslash{}item pvlib total-irradiance transposition API: \textbackslash{}\textbackslash{}url\{\{https://pvlib-python.readthedocs.io/en/latest/reference/generated/pvlib.irradiance.get\_total\_irradiance.html\}\}\textbackslash{}n  \textbackslash{}\textbackslash{}item Byrd et al. (1995), L-BFGS-B: \textbackslash{}\textbackslash{}url\{\{https://doi.org/10.1137/0916069\}\}\textbackslash{}n  \textbackslash{}\textbackslash{}item Liu et al. (2014), GPS road-grade quantification: \textbackslash{}\textbackslash{}url\{\{https://doi.org/10.1016/j.atmosenv.2013.12.025\}\}\textbackslash{}n  \textbackslash{}\textbackslash{}item Wang et al. (2018), synchronized link-level EV energy data: \textbackslash{}\textbackslash{}url\{\{https://escholarship.org/uc/item/2bp8x04q\}\}\textbackslash{}n\textbackslash{}\textbackslash{}end\{\{itemize\}\}\textbackslash{}n\textbackslash{}n\textbackslash{}\textbackslash{}section\{\{Terminal anchor\}\}\textbackslash{}n\textbackslash{}\textbackslash{}begin\{\{longtable\}\}\{\{p\{\{0.48\textbackslash{}\textbackslash{}linewidth\}\}p\{\{0.32\textbackslash{}\textbackslash{}linewidth\}\}\}\}\textbackslash{}n\textbackslash{}\textbackslash{}toprule\textbackslash{}n項目 \& 値 \textbackslash{}\textbackslash{}\textbackslash{}\textbackslash{}\textbackslash{}n\textbackslash{}\textbackslash{}midrule\textbackslash{}n\textbackslash{}\textbackslash{}endhead\textbackslash{}nanchor distance [km] \& \{float(terminal\_anchor.get('s\_km', float('nan'))):.3f\} \textbackslash{}\textbackslash{}\textbackslash{}\textbackslash{}\textbackslash{}nanchor time UTC \& \textbackslash{}\textbackslash{}path\{\{\{str(terminal\_anchor.get('time\_utc', ''))\}\}\} \textbackslash{}\textbackslash{}\textbackslash{}\textbackslash{}\textbackslash{}nanchor voltage [V] \& \{float(terminal\_anchor.get('voltage\_v', float('nan'))):.3f\} \textbackslash{}\textbackslash{}\textbackslash{}\textbackslash{}\textbackslash{}nanchor current [A] \& \{float(terminal\_anchor.get('current\_a', float('nan'))):.3f\} \textbackslash{}\textbackslash{}\textbackslash{}\textbackslash{}\textbackslash{}nanchor temperature [C] \& \{float(terminal\_anchor.get('temp\_c', float('nan'))):.3f\} \textbackslash{}\textbackslash{}\textbackslash{}\textbackslash{}\textbackslash{}nanchor SoC target [-] \& \{float(terminal\_anchor.get('soc\_target', float('nan'))):.6f\} \textbackslash{}\textbackslash{}\textbackslash{}\textbackslash{}\textbackslash{}n\textbackslash{}\textbackslash{}bottomrule\textbackslash{}n\textbackslash{}\textbackslash{}end\{\{longtable\}\}\textbackslash{}n\textbackslash{}n\textbackslash{}\textbackslash{}section\{\{Fit plan\}\}\textbackslash{}n\textbackslash{}\textbackslash{}begin\{\{itemize\}\}\textbackslash{}n  \textbackslash{}\textbackslash{}item quality: \{str(fit\_plan.get('quality', ''))\}\textbackslash{}n  \textbackslash{}\textbackslash{}item battery restart count: \{fit\_plan.get('battery\_restart\_count', '')\}\textbackslash{}n  \textbackslash{}\textbackslash{}item battery maxiter: \{fit\_plan.get('battery\_maxiter', '')\}\textbackslash{}n  \textbackslash{}\textbackslash{}item motion restart count: \{fit\_plan.get('motion\_restart\_count', '')\}\textbackslash{}n  \textbackslash{}\textbackslash{}item motion maxiter: \{fit\_plan.get('motion\_maxiter', '')\}\textbackslash{}n  \textbackslash{}\textbackslash{}item joint restart count: \{fit\_plan.get('joint\_restart\_count', '')\}\textbackslash{}n  \textbackslash{}\textbackslash{}item joint random start count: \{fit\_plan.get('joint\_random\_start\_count', '')\}\textbackslash{}n  \textbackslash{}\textbackslash{}item joint local topk: \{fit\_plan.get('joint\_local\_topk', '')\}\textbackslash{}n  \textbackslash{}\textbackslash{}item joint maxiter: \{fit\_plan.get('joint\_maxiter', '')\}\textbackslash{}n  \textbackslash{}\textbackslash{}item fit stride: \{fit\_plan.get('fit\_stride', '')\}\textbackslash{}n  \textbackslash{}\textbackslash{}item allow map-shape fit: \{fit\_plan.get('allow\_map\_shape\_fit', '')\}\textbackslash{}n  \textbackslash{}\textbackslash{}item post-refine enabled: \{fit\_plan.get('post\_refine\_enabled', '')\}\textbackslash{}n  \textbackslash{}\textbackslash{}item terminal anchor role: \textbackslash{}\textbackslash{}path\{\{\{str(fit\_plan.get('terminal\_anchor\_role', ''))\}\}\}\textbackslash{}n  \textbackslash{}\textbackslash{}item panel deployment stopped speed: \{fit\_plan.get('panel\_deployment\_stopped\_speed\_kmh', '')\} km/h\textbackslash{}n  \textbackslash{}\textbackslash{}item panel deployment minimum dwell: \{fit\_plan.get('panel\_deployment\_min\_dwell\_sec', '')\} s\textbackslash{}n  \textbackslash{}\textbackslash{}item panel deployment maximum sample gap: \{fit\_plan.get('panel\_deployment\_max\_sample\_gap\_sec', '')\} s\textbackslash{}n  \textbackslash{}\textbackslash{}item acceleration observation filter/alignment adopted: \{bool((fit\_plan.get('acceleration\_observation\_fit', \{\}) or \{\}).get('adopted', False))\}\textbackslash{}n  \textbackslash{}\textbackslash{}item acceleration selected filter: \textbackslash{}\textbackslash{}path\{\{\{str((fit\_plan.get('acceleration\_observation\_fit', \{\}) or \{\}).get('selected\_filter\_method', 'legacy'))\}\}\}, window=\{(fit\_plan.get('acceleration\_observation\_fit', \{\}) or \{\}).get('selected\_filter\_window\_samples', '')\} samples / \{(fit\_plan.get('acceleration\_observation\_fit', \{\}) or \{\}).get('selected\_filter\_window\_sec', '')\} s\textbackslash{}n  \textbackslash{}\textbackslash{}item acceleration timestamp selected lag: \{float((fit\_plan.get('acceleration\_observation\_fit', \{\}) or \{\}).get('selected\_lag\_sec', float('nan'))):.3f\} s\textbackslash{}n  \textbackslash{}\textbackslash{}item acceleration alignment train RMSE: \{float((fit\_plan.get('acceleration\_observation\_fit', \{\}) or \{\}).get('baseline\_training\_rmse\_w', float('nan'))):.3f\} \$\textbackslash{}\textbackslash{}rightarrow\$ \{float((fit\_plan.get('acceleration\_observation\_fit', \{\}) or \{\}).get('selected\_training\_rmse\_w', float('nan'))):.3f\} W\textbackslash{}n  \textbackslash{}\textbackslash{}item acceleration alignment held-out RMSE: \{float((fit\_plan.get('acceleration\_observation\_fit', \{\}) or \{\}).get('baseline\_validation\_rmse\_w', float('nan'))):.3f\} \$\textbackslash{}\textbackslash{}rightarrow\$ \{float((fit\_plan.get('acceleration\_observation\_fit', \{\}) or \{\}).get('selected\_validation\_rmse\_w', float('nan'))):.3f\} W\textbackslash{}n  \textbackslash{}\textbackslash{}item acceleration held-out RMSE ratio: \{float((fit\_plan.get('acceleration\_observation\_fit', \{\}) or \{\}).get('validation\_rmse\_ratio', float('nan'))):.6f\}\textbackslash{}n  \textbackslash{}\textbackslash{}item DEM grade observation adopted: \{bool(grade\_observation\_fit.get('adopted', False))\}\textbackslash{}n  \textbackslash{}\textbackslash{}item DEM grade smoothing window: \{grade\_observation\_fit.get('selected\_smoothing\_window\_km', '')\} km\textbackslash{}n  \textbackslash{}\textbackslash{}item DEM grade distance offset: \{grade\_observation\_fit.get('selected\_distance\_offset\_km', '')\} km\textbackslash{}n  \textbackslash{}\textbackslash{}item DEM grade train RMSE: \{float(grade\_observation\_fit.get('baseline\_training\_rmse\_w', float('nan'))):.3f\} \$\textbackslash{}\textbackslash{}rightarrow\$ \{float(grade\_observation\_fit.get('selected\_training\_rmse\_w', float('nan'))):.3f\} W\textbackslash{}n  \textbackslash{}\textbackslash{}item DEM grade held-out RMSE: \{float(grade\_observation\_fit.get('baseline\_validation\_rmse\_w', float('nan'))):.3f\} \$\textbackslash{}\textbackslash{}rightarrow\$ \{float(grade\_observation\_fit.get('selected\_validation\_rmse\_w', float('nan'))):.3f\} W\textbackslash{}n\textbackslash{}\textbackslash{}end\{\{itemize\}\}\textbackslash{}n\textbackslash{}n\textbackslash{}\textbackslash{}section\{\{Stage anchors\}\}\textbackslash{}n\textbackslash{}\textbackslash{}begin\{\{itemize\}\}\textbackslash{}n  \textbackslash{}\textbackslash{}item stage anchor count: \{len(stage\_anchors)\}\textbackslash{}n" + ('\textbackslash{}n' + '\textbackslash{}n'.join((f"  \textbackslash{}\textbackslash{}item anchor \{idx:02d\}: time=\textbackslash{}\textbackslash{}path\{\{\{str(anchor.get('time\_utc', ''))\}\}\}, s=\{float(anchor.get('s\_km', float('nan'))):.3f\} km, V=\{float(anchor.get('voltage\_v', float('nan'))):.3f\} V, I=\{float(anchor.get('current\_a', float('nan'))):.3f\} A, SoC=\{float(anchor.get('soc\_target', float('nan'))):.6f\}, dwell=\{float(anchor.get('dwell\_sec', float('nan'))):.1f\} s" for idx, anchor in enumerate(stage\_anchors, start=1))) if stage\_anchors else '\textbackslash{}n  \textbackslash{}\textbackslash{}item (none)') + f'\textbackslash{}n\textbackslash{}\textbackslash{}end\{\{itemize\}\}\textbackslash{}n\textbackslash{}n\textbackslash{}\textbackslash{}section\{\{Day metrics\}\}\textbackslash{}n\textbackslash{}\textbackslash{}begin\{\{longtable\}\}\{\{>\{\{\textbackslash{}\textbackslash{}raggedright\textbackslash{}\textbackslash{}arraybackslash\}\}p\{\{0.10\textbackslash{}\textbackslash{}linewidth\}\}>\{\{\textbackslash{}\textbackslash{}raggedright\textbackslash{}\textbackslash{}arraybackslash\}\}p\{\{0.16\textbackslash{}\textbackslash{}linewidth\}\}>\{\{\textbackslash{}\textbackslash{}raggedright\textbackslash{}\textbackslash{}arraybackslash\}\}p\{\{0.16\textbackslash{}\textbackslash{}linewidth\}\}>\{\{\textbackslash{}\textbackslash{}raggedright\textbackslash{}\textbackslash{}arraybackslash\}\}p\{\{0.16\textbackslash{}\textbackslash{}linewidth\}\}>\{\{\textbackslash{}\textbackslash{}raggedright\textbackslash{}\textbackslash{}arraybackslash\}\}p\{\{0.16\textbackslash{}\textbackslash{}linewidth\}\}>\{\{\textbackslash{}\textbackslash{}raggedright\textbackslash{}\textbackslash{}arraybackslash\}\}p\{\{0.10\textbackslash{}\textbackslash{}linewidth\}\}\}\}\textbackslash{}n\textbackslash{}\textbackslash{}toprule\textbackslash{}nday \& dist end [km] \& final SoC \& power RMSE [W] \& voltage RMSE [V] \& excl. power \textbackslash{}\textbackslash{}\textbackslash{}\textbackslash{}\textbackslash{}n\textbackslash{}\textbackslash{}midrule\textbackslash{}n\textbackslash{}\textbackslash{}endhead\textbackslash{}n' + '\textbackslash{}n'.join((f"\{int(row.get('day', -1))\} \& \{float(row.get('distance\_end\_km', float('nan'))):.1f\} \& \{float(row.get('final\_soc\_pred', float('nan'))):.4f\} \& \{float(row.get('power\_rmse\_clean\_w', float('nan'))):.2f\} \& \{float(row.get('voltage\_rmse\_clean\_v', float('nan'))):.3f\} \& \{int(row.get('excluded\_power\_points', 0))\} \textbackslash{}\textbackslash{}\textbackslash{}\textbackslash{}" for row in day\_metrics)) + f'\textbackslash{}n\textbackslash{}\textbackslash{}bottomrule\textbackslash{}n\textbackslash{}\textbackslash{}end\{\{longtable\}\}\textbackslash{}n\textbackslash{}n\textbackslash{}\textbackslash{}section\{\{Evidence bundle\}\}\textbackslash{}n\textbackslash{}\textbackslash{}begin\{\{itemize\}\}\textbackslash{}n' + '\textbackslash{}n'.join((f'  \textbackslash{}\textbackslash{}item \{tex\_path\_fragment(key)\}: \{tex\_path\_fragment(value)\}' for key, value in evidence\_rows if value)) + ('\textbackslash{}n' + '\textbackslash{}n'.join((f'  \textbackslash{}\textbackslash{}item explicit \{tex\_path\_fragment(key)\}: \{tex\_path\_fragment(value)\}' for key, value in explicit\_grounded\_assets.items())) if explicit\_grounded\_assets else '') + ('\textbackslash{}n' + '\textbackslash{}n'.join((f'  \textbackslash{}\textbackslash{}item external document: \{tex\_path\_fragment(value)\}' for value in external\_documents)) if external\_documents else '') + f'\textbackslash{}n\textbackslash{}\textbackslash{}end\{\{itemize\}\}\textbackslash{}n\textbackslash{}n\textbackslash{}\textbackslash{}section\{\{Active maps\}\}\textbackslash{}n\textbackslash{}\textbackslash{}begin\{\{itemize\}\}\textbackslash{}n' + '\textbackslash{}n'.join((f'  \textbackslash{}\textbackslash{}item \{tex\_path\_fragment(key)\}: \{tex\_path\_fragment(value)\}' for key, value in rel\_maps.items())) + f"\textbackslash{}n\textbackslash{}\textbackslash{}end\{\{itemize\}\}\textbackslash{}n\textbackslash{}n\textbackslash{}\textbackslash{}section\{\{Grounded map provenance\}\}\textbackslash{}n\textbackslash{}\textbackslash{}begin\{\{itemize\}\}\textbackslash{}n  \textbackslash{}\textbackslash{}item grounded map summary yaml: \textbackslash{}\textbackslash{}path\{\{\{grounded\_yaml\}\}\}\textbackslash{}n\textbackslash{}\textbackslash{}end\{\{itemize\}\}\textbackslash{}n\textbackslash{}n\textbackslash{}\textbackslash{}section\{\{Post refinement\}\}\textbackslash{}n\textbackslash{}\textbackslash{}begin\{\{longtable\}\}\{\{p\{\{0.48\textbackslash{}\textbackslash{}linewidth\}\}p\{\{0.22\textbackslash{}\textbackslash{}linewidth\}\}\}\}\textbackslash{}n\textbackslash{}\textbackslash{}toprule\textbackslash{}n項目 \& 値 \textbackslash{}\textbackslash{}\textbackslash{}\textbackslash{}\textbackslash{}n\textbackslash{}\textbackslash{}midrule\textbackslash{}n\textbackslash{}\textbackslash{}endhead\textbackslash{}naccepted \& \{post\_refine.accepted\} \textbackslash{}\textbackslash{}\textbackslash{}\textbackslash{}\textbackslash{}npanel gain factor \& \{post\_refine.panel\_gain\_factor:.5f\} \textbackslash{}\textbackslash{}\textbackslash{}\textbackslash{}\textbackslash{}nCdA factor \& \{post\_refine.cda\_factor:.5f\} \textbackslash{}\textbackslash{}\textbackslash{}\textbackslash{}\textbackslash{}nCrr factor \& \{post\_refine.crr\_factor:.5f\} \textbackslash{}\textbackslash{}\textbackslash{}\textbackslash{}\textbackslash{}ndrive efficiency factor \& \{post\_refine.drive\_eff\_factor:.5f\} \textbackslash{}\textbackslash{}\textbackslash{}\textbackslash{}\textbackslash{}nheadwind factor \& \{post\_refine.headwind\_gain\_factor:.5f\} \textbackslash{}\textbackslash{}\textbackslash{}\textbackslash{}\textbackslash{}n\$E\_\{\{nom\}\}\$ factor \& \{post\_refine.e\_nom\_factor:.5f\} \textbackslash{}\textbackslash{}\textbackslash{}\textbackslash{}\textbackslash{}n\$R\_\{\{int\}\}\$ factor \& \{post\_refine.rint\_factor:.5f\} \textbackslash{}\textbackslash{}\textbackslash{}\textbackslash{}\textbackslash{}n\textbackslash{}\textbackslash{}bottomrule\textbackslash{}n\textbackslash{}\textbackslash{}end\{\{longtable\}\}\textbackslash{}n\textbackslash{}n\textbackslash{}\textbackslash{}section\{\{Outputs\}\}\textbackslash{}n\textbackslash{}\textbackslash{}begin\{\{itemize\}\}\textbackslash{}n  \textbackslash{}\textbackslash{}item summary YAML: \textbackslash{}\textbackslash{}path\{\{\{os.path.relpath(summary\_yaml, report\_dir).replace('\textbackslash{}\textbackslash{}', '/')\}\}\}\textbackslash{}n  \textbackslash{}\textbackslash{}item current maps and coefficients: \textbackslash{}\textbackslash{}path\{\{\{current\_maps\_rel\_report\}\}\}\textbackslash{}n\textbackslash{}\textbackslash{}end\{\{itemize\}\}\textbackslash{}n\textbackslash{}n\textbackslash{}\textbackslash{}end\{\{document\}\}\textbackslash{}n" の結果を代入する。
\item tex\_path.write\_text(...) を実行する。
\item compile\_tex(...) を実行する。
\item (md\_path, pdf\_path) を返す。
            \end{enumerate}
            \paragraph{代表コード断片}
            \begin{lstlisting}[language=Python]
def write_generic_report(
    package_dir: Path,
    profile_yaml: Path,
    manifest_path: Path,
    summary_yaml: Path,
    observed_log_csv: Path,
    pv: PvFitResult,
    batt: BatteryFitResult,
    mot: MotionFitResult,
    metrics: Dict[str, float],
    post_refine: PostRefineResult,
    map_assets: Dict[str, Path],
    *,
    terminal_anchor: Dict[str, float] | None = None,
    stage_anchors: list[dict] | None = None,
    day_metrics: list[dict] | None = None,
    battery_dynamic_fit: Dict[str, Any] | None = None,
    fit_plan: dict | None = None,
    grounded_map_summary: Dict[str, object] | None = None,
    manifest_context: dict | None = None,
    terminal_consistency: Dict[str, Any] | None = None,
    report_dir: Path | None = None,
    current_maps_path: Path | None = None,
) -> Tuple[Path, Path]:
    report_dir = report_dir or package_dir / "outputs" / "reports"
    ensure_dir(report_dir)
    md_path = report_dir / f"{package_dir.name}_generic_identification_report.md"
    tex_path = report_dir / f"{package_dir.name}_generic_identification_report.tex"
    pdf_path = tex_path.with_suffix(".pdf")
    rel_maps = {key: os.path.relpath(path, package_dir).replace("\\", "/") for key, path in map_assets.items()}
    terminal_anchor = terminal_anchor or {}
    stage_anchors = stage_anchors or []
    day_metrics = day_metrics or []
    battery_dynamic_fit = battery_dynamic_fit or {}
    fit_plan = fit_plan or {}
...
\end{lstlisting}

\subsection{L3075 関数 main}
            \begin{itemize}[leftmargin=1.6em]
              \item 行範囲: L3075--L3841
              \item 所属: module直下
              \item 定義: \path|main() -> None|
              \item docstring: 明示されていない。
              \item このブロックが直接呼ぶ主な関数・メソッド: ArgumentParser, Path, PostRefineResult, R\_int, ValueError, add\_argument, add\_battery\_conditional\_metrics, add\_end\_to\_end\_metrics, apply\_battery\_polarization, attach\_archive\_pv\_model, bool, build\_grounded\_map\_assets
              \item 戻り値の要点: 戻り値が無いか、return 文が明示されていない。
              \item 主なローカル代入名: \_, acceleration\_observation\_fit, adopted\_map\_assets, adopted\_model, adopted\_ocv\_df, adopted\_route\_profile, adopted\_score, adopted\_stage\_anchors, allow\_map\_shape\_fit, ap, args, base\_cycle\_batt\_fit, base\_cycle\_battery\_conditioned\_replay\_df, base\_cycle\_end\_to\_end\_replay\_df, base\_cycle\_logs, base\_cycle\_metrics, base\_cycle\_mot\_fit, base\_cycle\_ocv\_df, base\_cycle\_pv\_fit, base\_cycle\_replay\_df
              \item 読み取る主なインスタンス属性: 特になし。
              \item 更新する主なインスタンス属性: 特になし。
              \item 明示的に送出する例外: ValueError(f'identification options disagree between profile and manifest: \{names\}. Keep duplicate YAML values identical or define each option in only one file.')
              \item 制御構造の規模: 条件分岐 19、ループ 1、try 0
            \end{itemize}
            \paragraph{この定義を読むためのPython構文}
            \begin{itemize}[leftmargin=1.6em]
            \item `->`以後は戻り値の型注釈であり、実行時の値を自動変換する命令ではない。
\item 内包表記は、反復・条件・値生成を一つの式で記述してcollectionを作る。
\item with文はcontext managerに開始・終了処理を任せ、ファイルなどの資源を確実に解放する。
            \end{itemize}
            \paragraph{上から順の処理}
            \begin{enumerate}[leftmargin=1.8em]
            \item ap に argparse.ArgumentParser() の結果を代入する。
\item ap.add\_argument(...) を実行する。
\item ap.add\_argument(...) を実行する。
\item ap.add\_argument(...) を実行する。
\item ap.add\_argument(...) を実行する。
\item ap.add\_argument(...) を実行する。
\item ap.add\_argument(...) を実行する。
\item ap.add\_argument(...) を実行する。
\item ap.add\_argument(...) を実行する。
\item ap.add\_argument(...) を実行する。
\item args に ap.parse\_args() の結果を代入する。
\item profile\_yaml に resolve\_relative(ROOT, args.profile) の結果を代入する。
\item package\_dir に profile\_yaml.parent の結果を代入する。
\item stage(...) を実行する。
\item profile\_cfg に load\_profile\_yaml(profile\_yaml) の結果を代入する。
\item (manifest\_path, manifest) に load\_manifest(package\_dir, args.manifest) の結果を代入する。
\item manifest\_context に resolve\_manifest\_context(package\_dir, manifest) の結果を代入する。
\item inputs に manifest\_context['inputs'] の結果を代入する。
\item options に dict(manifest\_context['options']) の結果を代入する。
\item identification\_cfg に profile\_cfg.get('identification', \{\}) or \{\} の結果を代入する。
\item option\_override\_keys に ('fit\_quality', 'rebuild\_grounded\_base\_maps', 'use\_grounded\_base\_maps', 'allow\_map\_shape\_fit', 'post\_refine\_enabled', 'sensor\_filter', 'acceleration\_observation', 'grade\_observation', 'panel\_deployment\_stopped\_speed\_kmh', 'panel\_deployment\_min\_dwell\_sec', 'panel\_deployment\_max\_sample\_gap\_sec', 'panel\_control\_stop\_tolerance\_km', 'battery\_restart\_count', 'battery\_maxiter', 'motion\_restart\_count', 'motion\_maxiter', 'joint\_restart\_count', 'joint\_random\_start\_count', 'joint\_local\_topk', 'joint\_maxiter', 'fit\_stride') の結果を代入する。
\item conflicting\_options に [key for key in option\_override\_keys if key in identification\_cfg and key in options and (identification\_cfg[key] != options[key])] の結果を代入する。
\item 条件 conflicting\_options を判定し、真なら内部処理を行う。
\item   names に ', '.join(conflicting\_options) の結果を代入する。
\item   ValueError(f'identification options disagree between profile and manifest: \{names\}. Keep duplicate YAML values identical or define each option in only one file.') を送出する。
\item option\_override\_keys を順に走査し、各要素を key に入れて処理する。
\item   条件 key in identification\_cfg を判定し、真なら内部処理を行う。
\item     options[key] に identification\_cfg[key] の結果を代入する。
\item fit\_plan に resolve\_fit\_plan(options, quality=args.quality) の結果を代入する。
\item panel\_deployment\_options に \{'stopped\_speed\_kmh': float(fit\_plan['panel\_deployment\_stopped\_speed\_kmh']), 'deployment\_min\_dwell\_sec': float(fit\_plan['panel\_deployment\_min\_dwell\_sec']), 'deployment\_max\_sample\_gap\_sec': float(fit\_plan['panel\_deployment\_max\_sample\_gap\_sec']), 'horizontal\_control\_stop\_km': declared\_control\_stop\_km(profile\_cfg, profile\_yaml), 'control\_stop\_tolerance\_km': float(fit\_plan['panel\_control\_stop\_tolerance\_km'])\} の結果を代入する。
\item output\_layout に resolve\_identification\_output\_layout(package\_dir, profile\_cfg, output\_tag\_override=args.output\_tag) の結果を代入する。
\item run\_output\_dir に Path(output\_layout['run\_root']) の結果を代入する。
\item report\_output\_dir に Path(output\_layout['report\_root']) の結果を代入する。
\item 条件 args.manifest\_only を判定し、真なら内部処理を行う。
\item   payload に \{'profile\_yaml': relpath\_from(ROOT, profile\_yaml), 'manifest\_yaml': relpath\_from(ROOT, manifest\_path), 'actual\_event\_yaml': relpath\_from(package\_dir, manifest\_context.get('actual\_event\_path')), 'counterfactual\_event\_yaml': relpath\_from(package\_dir, manifest\_context.get('counterfactual\_event\_path')), 'terminal\_anchor\_yaml': relpath\_from(package\_dir, manifest\_context.get('terminal\_anchor\_path')), 'grounded\_map\_summary\_yaml': relpath\_from(package\_dir, manifest\_context.get('grounded\_summary\_path')), 'source\_inventory\_json': relpath\_from(package\_dir, manifest\_context.get('source\_inventory\_path')), 'notes\_markdown': relpath\_from(package\_dir, manifest\_context.get('notes\_markdown\_path')), 'explicit\_grounded\_assets': \{key: relpath\_from(package\_dir, path) for key, path in manifest\_context.get('explicit\_grounded\_assets', \{\}).items()\}, 'external\_documents': list(manifest\_context.get('external\_documents', []))\} の結果を代入する。
\item   print(...) を実行する。
\item    を返す。
\item base\_model に build\_model\_from\_profile\_cfg(profile\_cfg, profile\_yaml) の結果を代入する。
\item source\_map\_assets に build\_source\_map\_assets(profile\_cfg, profile\_yaml) の結果を代入する。
\item grounded\_map\_summary に dict(manifest\_context.get('grounded\_summary\_payload', \{\}) or \{\}) を代入する。
\item explicit\_grounded\_assets に manifest\_context.get('explicit\_grounded\_assets', \{\}) の結果を代入する。
\item rebuild\_grounded に bool(args.rebuild\_grounded\_base\_maps or options.get('rebuild\_grounded\_base\_maps', False)) の結果を代入する。
\item 条件 rebuild\_grounded を判定し、真なら内部処理を行う。
\item   stage(...) を実行する。
\item   (source\_map\_assets, grounded\_map\_summary) に build\_grounded\_map\_assets(profile\_cfg, Path(output\_layout['grounded\_maps']), pv\_area\_m2=float((profile\_cfg.get('model', \{\}) or \{\}).get('pv\_area', 6.0)), base\_dir=package\_dir) の結果を代入する。
\item   grounded\_summary\_file に Path(str(grounded\_map\_summary.get('summary\_yaml', '') or '')).resolve() の結果を代入する。
\item   grounded\_map\_summary['summary\_yaml'] に relpath\_from(package\_dir, grounded\_summary\_file) の結果を代入する。
\item   evidence\_summary\_path に manifest\_context.get('grounded\_summary\_path') の結果を代入する。
\item   条件 evidence\_summary\_path is not None and (not str(output\_layout['tag'])) を判定し、真なら内部処理を行う。
\item     payload に dict(grounded\_map\_summary) の結果を代入する。
\item     payload.pop(...) を実行する。
\item     with 文で Path(evidence\_summary\_path).open('w', encoding='utf-8', newline='\textbackslash{}n') を管理しながら処理する。
\item       yaml.safe\_dump(...) を実行する。
\item     manifest\_context['grounded\_summary\_payload'] に payload の結果を代入する。
\item   base\_model に build\_model\_from\_map\_assets(base\_model, source\_map\_assets) の結果を代入する。
\item   ocv\_df に pd.read\_csv(source\_map\_assets['ocv\_soc\_map']) の結果を代入する。
\item   上の条件が偽の場合:
\item   条件 explicit\_grounded\_assets を判定し、真なら内部処理を行う。
\item     stage(...) を実行する。
\item     source\_map\_assets に explicit\_grounded\_assets の結果を代入する。
\item     base\_model に build\_model\_from\_map\_assets(base\_model, source\_map\_assets) の結果を代入する。
\item     ocv\_df に pd.read\_csv(source\_map\_assets['ocv\_soc\_map']) の結果を代入する。
\item     上の条件が偽の場合:
\item     条件 bool(options.get('use\_grounded\_base\_maps', False)) を判定し、真なら内部処理を行う。
\item       stage(...) を実行する。
\item       (source\_map\_assets, grounded\_map\_summary) に build\_grounded\_map\_assets(profile\_cfg, Path(output\_layout['grounded\_maps']), pv\_area\_m2=float((profile\_cfg.get('model', \{\}) or \{\}).get('pv\_area', 6.0)), base\_dir=package\_dir) の結果を代入する。
\item       grounded\_summary\_file に Path(str(grounded\_map\_summary.get('summary\_yaml', '') or '')).resolve() の結果を代入する。
\item       grounded\_map\_summary['summary\_yaml'] に relpath\_from(package\_dir, grounded\_summary\_file) の結果を代入する。
\item       base\_model に build\_model\_from\_map\_assets(base\_model, source\_map\_assets) の結果を代入する。
\item       ocv\_df に pd.read\_csv(source\_map\_assets['ocv\_soc\_map']) の結果を代入する。
\item       上の条件が偽の場合:
\item       ocv\_df に load\_ocv\_df(profile\_cfg, profile\_yaml) の結果を代入する。
\item base\_model に neutralize\_identification\_scalars(base\_model) の結果を代入する。
\item raw\_log に str(inputs.get('normalized\_replay\_log\_csv', 'data/identification/raw/observed\_replay\_log.csv') or '').strip() の結果を代入する。
\item observed\_log\_csv に resolve\_relative(package\_dir, raw\_log) の結果を代入する。
\item stage(...) を実行する。
\item logs に normalize\_generic\_log(observed\_log\_csv, actual\_event\_payload=manifest\_context.get('actual\_event\_payload', \{\}), base\_model=base\_model, options=fit\_plan) の結果を代入する。
\item stage(...) を実行する。
\item (logs, solar\_measurement\_calibration) に calibrate\_solar\_measurement\_to\_pack(logs, known\_aux\_power\_w=float((profile\_cfg.get('model', \{\}) or \{\}).get('P\_aux', 21.0))) の結果を代入する。
\item fit\_plan['solar\_measurement\_calibration'] に solar\_measurement\_calibration の結果を代入する。
            \end{enumerate}
            \paragraph{代表コード断片}
            \begin{lstlisting}[language=Python]
def main() -> None:
    ap = argparse.ArgumentParser()
    ap.add_argument("--profile", required=True)
    ap.add_argument("--manifest")
    ap.add_argument(
        "--quality",
        choices=sorted(FIT_QUALITY_PRESETS.keys()),
        default=None,
        help="Override identification.fit_quality; omit to use the profile/manifest YAML value.",
    )
    ap.add_argument("--output-tag", help="Override identification.output_tag for versioned artifacts.")
    ap.add_argument(
        "--adopt-profile",
        action="store_true",
        help="Allow a tagged run to update the canonical profile.yaml; otherwise write a candidate profile in the run directory.",
    )
    ap.add_argument("--allow-map-shape-fit", action="store_true")
    ap.add_argument("--skip-map-shape-fit", action="store_true")
    ap.add_argument(
        "--rebuild-grounded-base-maps",
        action="store_true",
        help="Rebuild product/theory grounded maps before fitting instead of reusing manifest assets.",
    )
    ap.add_argument("--manifest-only", action="store_true")
    args = ap.parse_args()

    profile_yaml = resolve_relative(ROOT, args.profile)
    package_dir = profile_yaml.parent
    stage(f"loading profile: {profile_yaml}")
    profile_cfg = load_profile_yaml(profile_yaml)
    manifest_path, manifest = load_manifest(package_dir, args.manifest)
    manifest_context = resolve_manifest_context(package_dir, manifest)
    inputs = manifest_context["inputs"]
    options = dict(manifest_context["options"])
    identification_cfg = profile_cfg.get("identification", {}) or {}
...
\end{lstlisting}







        \subsection{CLI 引数}
            \begin{longtable}{p{0.07\linewidth} p{0.78\linewidth}}
            \toprule
            行 & 引数 \\
            \midrule
            3077 & \path|--profile| \\
3078 & \path|--manifest| \\
3079 & \path|--quality| \\
3085 & \path|--output-tag| \\
3086 & \path|--adopt-profile| \\
3091 & \path|--allow-map-shape-fit| \\
3092 & \path|--skip-map-shape-fit| \\
3093 & \path|--rebuild-grounded-base-maps| \\
3098 & \path|--manifest-only| \\
            \bottomrule
            \end{longtable}



        \section{処理の流れ}
        \begin{enumerate}[leftmargin=1.7em]
        \item CLI 引数を解釈し、main() から処理を起動する。
        \end{enumerate}

        \section{このファイルの位置づけ}
        この資料群では、\path|scripts/run_vehicle_identification.py| を
        \textbf{identification}の中核として扱う。
        したがって、ここで扱う処理は単独で完結せず、
        前段の profile / launch / telemetry と後段の planner / logger / report へ繋がる。

        \end{document}
